% Oxydoréduction par voie sèche et nombres d'oxydation — Chimie 1ère C — Cours signé OnBuch+
% Programme officiel MINESEC, Première C, D, E, TI. Compiler avec : tectonic N10-oxydoreduction-voie-seche-nombres-oxydation.tex
\newcommand{\DOCMATIERE}{Chimie}
\newcommand{\DOCNIVEAU}{1ère C}
\newcommand{\DOCMODULE}{Module 2 : Oxydoréduction}
\newcommand{\DOCLECON}{10}
\newcommand{\DOCTITRE}{Oxydoréduction par voie sèche et nombres d'oxydation}
% OnBuch+ — préambule « cours signés » ENRICHI (compilé avec tectonic / XeLaTeX)
% Système visuel « Pop » d'OnBuch+ : fond crème (jamais blanc), bordures encre
% épaisses, ombres « dures » décalées, titres Archivo Black en capitales.
% Version enrichie pour les cours de chimie : chimie (mhchem, chemfig),
% graphes (pgfplots), boîtes pédagogiques supplémentaires (définition,
% propriété, méthode, attention, le savais-tu, exercice résolu, exercice,
% TP, bilan), page de garde refaite et signature OnBuch+ ancrée en bas.
%
% Macros attendues dans main.tex AVANT \input{preamble} :
%   \DOCMATIERE  \DOCNIVEAU  \DOCMODULE  \DOCLECON (numéro)  \DOCTITRE
\documentclass[11pt]{article}

\usepackage{fontspec}
\usepackage[french]{babel}
\usepackage[a4paper,margin=2cm,top=3.4cm,bottom=2.8cm,headheight=1.4cm,footskip=1.3cm]{geometry}
\usepackage{amsmath,amssymb,amsfonts}
\usepackage[table]{xcolor} % [table] : fournit aussi \rowcolors (alternance de lignes)
\usepackage{pagecolor}
\usepackage{tikz}
\usetikzlibrary{arrows.meta,positioning,calc,shapes.geometric,shapes.misc,shapes.arrows,decorations.pathmorphing,decorations.markings,patterns,shadows,mindmap,trees,fit,backgrounds,matrix,intersections}
\usepackage{pgfplots}
\pgfplotsset{compat=1.18}
\usepackage[version=4]{mhchem}
\usepackage{chemfig}
\usepackage{enumitem}
\usepackage{fancyhdr}
% [most] charge toutes les bibliothèques tcolorbox (listings, minted, raster,
% documentation...) et gonfle inutilement la mémoire TeX sur un document long
% (~300 boîtes) : on ne charge que ce qui est réellement utilisé.
\usepackage[skins,breakable]{tcolorbox}
\usepackage{graphicx}
\usepackage{colortbl}
\usepackage{booktabs}
\usepackage{tabularx}
\usepackage{array}
\usepackage{multicol}
\usepackage{siunitx}
% parse-numbers=false : accepte aussi \SI{2,5 \times 10^{-3}}{...}
\sisetup{locale=FR,output-decimal-marker={,},group-minimum-digits=4,parse-numbers=false}
\DeclareSIUnit\ohm{\ensuremath{\symup{\Omega}}} % U+2126 absent de la police
\usepackage{qrcode}
% \degree : commande fréquemment utilisée par les modèles pour un angle ou une
% température en dehors de \SI{}{} (non fournie par les paquets chargés).
\providecommand{\degree}{^{\circ}}
\usepackage{lastpage}
\usepackage{hyperref}
\hypersetup{hidelinks}

% ============================================================
% Palette « Pop » OnBuch+ — mêmes hex que la classe P (plus_ui.dart)
% ============================================================
\definecolor{poporange}{HTML}{F59321}
\definecolor{popdark}{HTML}{E07A0C}
\definecolor{popcream}{HTML}{FFF6E8}
\definecolor{popink}{HTML}{1C1714}
\definecolor{poppurple}{HTML}{7A5AE0}
\definecolor{popgreen}{HTML}{2E9E6A}
\definecolor{poppink}{HTML}{E0457B}
\definecolor{popblue}{HTML}{2F7DE1}
\definecolor{popgold}{HTML}{C98A12}
\definecolor{popmuted}{HTML}{8A7F76}
\definecolor{popline}{HTML}{EADFD2}
\definecolor{popgray}{HTML}{8A7F76}
\colorlet{popgrey}{popgray}
% Alias tolérants (le modèle nomme parfois les couleurs différemment)
\colorlet{popwhite}{white}
\colorlet{popblack}{popink}
\colorlet{popred}{poppink}
\colorlet{popyellow}{popgold}
% Teintes claires pour fonds de tableaux / zones
\colorlet{popblueL}{popblue!10!white}
\colorlet{poporangeL}{poporange!14!white}
\colorlet{popgreenL}{popgreen!12!white}
\colorlet{poppurpleL}{poppurple!12!white}
\colorlet{poppinkL}{poppink!12!white}
\colorlet{popgoldL}{popgold!14!white}

\pagecolor{popcream}

% ============================================================
% Typographie
% ============================================================
\defaultfontfeatures{Ligatures=TeX}
\setmainfont{PlusJakartaSans-Medium.ttf}[Path=../../../fonts/,BoldFont=PlusJakartaSans-Bold.ttf]
% Maths en sans-serif (Fira Math) avec chiffres et lettres droites de Plus
% Jakarta, pour que les formules se fondent dans le texte.
\usepackage[math-style=ISO,bold-style=ISO]{unicode-math}
\setmathfont{FiraMath-Regular.otf}[Path=../../../fonts/]
\setmathfont{PlusJakartaSans-Medium.ttf}[Path=../../../fonts/,range={up/num,up/{Latin,latin},\mathup}]
\usepackage{icomma} % virgule décimale sans espace en mode math
% Caractères absents de la police de texte (sinon carré vide) : rendus par
% leur équivalent mathématique, en texte comme en formule.
\usepackage{newunicodechar}
\newunicodechar{α}{\ensuremath{\alpha}}
\newunicodechar{Α}{\ensuremath{\alpha}}
\newunicodechar{β}{\ensuremath{\beta}}
\newunicodechar{δ}{\ensuremath{\delta}}
\newunicodechar{✓}{\popcheck{popgreen}}
\newfontfamily\popdisplay{ArchivoBlack-Regular.ttf}[Path=../../../fonts/]
\newfontfamily\poplogo{SpaceGrotesk-Bold.ttf}[Path=../../../fonts/]
\newfontfamily\popsemi{PlusJakartaSans-SemiBold.ttf}[Path=../../../fonts/]
\newfontfamily\popxbold{PlusJakartaSans-ExtraBold.ttf}[Path=../../../fonts/]
\newfontfamily\popxboldls{PlusJakartaSans-ExtraBold.ttf}[Path=../../../fonts/,LetterSpace=3.0]

\setlist[enumerate]{leftmargin=1.7em,itemsep=5pt,topsep=4pt}
\setlist[itemize]{leftmargin=1.7em,itemsep=4pt,topsep=4pt,label={\color{poporange}\textbullet}}
\setlength{\parindent}{0pt}
\setlength{\parskip}{5pt}
\renewcommand{\arraystretch}{1.25}

% chemfig : liaisons un peu plus épaisses, encre Pop
\setchemfig{atom sep=2em,bond style={line width=0.9pt,popink},cram width=3pt}

% pgfplots : style Pop par défaut
\pgfplotsset{
  every axis/.append style={axis line style={popink,line width=1pt},tick style={popink},
    grid style={popline,line width=0.5pt},label style={font=\small},tick label style={font=\footnotesize}},
  popcurve/.style={poporange,line width=1.8pt,no markers,smooth},
}

% ============================================================
% Pill / squiggle
% ============================================================
\newcommand{\poppill}[3]{%
  \tikz[baseline=-0.55ex]\node[draw=popink,line width=1.2pt,rounded corners=2.6pt,%
    fill=#2,inner xsep=6.5pt,inner ysep=3pt]{%
    \popxboldls\fontsize{7}{7}\selectfont\color{#3}\MakeUppercase{#1}};%
}
\newcommand{\popsquiggle}[1]{%
  \tikz[baseline=-0.4ex]{
    \draw[#1,line width=2.6pt,line cap=round]
      plot [smooth] coordinates {(0,0.09) (0.34,0.01) (0.68,0.21) (1.02,0.01) (1.36,0.10)};
  }%
}
% Mot-clé surligné (marqueur orange sous le texte)
\newcommand{\cle}[1]{{\popxbold\color{popdark}#1}}

% ============================================================
% En-tête / pied
% ============================================================
\pagestyle{fancy}
\fancyhf{}
\renewcommand{\headrulewidth}{0pt}
\renewcommand{\footrulewidth}{0pt}
\lhead{\raisebox{-0.15ex}{\poplogo\fontsize{13}{13}\selectfont\color{popink}OnBuch\color{poporange}+}}
\rhead{\poppill{\DOCMATIERE\ \textbullet\ \DOCNIVEAU\ \textbullet\ Leçon \DOCLECON}{white}{popink}}
\lfoot{\popxbold\fontsize{7.6}{7.6}\selectfont\color{popmuted}\MakeUppercase{Cours signé OnBuch+}\ \textbullet\ {\popsemi\DOCTITRE}}
\rfoot{\tikz[baseline=-0.6ex]\node[rounded corners=8.5pt,fill=popink,minimum height=17pt,minimum width=17pt,inner xsep=5pt,inner sep=0]{\popxbold\fontsize{8}{8}\selectfont\color{popcream}\thepage\,/\,\pageref*{LastPage}};}

% ============================================================
% Titres
% ============================================================
\newcommand{\chaptitre}[2]{%
  \begin{tcolorbox}[enhanced,colback=poporange,colframe=popink,boxrule=2.6pt,arc=11pt,
      drop shadow=popink,left=15pt,right=15pt,top=12pt,bottom=11pt,before skip=3pt,after skip=18pt]
    \poppill{#2}{popink}{popcream}\par\vspace{9pt}
    {\raggedright\hyphenpenalty=10000\exhyphenpenalty=10000\popdisplay\fontsize{22}{24}\selectfont\color{popcream}\MakeUppercase{#1}\par}
  \end{tcolorbox}
}
% Section numérotée : pastille orange avec le numéro + titre Archivo
\newcounter{popsec}
\newcommand{\coursec}[1]{%
  \stepcounter{popsec}\setcounter{popsub}{0}%
  \par\vspace{16pt}\noindent
  \begin{minipage}{\linewidth}
  \tikz[baseline=-0.7ex]\node[draw=popink,line width=1.6pt,fill=poporange,rounded corners=4pt,minimum size=22pt,inner sep=0]{\popdisplay\fontsize{12}{12}\selectfont\color{popcream}\thepopsec};%
  \hspace{8pt}{\popdisplay\fontsize{15}{17}\selectfont\color{popink}\MakeUppercase{#1}}\par
  \vspace{4pt}\noindent{\color{poporange}\rule{36pt}{3.6pt}}
  \end{minipage}\par\nopagebreak\vspace{8pt}
}
\newcounter{popsub}
\newcommand{\courssub}[1]{%
  \stepcounter{popsub}%
  \par\vspace{9pt}\noindent{\popxbold\fontsize{11.5}{13}\selectfont\color{popink}{\color{poporange}\thepopsec.\thepopsub}\hspace{6pt}#1}\par\nopagebreak\vspace{4pt}
}

% ============================================================
% Boîtes pédagogiques — même recette : fond blanc, bordure épaisse colorée,
% ombre dure encre, badge plein. Titre optionnel [#1].
% ============================================================
\tcbset{popbase/.style={breakable,enhanced,colback=white,boxrule=2.3pt,arc=9pt,
  shadow={3pt}{-3pt}{0pt}{popink},left=14pt,right=14pt,top=11pt,bottom=11pt,before skip=11pt,after skip=15pt,
  fontupper=\popsemi\fontsize{10.6}{14.5}\selectfont\color{popink},
  fontlower=\popsemi\fontsize{10.6}{14.5}\selectfont\color{popink}}}
% Coche dessinée (le glyphe ✓ n'existe pas dans les polices du document)
\newcommand{\popcheck}[1]{\tikz[baseline=-0.5ex]\draw[#1,line width=1.6pt,line cap=round,line join=round](-0.13,0.0)--(-0.04,-0.09)--(0.13,0.1);}
\newcommand{\popboxhead}[3]{\poppill{#1}{#2}{white}\ifx\relax#3\relax\else\hspace{6pt}{\popxbold\fontsize{10.6}{12}\selectfont\color{popink}#3}\fi\par\vspace{7pt}}

\newtcolorbox{aretenir}[1][]{popbase,colframe=popgreen,before upper={\popboxhead{À retenir}{popgreen}{#1}}}
\newtcolorbox{exemplebox}[1][]{popbase,colframe=poporange,before upper={\popboxhead{Exemple}{poporange}{#1}}}
\newtcolorbox{definition}[1][]{popbase,colframe=popblue,before upper={\popboxhead{Définition}{popblue}{#1}}}
\newtcolorbox{propriete}[1][]{popbase,colframe=poppurple,before upper={\popboxhead{Propriété}{poppurple}{#1}}}
\newtcolorbox{methode}[1][]{popbase,colframe=popgold,before upper={\popboxhead{Méthode}{popgold}{#1}}}
\newtcolorbox{attention}[1][]{popbase,colframe=poppink,before upper={\popboxhead{Attention}{poppink}{#1}}}
\newtcolorbox{experience}[1][]{popbase,colframe=popblue,colback=popblueL,before upper={\popboxhead{Expérience}{popblue}{#1}}}
% « Le savais-tu ? » : carte violette pleine (contexte, culture, Cameroun)
\newtcolorbox{savaistu}[1][]{popbase,colback=poppurple,colframe=popink,
  fontupper=\popsemi\fontsize{10.4}{14.2}\selectfont\color{white},
  before upper={\poppill{Le savais-tu\,?}{white}{poppurple}\ifx\relax#1\relax\else\hspace{6pt}{\popxbold\color{white}#1}\fi\par\vspace{7pt}}}
% Exercice résolu : énoncé en haut, \tcblower puis la solution sur fond crème
\newcounter{popexo}\newcounter{popexr}
\newtcolorbox{exoresolu}[1][]{popbase,colframe=poporange,colbacklower=popcream,
  segmentation style={popink,line width=1.2pt,dashed},
  before upper={\stepcounter{popexr}\popboxhead{Exercice résolu \thepopexr}{poporange}{#1}},
  before lower={\poppill{Solution}{popink}{popcream}\par\vspace{6pt}}}
% Exercice (non résolu en place) — la correction est dans la partie Corrigés
\newtcolorbox{exercice}[1][]{popbase,colframe=popink,
  before upper={\stepcounter{popexo}\popboxhead{Exercice \thepopexo}{popink}{#1}}}
% Corrigé
\newtcolorbox{corrige}[1][]{popbase,colframe=popgreen,colback=popgreenL,
  before upper={\popboxhead{Corrigé}{popgreen}{#1}}}
% Objectifs / prérequis / bilan
\newtcolorbox{objectifs}{popbase,colframe=popink,colback=poporangeL,
  before upper={\popboxhead{Objectifs de la leçon}{popink}{}}}
\newtcolorbox{prerequis}{popbase,colframe=popink,colback=popblueL,
  before upper={\popboxhead{Prérequis}{popblue}{}}}
\newtcolorbox{bilan}{popbase,colframe=popink,colback=white,boxrule=2.8pt,
  before upper={\popboxhead{Fiche bilan}{poporange}{L'essentiel à connaître pour l'examen}}}
% Figure encadrée avec légende
\newtcolorbox{popfigure}[1][]{enhanced,breakable=false,colback=white,colframe=popline,boxrule=1.4pt,arc=8pt,
  left=10pt,right=10pt,top=10pt,bottom=8pt,before skip=10pt,after skip=12pt,halign=center,
  lower separated=false,halign lower=center,
  fontlower=\popsemi\fontsize{9}{11.5}\selectfont\color{popmuted},
  #1}
\newcommand{\legende}[1]{\par\vspace{6pt}{\popsemi\fontsize{9}{11.5}\selectfont\color{popmuted}#1}}

% ============================================================
% Signature OnBuch+
% ============================================================
\newcommand{\poplogoinline}[1]{{\poplogo\fontsize{#1}{#1}\selectfont\color{popink}OnBuch\color{poporange}+}}
\newcommand{\signatureonbuch}{%
  \begin{tcolorbox}[enhanced,colback=popink,colframe=popink,boxrule=2.2pt,arc=10pt,
      drop shadow=poporange,left=14pt,right=14pt,top=10pt,bottom=10pt,before skip=0pt,after skip=0pt]
    \tikz[baseline=-0.7ex]\node[circle,fill=popgreen,draw=popcream,line width=1.3pt,minimum size=22pt,inner sep=0]{\popcheck{white}};%
    \hspace{9pt}\begin{minipage}[c]{0.62\linewidth}
      {\popxbold\fontsize{12}{13}\selectfont\color{popcream}Signé\ \textbullet\ {\poplogo OnBuch\color{poporange}+}}\par\vspace{2pt}
      {\raggedright\popsemi\fontsize{8}{10.5}\selectfont\color{popline}\MakeUppercase{Équipe pédagogique OnBuch+}\\\MakeUppercase{Conforme au programme officiel MINESEC}\par}
    \end{minipage}\hfill
    {\poplogo\fontsize{22}{22}\selectfont\color{popcream}OnBuch\color{poporange}+}
  \end{tcolorbox}%
}

% ============================================================
% Page de garde
% #1 titre · #2 matière · #3 niveau · #4 module · #5 n° leçon · #6 durée ·
% #7 description / objectifs (texte ou itemize)
% ============================================================
\newcommand{\pagegarde}[7]{%
  \thispagestyle{empty}
  \newgeometry{a4paper,margin=1.7cm,top=1.6cm,bottom=1.4cm}%
  \begin{tikzpicture}[remember picture,overlay]
    \fill[poporange!18] ([xshift=-3.2cm,yshift=-3.6cm]current page.north east) circle (4.6cm);
    \fill[poppurple!14] ([xshift=2.4cm,yshift=7.6cm]current page.south west) circle (3.8cm);
  \end{tikzpicture}%
  {\poplogo\fontsize{30}{30}\selectfont\color{popink}OnBuch\color{poporange}+}\hfill
  \raisebox{0.6ex}{\poppill{Cours signé \textbullet\ Premium}{popink}{popcream}}\par
  \vspace{1pt}\popsquiggle{poppurple}\par
  \vspace{8pt}
  \begin{tcolorbox}[enhanced,colback=poporange,colframe=popink,boxrule=2.8pt,arc=13pt,
      drop shadow=popink,left=18pt,right=18pt,top=12pt,bottom=13pt,after skip=10pt]
    \poppill{#2 \textbullet\ #3}{popink}{popcream}\hspace{5pt}\poppill{#4}{white}{popink}\par\vspace{8pt}
    {\popdisplay\fontsize{14}{14}\selectfont\color{popink}LEÇON #5}\par\vspace{4pt}
    {\raggedright\hyphenpenalty=10000\exhyphenpenalty=10000\popdisplay\fontsize{25}{27}\selectfont\color{popcream}\MakeUppercase{#1}\par}
    \vspace{8pt}
    {\popxbold\fontsize{9.5}{9.5}\selectfont\color{popink}\MakeUppercase{Durée conseillée : #6}}
  \end{tcolorbox}
  \begin{tcolorbox}[enhanced,colback=white,colframe=popink,boxrule=2.2pt,arc=10pt,
      drop shadow=popink,left=16pt,right=16pt,top=10pt,bottom=10pt,after skip=8pt]
    \poppill{Dans cette leçon}{poppurple}{white}\par\vspace{5pt}
    \popsemi\fontsize{9.3}{12.6}\selectfont\color{popink}#7
  \end{tcolorbox}
  \noindent\begin{minipage}[t]{0.487\linewidth}
    \begin{tcolorbox}[enhanced,colback=popgreen,colframe=popink,boxrule=2.2pt,arc=10pt,
        drop shadow=popink,left=11pt,right=11pt,top=10pt,bottom=10pt,height=3.1cm,valign=center]
      \tcbox[colback=white,colframe=popink,boxrule=1.3pt,arc=5pt,boxsep=0pt,nobeforeafter,
        left=3pt,right=3pt,top=3pt,bottom=3pt,valign=center]{\qrcode[height=1.9cm]{https://play.google.com/store/apps/details?id=com.luvvix.onbuch&hl=fr}}%
      \hspace{8pt}\begin{minipage}[c]{0.5\linewidth}
        {\popdisplay\fontsize{11}{12.5}\selectfont\color{white}\MakeUppercase{Télécharge OnBuch}}\par\vspace{4pt}
        {\raggedright\popsemi\fontsize{8.4}{11}\selectfont\color{white}Gratuit \textbullet\ Google Play\\Tuteur IA, annales, concours, résultats\par}
      \end{minipage}
    \end{tcolorbox}
  \end{minipage}\hfill
  \begin{minipage}[t]{0.487\linewidth}
    \begin{tcolorbox}[enhanced,colback=poppurple,colframe=popink,boxrule=2.2pt,arc=10pt,
        drop shadow=popink,left=11pt,right=11pt,top=10pt,bottom=10pt,height=3.1cm,valign=center]
      \tikz[baseline=-0.7ex]\node[circle,fill=white,draw=popink,line width=1.3pt,minimum size=1.6cm,inner sep=0]{\popdisplay\fontsize{24}{24}\selectfont\color{poppurple}+};%
      \hspace{8pt}\begin{minipage}[c]{0.6\linewidth}
        {\popdisplay\fontsize{11}{12.5}\selectfont\color{white}\MakeUppercase{Passe à OnBuch+}}\par\vspace{4pt}
        {\raggedright\popsemi\fontsize{8.4}{11}\selectfont\color{white}Tous les cours signés, Tuteur IA illimité, fiches PDF — onglet OnBuch+ de l'app\par}
      \end{minipage}
    \end{tcolorbox}
  \end{minipage}
  \vspace{8pt}
  \popcontenu
  \vfill
  \signatureonbuch
  \restoregeometry
  \newpage
}

% Rangée « Ce cours contient » (page de garde)
\newcommand{\poptile}[3]{% #1 couleur #2 gros texte #3 légende
  \begin{tcolorbox}[enhanced,colback=#1,colframe=popink,boxrule=2pt,arc=9pt,shadow={2.5pt}{-2.5pt}{0pt}{popink},
      left=6pt,right=6pt,top=9pt,bottom=9pt,halign=center,valign=center,height=1.75cm,width=\linewidth,nobeforeafter]
    {\popdisplay\fontsize{12.5}{13}\selectfont\color{white}\MakeUppercase{#2}}\par\vspace{3pt}
    {\centering\popsemi\fontsize{7.8}{9.5}\selectfont\color{white}#3\par}
  \end{tcolorbox}}
\newcommand{\popcontenu}{%
  \par\noindent{\popxboldls\fontsize{7.5}{7.5}\selectfont\color{popmuted}CE COURS SIGNÉ CONTIENT}\par\vspace{7pt}
  \noindent\begin{minipage}[t]{0.235\linewidth}\poptile{popblue}{Cours}{complet, illustré, pas à pas}\end{minipage}\hfill
  \begin{minipage}[t]{0.235\linewidth}\poptile{poporange}{Méthodes}{et exercices résolus}\end{minipage}\hfill
  \begin{minipage}[t]{0.235\linewidth}\poptile{poppink}{Exercices}{type examen + corrigés}\end{minipage}\hfill
  \begin{minipage}[t]{0.235\linewidth}\poptile{popgreen}{Fiche bilan}{l'essentiel à retenir}\end{minipage}\par}

% Sommaire
\newtcolorbox{sommaire}{enhanced,colback=white,colframe=popink,boxrule=2.2pt,arc=10pt,
  drop shadow=popink,left=18pt,right=18pt,top=13pt,bottom=13pt,before skip=6pt,after skip=20pt,
  before upper={\poppill{Sommaire}{poppurple}{white}\par\vspace{9pt}\popsemi\fontsize{10.6}{14.5}\selectfont\color{popink}}}

% Signature de fin de cours — poussée tout en bas de la dernière page
\newcommand{\findecours}{%
  \par\vfill\nopagebreak
  \begin{center}\popsquiggle{poporange}\end{center}\vspace{4pt}
  \signatureonbuch
}
% Glyphes absents de Fira Math : redéfinis à partir de glyphes disponibles
\AtBeginDocument{%
  \def\setminus{\mathbin{\backslash}}\let\smallsetminus\setminus
  \def\vdots{\mathord{\vbox{\baselineskip3.2pt\lineskiplimit0pt\kern4pt\hbox{.}\hbox{.}\hbox{.}}}}%
  \def\ddots{\mathinner{\mkern1mu\raise6.5pt\hbox{.}\mkern2mu\raise3.6pt\hbox{.}\mkern2mu\raise0.7pt\hbox{.}\mkern1mu}}%
  \def\triangle{\mathord{\scalebox{0.9}{$\symup{\Delta}$}}}%
  \def\nabla{\mathord{\raisebox{\depth}{\scalebox{1}[-1]{$\symup{\Delta}$}}}}%
  \def\checkmark{\popcheck{popdark}}%
  \def\star{\mathbin{\ast}}\def\bigstar{\mathord{\ast}}%
  \def\diamond{\mathbin{\scalebox{0.6}{$\Diamond$}}}%
  \def\bigcup{\mathop{\mathchoice{\raisebox{-0.3em}{\scalebox{1.7}{$\cup$}}}{\scalebox{1.25}{$\cup$}}{\cup}{\cup}}\displaylimits}%
  \def\bigcap{\mathop{\mathchoice{\raisebox{-0.3em}{\scalebox{1.7}{$\cap$}}}{\scalebox{1.25}{$\cap$}}{\cap}{\cap}}\displaylimits}%
  \def\lesseqqgtr{\mathrel{\lessgtr}}\def\bigoplus{\mathop{\oplus}\displaylimits}%
}
\providecommand{\ssi}{si et seulement si} % abréviation fréquente

%% FIGFIX-BEGIN v5 — correctifs globaux des figures (docs/onbuch-generation/tools/figfix)
\usepackage{adjustbox}\usepackage{etoolbox}\tcbuselibrary{raster}
% toute figure TikZ/pgfplots plus large que la ligne est réduite à la largeur disponible
\BeforeBeginEnvironment{tikzpicture}{\begin{adjustbox}{max width=\linewidth}}
\AfterEndEnvironment{tikzpicture}{\end{adjustbox}}
% légendes pgfplots lisibles par-dessus les courbes
\pgfplotsset{every axis/.append style={legend style={fill=white,fill opacity=0.92,text opacity=1,draw=black!15,inner sep=3pt}}}
% étiquettes d'arêtes (syntaxe edge["..."]) avec fond blanc
\tikzset{every edge quotes/.append style={fill=white,inner sep=1.2pt}}
%% FIGFIX-END
\begin{document}

\pagegarde{Oxydoréduction par voie sèche et nombres d'oxydation}{Chimie}{1ère C}{Module 2 : Oxydoréduction}{10}{4 h}{Cette leçon étend la notion d'oxydoréduction aux réactions en phase gazeuse ou solide (voie sèche) et introduit le nombre d'oxydation, outil universel pour identifier et équilibrer les réactions redox. Elle s'achève par les grandes applications industrielles : sidérurgie, aluminothermie, préparation des acides sulfurique et nitrique.
\vspace{4pt}
\begin{multicols}{2}\small
\begin{itemize}
\item À la fin de cette leçon, je sais décrire et équilibrer les combustions du magnésium, du dihydrogène et du carbone dans le dioxygène et le dichlore.
\item À la fin de cette leçon, je sais définir une oxydation et une réduction en termes de variation du nombre d'oxydation.
\item À la fin de cette leçon, je sais appliquer les règles conventionnelles de détermination du nombre d'oxydation d'un élément dans un corps simple ou un composé.
\item À la fin de cette leçon, je sais montrer qu'une réaction est une réaction d'oxydoréduction en comparant les nombres d'oxydation avant et après réaction.
\item À la fin de cette leçon, je sais identifier l'oxydant, le réducteur, l'oxydation et la réduction à l'aide des nombres d'oxydation.
\item À la fin de cette leçon, je sais utiliser les nombres d'oxydation pour équilibrer une équation-bilan de réaction redox.
\end{itemize}\end{multicols}}

\begin{sommaire}
\begin{multicols}{2}
\begin{enumerate}
\item Les réactions d'oxydoréduction par voie sèche
\item Le nombre d'oxydation : définition et règles de calcul
\item Nombres d'oxydation et réactions redox
\item Équilibrer une équation redox grâce aux nombres d'oxydation
\item Applications industrielles : sidérurgie et aluminothermie
\item Préparation industrielle des acides sulfurique et nitrique
\item Travaux pratiques
\item Méthodes et exercices résolus
\item Exercices
\item Corrigés des exercices
\item Fiche bilan
\end{enumerate}
\end{multicols}
\end{sommaire}

\begin{objectifs}
\begin{itemize}
\item À la fin de cette leçon, je sais décrire et équilibrer les combustions du magnésium, du dihydrogène et du carbone dans le dioxygène et le dichlore.
\item À la fin de cette leçon, je sais définir une oxydation et une réduction en termes de variation du nombre d'oxydation.
\item À la fin de cette leçon, je sais appliquer les règles conventionnelles de détermination du nombre d'oxydation d'un élément dans un corps simple ou un composé.
\item À la fin de cette leçon, je sais montrer qu'une réaction est une réaction d'oxydoréduction en comparant les nombres d'oxydation avant et après réaction.
\item À la fin de cette leçon, je sais identifier l'oxydant, le réducteur, l'oxydation et la réduction à l'aide des nombres d'oxydation.
\item À la fin de cette leçon, je sais utiliser les nombres d'oxydation pour équilibrer une équation-bilan de réaction redox.
\item À la fin de cette leçon, je sais indiquer le principe de la sidérurgie et de l'aluminothermie.
\item À la fin de cette leçon, je sais écrire les équations-bilan des étapes de préparation de l'acide sulfurique et de l'acide nitrique.
\end{itemize}
\end{objectifs}

\begin{prerequis}
\begin{itemize}
\item Couples oxydant/réducteur, demi-équations électroniques et réactions redox en solution aqueuse (leçon 9)
\item Écriture et équilibrage d'une équation-bilan (Seconde)
\item Formules des corps simples (O2, H2, Cl2, Mg, C) et de quelques composés courants (H2O, CO2, MgO)
\item Structure électronique et valence des éléments courants (Seconde)
\item Notion de combustion et réactif comburant (Seconde)
\end{itemize}
\end{prerequis}

\newpage

\coursec{Les réactions d'oxydoréduction par voie sèche}

À la fonderie d'aluminium d'Édéa (Alucam) comme dans les petites forges artisanales de Yaoundé, on extrait les métaux de leurs minerais par chauffage intense : c'est la \cle{métallurgie}. De même, lorsqu'un morceau de magnésium brûle avec une flamme éblouissante dans le dioxygène, ou lorsque le charbon de bois rougeoie dans la cuisine, il se produit des réactions sans ions en solution : ce sont des oxydoréductions \og par voie sèche \fg{}. Comment reconnaître ces réactions alors qu'il n'y a pas d'ions pour écrire des demi-équations ? Les chimistes ont inventé un outil puissant : le \cle{nombre d'oxydation}. C'est lui qui permet de comprendre la fabrication de l'acide sulfurique, produit industriel le plus important au monde par tonnage.

\textbf{Problématique.}\par
Dans les leçons précédentes, nous avons étudié les réactions d'oxydoréduction \emph{en solution aqueuse} : les ions y circulent, et l'écriture des demi-équations électroniques vient naturellement. Mais qu'en est-il lorsque les réactifs sont des \emph{solides} et des \emph{gaz}, chauffés, sans aucune goutte d'eau ? La notion d'oxydoréduction reste-t-elle valable ? C'est toute l'ambition de cette section.

\courssub{Qu'est-ce qu'une oxydoréduction par voie sèche ?}

Jusqu'ici, tu as rencontré des couples redox en solution : \ce{Cu^2+}/\ce{Cu}, \ce{Zn^2+}/\ce{Zn}, \ce{MnO4-}/\ce{Mn^2+}... L'eau était partout. Or une immense famille de réactions se déroule \emph{hors de toute solution} : les combustions, les réactions des hauts fourneaux, le soudage par aluminothermie. Elles mettent en jeu des solides et des gaz, souvent portés à haute température.

\begin{definition}[Oxydoréduction par voie sèche]
Une \cle{oxydoréduction par voie sèche} est une réaction d'oxydoréduction réalisée entre des réactifs \textbf{solides et/ou gazeux}, \textbf{sans solvant} (en particulier sans eau), le plus souvent \textbf{à chaud}. Il y a toujours transfert d'électrons d'un \cle{réducteur} (qui perd des électrons) vers un \cle{oxydant} (qui les gagne), même si aucun ion n'existe dans le milieu.
\end{definition}

L'exemple type est la \cle{combustion} : réaction vive, exothermique, entre un \cle{combustible} (le réducteur : métal, dihydrogène, carbone...) et un \cle{comburant} (l'oxydant : le plus souvent le dioxygène \ce{O2}, parfois le dichlore \ce{Cl2}).

\begin{methode}[Reconnaître une combustion]
\begin{enumerate}
\item Identifier le \textbf{comburant} : c'est l'oxydant, presque toujours le dioxygène \ce{O2} (parfois le dichlore \ce{Cl2}).
\item Identifier le \textbf{combustible} : c'est le réducteur, l'espèce qui brûle (métal, \ce{H2}, carbone, hydrocarbure...).
\item Vérifier les trois conditions du \emph{triangle du feu} : combustible, comburant, énergie d'activation (étincelle, flamme, chaleur).
\item Écrire les produits : les éléments du combustible se retrouvent \emph{combinés à l'oxygène} (oxydes) ou au chlore (chlorures).
\end{enumerate}
\end{methode}

\courssub{Combustion du magnésium dans le dioxygène}

\begin{experience}[Le ruban de magnésium brûle dans le dioxygène]
\textbf{Dispositif et protocole.} On tient avec une pince un ruban de magnésium (métal gris brillant) préalablement décapé. On l'enflamme à la flamme d'un bec Bunsen, puis on le plonge dans un flacon rempli de dioxygène pur.

\textbf{Observations.} Le magnésium brûle avec une \textbf{flamme blanche éblouissante} (ne jamais la regarder directement !). Il se forme une \textbf{poudre blanche} qui se dépose sur les parois du flacon : l'oxyde de magnésium \ce{MgO}. La réaction dégage beaucoup de chaleur.

\textbf{Interprétation.} Le magnésium (combustible, réducteur) s'est combiné au dioxygène (comburant, oxydant) :
\[ \ce{2 Mg + O2 -> 2 MgO} \]
\end{experience}

\begin{popfigure}
\begin{center}
\begin{tikzpicture}[scale=0.95, every node/.style={font=\small}]
% Flacon
\draw[thick, popblue] (-1.6,-2.2) -- (-1.6,1.0) -- (-1.2,1.6) -- (1.2,1.6) -- (1.6,1.0) -- (1.6,-2.2) -- cycle;
\draw[popblue, thick] (-1.35,1.6) -- (-1.35,2.0) -- (1.35,2.0) -- (1.35,1.6);
\node[popblue, below] at (0,-2.35) {flacon de dioxygène \ce{O2}};
% Poudre blanche
\fill[popcream] (-1.5,-2.15) -- (1.5,-2.15) -- (1.5,-1.75) -- (-1.5,-1.75) -- cycle;
\node[popdark, font=\footnotesize] at (2.9,-1.95) {poudre blanche \ce{MgO}};
\draw[->, popmuted] (1.55,-1.95) -- (2.1,-1.95);
% Ruban de Mg
\draw[very thick, popmuted] (-0.12,2.6) -- (-0.12,0.4);
\draw[very thick, popmuted] (0.12,2.6) -- (0.12,0.4);
\draw[very thick, popmuted] (-0.12,0.4) -- (0.12,0.4);
\node[popdark, right] at (0.2,2.3) {ruban de Mg};
% Pince
\draw[very thick, popdark] (-0.35,3.4) -- (-0.12,2.55);
\draw[very thick, popdark] (0.35,3.4) -- (0.12,2.55);
\node[popdark, above] at (0,3.45) {pince};
% Flamme éblouissante
\fill[popgold!70] (0,0.35) .. controls (0.5,-0.1) and (0.35,-0.7) .. (0,-1.0) .. controls (-0.35,-0.7) and (-0.5,-0.1) .. (0,0.35);
\fill[white] (0,0.1) .. controls (0.25,-0.15) and (0.18,-0.5) .. (0,-0.7) .. controls (-0.18,-0.5) and (-0.25,-0.15) .. (0,0.1);
\foreach \a in {30,90,150,210,270,330} {\draw[poporange, thick] (\a:0.75) -- (\a:1.15);}
\node[poporange, font=\footnotesize] at (0,-1.45) {flamme blanche éblouissante};
\end{tikzpicture}
\end{center}
\legende{Combustion du ruban de magnésium dans un flacon de dioxygène : flamme blanche intense et dépôt d'oxyde de magnésium \ce{MgO}.}
\end{popfigure}

\begin{exemplebox}[Vérification de la conservation des atomes]
Équilibrons l'équation de combustion du magnésium et vérifions la conservation.
\[ \ce{2 Mg + O2 -> 2 MgO} \]
\begin{center}
\begin{tabularx}{\linewidth}{|l|X|X|}
\hline
\rowcolor{popblueL} \textbf{Élément} & \textbf{Avant (réactifs)} & \textbf{Après (produits)} \\
\hline
Magnésium Mg & $2 \times 1 = 2$ & $2 \times 1 = 2$ \\
\hline
Oxygène O & $1 \times 2 = 2$ & $2 \times 1 = 2$ \\
\hline
\end{tabularx}
\end{center}
Les nombres d'atomes de chaque élément sont conservés : l'équation est correctement équilibrée. Le coefficient $2$ devant \ce{Mg} et devant \ce{MgO} est indispensable : écrire \ce{Mg + O2 -> MgO} serait faux car il y aurait $2$ atomes d'oxygène à gauche et un seul à droite.
\end{exemplebox}

\begin{savaistu}[Le flash des premiers photographes]
Au XIX\textsuperscript{e} siècle, les photographes déclenchaient la combustion d'une poudre de magnésium pour obtenir un éclair blanc très bref et très intense : c'était le premier \emph{flash} de l'histoire de la photographie ! La lumière éblouissante que tu observes au laboratoire est exactement celle qui illuminait les portraits d'autrefois. Aujourd'hui encore, le magnésium entre dans la composition des feux d'artifice blancs et des fusées éclairantes.
\end{savaistu}

\courssub{Combustion du dihydrogène dans le dioxygène}

Le dihydrogène \ce{H2} est un gaz combustible extrêmement inflammable. Mélangé au dioxygène, il forme un mélange \emph{détonant}.

\begin{experience}[La détonation du mélange \ce{H2}/\ce{O2}]
\textbf{Protocole.} On approche une flamme d'un tube contenant du dihydrogène (obtenu par action de l'acide chlorhydrique sur le zinc). En présence de dioxygène de l'air, le gaz s'enflamme.

\textbf{Observations.} On entend une \textbf{détonation} (petite explosion, aboiement caractéristique). Des \textbf{gouttelettes d'eau} se condensent sur les parois froides du récipient.

\textbf{Interprétation.} Le dihydrogène (réducteur) est oxydé par le dioxygène (oxydant) ; le produit est l'eau :
\[ \ce{2 H2 + O2 -> 2 H2O} \]
Cette réaction, très exothermique, propulse les fusées (moteurs à hydrogène liquide) et alimente les piles à combustible.
\end{experience}

\begin{attention}[Sécurité : le mélange détonant]
Le mélange dihydrogène + dioxygène dans les proportions de l'équation-bilan ($2$ volumes de \ce{H2} pour $1$ volume de \ce{O2}) explose violemment à la moindre étincelle. C'est pourquoi on n'approche \textbf{jamais} une flamme d'un récipient de dihydrogène sans précaution, et pourquoi les bonbonnes d'hydrogène sont stockées loin de toute source de chaleur.
\end{attention}

\courssub{Combustion du carbone dans le dioxygène}

Le charbon de bois, utilisé chaque jour dans les cuisines camerounaises, est presque uniquement constitué de carbone. Sa combustion illustre un point capital : \emph{le produit formé dépend de la quantité de dioxygène disponible}.

\textbf{Combustion complète} (dioxygène en excès, braises bien aérées) : le carbone est entièrement transformé en dioxyde de carbone :
\[ \ce{C + O2 -> CO2} \]

\textbf{Combustion incomplète} (dioxygène en défaut, pièce mal aérée, braises qui couvent) : il se forme du monoxyde de carbone :
\[ \ce{2 C + O2 -> 2 CO} \]

\begin{attention}[Le monoxyde de carbone, tueur silencieux]
La combustion \textbf{incomplète} du carbone produit le monoxyde de carbone \ce{CO}, gaz \textbf{incolore, inodore et mortel}. Il se fixe sur l'hémoglobine du sang à la place du dioxygène et provoque l'asphyxie sans aucun signe avant-coureur. Chaque année, des intoxications surviennent lorsque des braseros ou des cuisines au charbon fonctionnent dans des pièces fermées. Règle d'or : \textbf{aérer obligatoirement} toute pièce où brûle du charbon de bois, et ne jamais laisser couver un brasero dans une chambre pendant la nuit.
\end{attention}

\courssub{Action du dichlore sur le magnésium et sur le dihydrogène}

Le dioxygène n'est pas le seul oxydant capable de réagir à sec. Le \cle{dichlore} \ce{Cl2}, gaz jaune-vert toxique, joue un rôle très comparable.

\textbf{Action sur le magnésium.} Le ruban de magnésium chauffé brûle dans le dichlore en donnant une poudre blanche de chlorure de magnésium :
\[ \ce{Mg + Cl2 -> MgCl2} \]

\textbf{Action sur le dihydrogène.} Le mélange \ce{H2}/\ce{Cl2}, exposé à la lumière ou enflammé, réagit avec une vive détonation ; il se forme du chlorure d'hydrogène, gaz âcre qui se dissout dans l'eau pour donner l'acide chlorhydrique :
\[ \ce{H2 + Cl2 -> 2 HCl} \]

L'analogie est frappante : le dichlore \og brûle \fg{} les mêmes réducteurs que le dioxygène. C'est bien le signe que ces réactions appartiennent à la même famille : les oxydoréductions.

\courssub{Interprétation électronique et généralisation}

Comment prouver qu'il s'agit bien de transferts d'électrons, alors qu'il n'y a ni solution ni ions ? Regardons de près le cas du magnésium.

Dans le métal magnésium, l'atome \ce{Mg} possède $2$ électrons périphériques. Dans l'oxyde de magnésium \ce{MgO}, solide \emph{ionique}, le magnésium existe sous forme d'ion \ce{Mg^2+} et l'oxygène sous forme d'ion \ce{O^2-}. Le bilan électronique est donc clair :
\begin{align*}
\text{oxydation : } & \ce{Mg -> Mg^2+ + 2 e-} \quad \text{(le magnésium \emph{perd} 2 électrons)}\\
\text{réduction : } & \ce{O2 + 4 e- -> 2 O^2-} \quad \text{(le dioxygène \emph{gagne} des électrons)}
\end{align*}
En multipliant la première demi-équation par $2$ et en additionnant, on retrouve : \ce{2 Mg + O2 -> 2 MgO}. Il y a bien eu \textbf{oxydation} du magnésium et \textbf{réduction} du dioxygène, exactement comme en solution aqueuse.

De même, dans \ce{MgCl2} (solide ionique \ce{Mg^2+} + $2$\,\ce{Cl-}), le magnésium a perdu ses électrons au profit du chlore : \ce{Cl2 + 2 e- -> 2 Cl-}.

\begin{propriete}[Généralisation]
\begin{itemize}
\item Le \textbf{dioxygène} \ce{O2} et le \textbf{dichlore} \ce{Cl2} sont des \cle{oxydants} : ils captent des électrons.
\item Les \textbf{métaux} (comme \ce{Mg}) et le \textbf{dihydrogène} \ce{H2} sont des \cle{réducteurs} : ils cèdent des électrons.
\item Une combustion est donc toujours une \textbf{oxydoréduction par voie sèche} : le combustible est oxydé, le comburant est réduit.
\end{itemize}
\end{propriete}

\begin{popfigure}
\begin{center}
\begin{tabularx}{\linewidth}{|l|X|X|X|}
\hline
\rowcolor{popblueL} \textbf{Réactifs} & \textbf{Observations} & \textbf{Produit} & \textbf{Équation-bilan} \\
\hline
\ce{Mg} + \ce{O2} & flamme blanche éblouissante & poudre blanche \ce{MgO} & \ce{2 Mg + O2 -> 2 MgO} \\
\hline
\ce{H2} + \ce{O2} & détonation, buée & eau \ce{H2O} & \ce{2 H2 + O2 -> 2 H2O} \\
\hline
\ce{C} + \ce{O2} (excès) & rougeoiement vif & \ce{CO2} & \ce{C + O2 -> CO2} \\
\hline
\ce{C} + \ce{O2} (défaut) & combustion lente & \ce{CO} (mortel) & \ce{2 C + O2 -> 2 CO} \\
\hline
\ce{Mg} + \ce{Cl2} & combustion vive & poudre blanche \ce{MgCl2} & \ce{Mg + Cl2 -> MgCl2} \\
\hline
\ce{H2} + \ce{Cl2} & détonation à la lumière & gaz âcre \ce{HCl} & \ce{H2 + Cl2 -> 2 HCl} \\
\hline
\end{tabularx}
\end{center}
\legende{Bilan des combustions étudiées : dans chaque cas, le combustible (réducteur) est oxydé par le comburant (oxydant : \ce{O2} ou \ce{Cl2}).}
\end{popfigure}

\courssub{Limite de l'approche par demi-équations : vers un nouvel outil}

Pour \ce{MgO} et \ce{MgCl2}, tout va bien : ce sont des composés \emph{ioniques}, on peut écrire les ions \ce{Mg^2+}, \ce{O^2-}, \ce{Cl-} et raisonner sur les demi-équations. Mais que faire de l'eau \ce{H2O}, du dioxyde de carbone \ce{CO2} ou du chlorure d'hydrogène \ce{HCl} ? Ce sont des molécules \emph{covalentes} : il n'existe \textbf{aucun ion} \ce{H+} ou \ce{O^2-} réel dans l'eau formée par combustion ! Écrire \ce{H2 -> 2 H+ + 2 e-} serait une fiction commode mais physiquement incorrecte ici.

\begin{attention}[Piège à éviter]
Ne confonds pas la \emph{réalité physique} (les molécules covalentes ne contiennent pas d'ions) avec les \emph{outils de calcul} du chimiste. Pour traiter \textbf{toutes} les réactions redox — y compris celles des composés covalents gazeux — les chimistes ont imaginé une convention : le \cle{nombre d'oxydation} (ou degré d'oxydation). C'est une charge \emph{fictive} attribuée à chaque atome, qui permet de repérer et d'équilibrer n'importe quelle oxydoréduction, en solution ou à sec. C'est l'objet de la section suivante.
\end{attention}

\begin{aretenir}[L'essentiel de la section]
\begin{itemize}
\item Une oxydoréduction \textbf{par voie sèche} se déroule entre solides et/ou gaz, sans solvant, souvent à chaud.
\item Les combustions de \ce{Mg}, \ce{H2} et \ce{C} dans \ce{O2}, et celles de \ce{Mg} et \ce{H2} dans \ce{Cl2}, en sont des exemples : \ce{2 Mg + O2 -> 2 MgO} ; \ce{2 H2 + O2 -> 2 H2O} ; \ce{C + O2 -> CO2} ; \ce{2 C + O2 -> 2 CO} ; \ce{Mg + Cl2 -> MgCl2} ; \ce{H2 + Cl2 -> 2 HCl}.
\item Le \textbf{comburant} (\ce{O2}, \ce{Cl2}) est l'oxydant ; le \textbf{combustible} (métal, \ce{H2}, \ce{C}) est le réducteur : il y a bien transfert d'électrons.
\item La combustion incomplète du carbone produit \ce{CO}, gaz mortel : aération indispensable.
\item Les demi-équations ioniques échouent pour les composés covalents : d'où la nécessité du \textbf{nombre d'oxydation}.
\end{itemize}
\end{aretenir}

\coursec{Le nombre d'oxydation : définition et règles de calcul}

Dans la section précédente, nous avons découvert les réactions d'oxydoréduction par voie sèche : le magnésium qui brûle dans le dioxygène, le dihydrogène qui s'enflamme dans le dichlore, le carbone qui rougeoit dans l'air. Nous avons reconnu ces réactions en repérant des \emph{transferts d'électrons} entre les espèces. Mais cette méthode atteint vite ses limites : comment repérer un transfert d'électrons dans une réaction comme la transformation du soufre en trioxyde de soufre, étape clé de la fabrication de l'acide sulfurique ? Aucun ion n'apparaît clairement, aucune demi-équation évidente ne se dessine. Il nous faut un outil plus général, plus puissant, capable de traiter \emph{toutes} les réactions redox, même les plus complexes. Cet outil existe : c'est le \cle{nombre d'oxydation}. C'est lui qui va nous servir de << détecteur universel >> d'oxydoréduction dans toute la suite du chapitre.

\courssub{Pourquoi un nouvel outil ? L'idée intuitive}

Reprenons l'exemple de la combustion du magnésium :
\[ \ce{2Mg + O2 -> 2MgO} \]
Ici, tout est simple : le magnésium perd deux électrons (il devient \ce{Mg^2+}), l'oxygène en gagne deux (il devient \ce{O^2-}). Le transfert d'électrons est visible parce que le produit, l'oxyde de magnésium, est un solide \emph{ionique}.

Mais considérons maintenant la combustion du carbone :
\[ \ce{C + O2 -> CO2} \]
Le dioxyde de carbone est une molécule \emph{covalente} : il n'y a ni \ce{C^4+} ni \ce{O^2-} à l'intérieur ! Les électrons des liaisons sont \emph{partagés} entre les atomes, pas transférés. Pourtant, cette réaction est bien une oxydoréduction : le carbone est oxydé, l'oxygène est réduit. Pour le prouver, les chimistes ont inventé une convention astucieuse : faire << comme si >> toutes les liaisons étaient ioniques, en attribuant fictivement les électrons de chaque liaison à l'atome le plus électronégatif. La charge que porterait alors chaque atome est son nombre d'oxydation.

\begin{definition}[Nombre d'oxydation]
Le \cle{nombre d'oxydation} (noté n.o.) d'un élément dans une espèce chimique est la charge, positive ou négative, qu'aurait l'atome de cet élément si tous les électrons de chaque liaison covalente étaient attribués à l'atome le plus électronégatif des deux. C'est une \textbf{charge fictive, purement conventionnelle}. On l'exprime en chiffres romains précédés d'un signe : $+\mathrm{I}$, $-\mathrm{II}$, $+\mathrm{VI}$, etc.
\end{definition}

\begin{attention}[Nombre d'oxydation et charge réelle : ne pas confondre !]
Le nombre d'oxydation est une \textbf{charge fictive} ; la charge d'un ion est une \textbf{charge réelle}. Dans \ce{CO2}, le carbone a pour n.o. $+\mathrm{IV}$, mais il n'existe aucun ion \ce{C^4+} dans cette molécule ! Autre piège classique : dans le dioxygène \ce{O2}, le n.o. de l'oxygène vaut \textbf{zéro} (corps simple), et non $-\mathrm{II}$. La valeur $-\mathrm{II}$ ne vaut que pour l'oxygène \emph{combiné} dans un composé.
\end{attention}

\courssub{Les cinq règles de détermination du nombre d'oxydation}

La définition par l'électronégativité serait fastidieuse à appliquer systématiquement. Heureusement, on en déduit cinq règles simples qui permettent de calculer n'importe quel n.o. en quelques secondes. Apprends-les dans l'ordre : les règles 1 et 2 sont des cas particuliers, les règles 3 et 4 donnent des valeurs de référence, et la règle 5 est la clé de tous les calculs.

\textbf{Règle 1 : corps simples.} Le nombre d'oxydation d'un élément dans un \cle{corps simple} est \textbf{nul}. En effet, entre deux atomes identiques, l'électronégativité est la même : aucun ne << gagne >> les électrons de la liaison. Ainsi : n.o.(Fe) $= 0$ dans le fer métal, n.o.(O) $= 0$ dans \ce{O2}, n.o.(H) $= 0$ dans \ce{H2}, n.o.(Cl) $= 0$ dans \ce{Cl2}, n.o.(S) $= 0$ dans \ce{S8}.

\textbf{Règle 2 : ions monoatomiques.} Le nombre d'oxydation d'un élément dans un \cle{ion monoatomique} est égal à la \textbf{charge de cet ion}. Ici, la charge fictive coïncide avec la charge réelle : n.o.(Na) $= +\mathrm{I}$ dans \ce{Na+}, n.o.(Cl) $= -\mathrm{I}$ dans \ce{Cl-}, n.o.(Al) $= +\mathrm{III}$ dans \ce{Al^3+}.

\textbf{Règle 3 : l'oxygène dans les composés.} Dans un composé, le n.o. de l'\cle{oxygène} vaut en général $-\mathrm{II}$, car c'est l'un des éléments les plus électronégatifs. \emph{Exception} : dans les peroxydes (qui contiennent le groupe $-\mathrm{O}-\mathrm{O}-$, comme l'eau oxygénée \ce{H2O2}), le n.o. de l'oxygène vaut $-\mathrm{I}$.

\textbf{Règle 4 : l'hydrogène dans les composés.} Dans un composé, le n.o. de l'\cle{hydrogène} vaut en général $+\mathrm{I}$, car il est moins électronégatif que la plupart des non-métaux. \emph{Exception} : dans les hydrures métalliques (comme \ce{NaH} ou \ce{CaH2}), l'hydrogène est l'atome le plus électronégatif du couple, et son n.o. vaut $-\mathrm{I}$.

\textbf{Règle 5 : règle de la somme.} La somme algébrique des nombres d'oxydation de tous les atomes d'une espèce est égale à la \textbf{charge globale} de cette espèce : zéro pour une molécule neutre, et la charge de l'ion pour un ion polyatomique.

\begin{methode}[Calculer un nombre d'oxydation inconnu]
Pour déterminer le n.o. d'un élément dans une espèce chimique :
\begin{enumerate}
\item \textbf{Poser} $x$ le nombre d'oxydation cherché ;
\item \textbf{Écrire} la somme algébrique des n.o. de tous les atomes, en utilisant les valeurs de référence : oxygène $-\mathrm{II}$, hydrogène $+\mathrm{I}$ (règles 3 et 4) ;
\item \textbf{Égaler} cette somme à la charge globale de l'espèce (règle 5) ;
\item \textbf{Résoudre} l'équation obtenue et exprimer le résultat en chiffres romains signés.
\end{enumerate}
\end{methode}

\begin{center}
\begin{tabularx}{\linewidth}{|l|X|l|}
\hline
\rowcolor{popblueL} \textbf{Règle} & \textbf{Énoncé} & \textbf{Exemple} \\
\hline
1 & Corps simple : n.o. $= 0$ & \ce{O2}, \ce{Fe}, \ce{S8} : $0$ \\
\hline
2 & Ion monoatomique : n.o. $=$ charge & \ce{Al^3+} : $+\mathrm{III}$ \\
\hline
3 & Oxygène combiné : $-\mathrm{II}$ (peroxydes : $-\mathrm{I}$) & O dans \ce{H2O} : $-\mathrm{II}$ \\
\hline
4 & Hydrogène combiné : $+\mathrm{I}$ (hydrures : $-\mathrm{I}$) & H dans \ce{H2O} : $+\mathrm{I}$ \\
\hline
5 & Somme des n.o. $=$ charge de l'espèce & \ce{MnO4-} : somme $= -\mathrm{I}$ \\
\hline
\end{tabularx}
\end{center}

\courssub{Exemples guidés : la méthode en action}

\begin{exemplebox}[Le manganèse dans l'ion permanganate]
Déterminons le n.o. du manganèse dans l'ion permanganate \ce{MnO4-} (le fameux ion violet de Kalinangan, utilisé comme désinfectant dans les pharmacies camerounaises).

Posons $x$ le n.o. du manganèse. L'ion contient quatre atomes d'oxygène, chacun de n.o. $-\mathrm{II}$, et sa charge globale est $-\mathrm{I}$. La règle 5 s'écrit :
\[ x + 4 \times (-\mathrm{II}) = -\mathrm{I} \quad\Longrightarrow\quad x = -\mathrm{I} + \mathrm{VIII} = +\mathrm{VII} \]
Le manganèse est au degré d'oxydation $+\mathrm{VII}$ : on écrit Mn(VII). C'est le degré d'oxydation maximal du manganèse, ce qui explique le grand pouvoir oxydant de l'ion permanganate.
\end{exemplebox}

\begin{exemplebox}[Le soufre dans l'acide sulfurique]
Déterminons le n.o. du soufre dans \ce{H2SO4}, l'acide sulfurique, produit phare de l'industrie chimique mondiale.

Posons $x$ le n.o. du soufre. La molécule est neutre (charge globale nulle), avec deux hydrogènes à $+\mathrm{I}$ et quatre oxygènes à $-\mathrm{II}$ :
\[ 2 \times (+\mathrm{I}) + x + 4 \times (-\mathrm{II}) = 0 \quad\Longrightarrow\quad x = \mathrm{VIII} - \mathrm{II} = +\mathrm{VI} \]
Le soufre est au degré $+\mathrm{VI}$ dans l'acide sulfurique : c'est son degré d'oxydation maximal.
\end{exemplebox}

\begin{exemplebox}[L'azote dans l'ammoniac et dans l'ion nitrate]
\textbf{Dans l'ammoniac \ce{NH3}} (molécule neutre, H à $+\mathrm{I}$) :
\[ x + 3 \times (+\mathrm{I}) = 0 \quad\Longrightarrow\quad x = -\mathrm{III} \]
L'azote est au degré $-\mathrm{III}$, son degré le plus \emph{réduit}.

\textbf{Dans l'ion nitrate \ce{NO3-}} (charge $-\mathrm{I}$, O à $-\mathrm{II}$) :
\[ x + 3 \times (-\mathrm{II}) = -\mathrm{I} \quad\Longrightarrow\quad x = +\mathrm{V} \]
L'azote est au degré $+\mathrm{V}$, son degré le plus \emph{oxydé}. Entre $-\mathrm{III}$ et $+\mathrm{V}$, l'azote parcourt huit degrés d'oxydation : c'est l'un des éléments les plus polyvalents du tableau périodique, et c'est toute la richesse de la chimie des engrais et de l'acide nitrique.
\end{exemplebox}

\begin{exemplebox}[Le carbone dans le dioxyde de carbone]
Dans \ce{CO2} (molécule neutre, O à $-\mathrm{II}$) :
\[ x + 2 \times (-\mathrm{II}) = 0 \quad\Longrightarrow\quad x = +\mathrm{IV} \]
Le carbone est au degré $+\mathrm{IV}$. Lors de la combustion $\ce{C + O2 -> CO2}$, le carbone passe de $0$ à $+\mathrm{IV}$ : son n.o. \emph{augmente}, il est donc \emph{oxydé} — exactement ce que nous avions pressenti sans pouvoir le prouver par les ions !
\end{exemplebox}

\courssub{L'escalier des nombres d'oxydation du soufre}

Un même élément peut exister à plusieurs nombres d'oxydation. Le soufre en est l'illustration parfaite : on le rencontre au nombre d'oxydation $-\mathrm{II}$ dans le sulfure d'hydrogène \ce{H2S} (gaz à l'odeur d'œuf pourri), au nombre d'oxydation $0$ dans le soufre natif, au nombre d'oxydation $+\mathrm{IV}$ dans le dioxyde de soufre \ce{SO2}, et au nombre d'oxydation $+\mathrm{VI}$ dans le trioxyde \ce{SO3} et l'acide sulfurique \ce{H2SO4}. Chaque << marche >> de cet escalier correspond à une étape de la fabrication industrielle de l'acide sulfurique, que nous étudierons dans la dernière section.

\begin{popfigure}
\begin{center}
\begin{tikzpicture}[scale=1.0, every node/.style={font=\small}]
% Marches de l'escalier
\fill[popgreenL] (0,0) rectangle (2.6,0.8);
\fill[popgoldL] (2.6,0.8) rectangle (5.2,1.8);
\fill[poporangeL] (5.2,1.8) rectangle (7.8,2.8);
\fill[poppinkL] (7.8,2.8) rectangle (10.4,3.8);
% Contours
\draw[popline, thick] (0,0) rectangle (2.6,0.8);
\draw[popline, thick] (2.6,0.8) rectangle (5.2,1.8);
\draw[popline, thick] (5.2,1.8) rectangle (7.8,2.8);
\draw[popline, thick] (7.8,2.8) rectangle (10.4,3.8);
% Textes des marches
\node[popdark] at (1.3,0.4) {\textbf{\ce{H2S}} : n.o. $= -\mathrm{II}$};
\node[popdark] at (3.9,1.3) {\textbf{\ce{S}} : n.o. $= 0$};
\node[popdark] at (6.5,2.3) {\textbf{\ce{SO2}} : n.o. $= +\mathrm{IV}$};
\node[popdark] at (9.1,3.3) {\textbf{\ce{SO3}, \ce{H2SO4}} : $+\mathrm{VI}$};
% Flèches d'oxydation
\draw[-{Stealth}, very thick, poporange] (2.2,1.0) -- (3.0,1.6);
\draw[-{Stealth}, very thick, poporange] (4.8,2.0) -- (5.6,2.6);
\draw[-{Stealth}, very thick, poporange] (7.4,3.0) -- (8.2,3.6);
\node[poporange, font=\small\bfseries] at (5.2,4.5) {Oxydation : le n.o. augmente};
\draw[-{Stealth}, very thick, popblue] (8.2,4.1) -- (2.2,0.9);
\node[popblue, font=\small\bfseries] at (8.6,0.2) {Réduction : le n.o. diminue};
\end{tikzpicture}
\end{center}
\legende{L'escalier des nombres d'oxydation du soufre : monter une marche, c'est oxyder le soufre ; descendre, c'est le réduire.}
\end{popfigure}

\begin{aretenir}[L'essentiel sur le nombre d'oxydation]
\begin{itemize}
\item Le n.o. est une \textbf{charge fictive} conventionnelle, notée en chiffres romains signés.
\item Corps simple : n.o. $= 0$ ; ion monoatomique : n.o. $=$ charge de l'ion.
\item Oxygène combiné : $-\mathrm{II}$ (sauf peroxydes : $-\mathrm{I}$) ; hydrogène combiné : $+\mathrm{I}$ (sauf hydrures : $-\mathrm{I}$).
\item La somme des n.o. de tous les atomes égale la charge globale de l'espèce : c'est la règle qui permet tous les calculs.
\item Une \textbf{augmentation} du n.o. traduit une \textbf{oxydation} ; une \textbf{diminution} traduit une \textbf{réduction}.
\end{itemize}
\end{aretenir}

\begin{exoresolu}[Nombres d'oxydation du chlore dans quatre espèces]
\textbf{Énoncé.} Calculer le nombre d'oxydation du chlore dans chacune des espèces suivantes : \ce{HCl}, \ce{Cl2}, \ce{HClO} (acide hypochloreux, principe actif de l'eau de Javel) et \ce{ClO3-} (ion chlorate).
\tcblower
\textbf{Solution détaillée.}

\textbf{Dans \ce{HCl}} (molécule neutre, H à $+\mathrm{I}$) :
\[ x + (+\mathrm{I}) = 0 \quad\Longrightarrow\quad x = -\mathrm{I} \]

\textbf{Dans \ce{Cl2}} (corps simple, règle 1) :
\[ \text{n.o.}(\ce{Cl}) = 0 \]

\textbf{Dans \ce{HClO}} (molécule neutre, H à $+\mathrm{I}$, O à $-\mathrm{II}$) :
\[ (+\mathrm{I}) + x + (-\mathrm{II}) = 0 \quad\Longrightarrow\quad x = +\mathrm{I} \]

\textbf{Dans \ce{ClO3-}} (ion de charge $-\mathrm{I}$, O à $-\mathrm{II}$) :
\[ x + 3 \times (-\mathrm{II}) = -\mathrm{I} \quad\Longrightarrow\quad x = +\mathrm{V} \]

\textbf{Bilan.} Le chlore existe donc aux degrés $-\mathrm{I}$, $0$, $+\mathrm{I}$ et $+\mathrm{V}$ selon les espèces : c'est un élément aux multiples visages, ce qui explique la richesse de sa chimie (désinfectants, eau de Javel, traitement de l'eau potable).
\end{exoresolu}

\begin{exercice}[À toi de jouer]
Calculer le nombre d'oxydation : (a) du chrome dans l'ion dichromate \ce{Cr2O7^2-} ; (b) du phosphore dans l'acide phosphorique \ce{H3PO4} ; (c) du carbone dans le méthane \ce{CH4} ; (d) du fer dans l'ion \ce{Fe^3+}.
\end{exercice}

\begin{corrige}[Exercice : nombres d'oxydation]
(a) Dans \ce{Cr2O7^2-} : $2x + 7 \times (-\mathrm{II}) = -\mathrm{II}$, donc $2x = +\mathrm{XII}$ et $x = +\mathrm{VI}$.

(b) Dans \ce{H3PO4} : $3 \times (+\mathrm{I}) + x + 4 \times (-\mathrm{II}) = 0$, donc $x = +\mathrm{V}$.

(c) Dans \ce{CH4} : $x + 4 \times (+\mathrm{I}) = 0$, donc $x = -\mathrm{IV}$.

(d) \ce{Fe^3+} est un ion monoatomique : n.o. $= +\mathrm{III}$ (règle 2).
\end{corrige}

\begin{savaistu}[Le fer, un élément à plusieurs couleurs]
Un même élément peut exister à plusieurs nombres d'oxydation, et cela se voit à l'œil nu ! Le fer existe couramment aux degrés $+\mathrm{II}$ et $+\mathrm{III}$ : les sels de fer(II) sont vert pâle (comme le sulfate de fer(II) vendu en pharmacie contre l'anémie), tandis que les sels de fer(III) sont jaune-orangé, couleur que l'on retrouve dans les latérites rouges qui teintent les routes du Cameroun — cette couleur est celle de l'oxyde de fer(III) \ce{Fe2O3}. Mieux encore : l'aimant naturel, la magnétite \ce{Fe3O4}, contient le fer aux \emph{deux} degrés à la fois, $+\mathrm{II}$ et $+\mathrm{III}$ ! Vérifie-le toi-même : $2 \times (+\mathrm{III}) + (+\mathrm{II}) + 4 \times (-\mathrm{II}) = 0$.
\end{savaistu}

Nous disposons désormais d'un outil de calcul simple et universel. Dans la prochaine section, nous allons l'exploiter pour répondre à la question essentielle : comment reconnaître, à coup sûr, qu'une réaction est une réaction d'oxydoréduction, et comment identifier l'oxydant et le réducteur ?

\coursec{Oxydoréduction par voie sèche et nombres d'oxydation}

\courssub{Les réactions d'oxydoréduction par voie sèche : les combustions}

On appelle \cle{voie sèche} les réactions qui se produisent à haute température, en l'absence de solvant (souvent à l'état gazeux ou solide fondu). Les exemples types au programme sont les \textbf{combustions} dans le dioxygène $\ce{O2}$ ou le dichlore $\ce{Cl2}$. Rappelons qu'une combustion est une réaction redox rapide, très exothermique, accompagnée de flamme ou d'incandescence.

\begin{experience}[Combustions caractéristiques en voie sèche]
\begin{itemize}
\item \textbf{Magnésium dans le dioxygène :} Un ruban de Mg tenu à la pince à creuset s'enflamme au chalumeau ; il brûle avec une \textbf{lumière blanche éblouissante} (danger pour les yeux !) en formant une poudre blanche, l'oxyde de magnésium \ce{MgO}.
\[ \ce{2 Mg(s) + O2(g) -> 2 MgO(s)} \]
\item \textbf{Dihydrogène dans le dioxygène :} Un tube à essai rempli de $\ce{H2}$ (test de l'étincelle : \og pop \fg{} caractéristique) ; la flamme est pâle, presque invisible le jour, produisant de la vapeur d'eau.
\[ \ce{2 H2(g) + O2(g) -> 2 H2O(g)} \]
\item \textbf{Carbone (coke, graphite) dans le dioxygène :} Selon la quantité d'air, combustion \textbf{complète} (dioxyde de carbone) ou \textbf{incomplète} (monoxyde de carbone, toxique).
\[ \ce{C(s) + O2(g) -> CO2(g)} \quad \text{(complète)} \]
\[ \ce{2 C(s) + O2(g) -> 2 CO(g)} \quad \text{(incomplète)} \]
\item \textbf{Dans le dichlore :} Le magnésium et l'hydrogène brûlent aussi dans $\ce{Cl2}$, formant des chlorures.
\[ \ce{Mg(s) + Cl2(g) -> MgCl2(s)} \qquad \ce{H2(g) + Cl2(g) -> 2 HCl(g)} \]
\end{itemize}
\end{experience}

\begin{aretenir}[Analyse électronique (rappel de Seconde)]
Dans ces combustions, il y a \textbf{transfert d'électrons} :
\begin{itemize}
\item Le métal (Mg) ou l'élément réducteur (H, C) \textbf{perd des électrons} : c'est l'\textbf{oxydation}.
\item Le dioxygène ou le dichlore \textbf{gagne des électrons} : c'est la \textbf{réduction}.
\end{itemize}
Exemple : $\ce{Mg -> Mg^2+ + 2 e-}$ (oxydation) ; $\ce{O2 + 4 e- -> 2 O^2-}$ (réduction).
Mais cette approche électronique montre ses limites dès que les liaisons deviennent covalentes (ex. $\ce{HCl}$, $\ce{CH4}$). Nous avons besoin d'un outil plus universel : le \textbf{nombre d'oxydation}.
\end{aretenir}

\courssub{Le nombre d'oxydation : définition et règles de calcul}

\begin{definition}[Nombre d'oxydation (n.o.)]
Le \cle{nombre d'oxydation} d'un atome dans une espèce chimique est la \textbf{charge électrique formelle} que cet atome porterait si l'on considérait toutes les liaisons comme \textbf{purement ioniques}, en attribuant les électrons de chaque liaison à l'atome le plus \textbf{électronégatif}. C'est un entier (positif, négatif ou nul), noté \og n.o. \fg{} ou \og N.ox. \fg{}.
\end{definition}

\begin{propriete}[Règles conventionnelles de calcul du n.o.]
On applique les règles dans l'ordre de priorité suivant :
\begin{enumerate}
\item \textbf{Corps simple :} n.o. $= 0$ (ex. $\ce{O2}$, $\ce{Cl2}$, $\ce{Fe}$, $\ce{S8}$, $\ce{P4}$).
\item \textbf{Ion monoatomique :} n.o. $=$ charge de l'ion (ex. $\ce{Na+}$ : $+I$ ; $\ce{Fe^3+}$ : $+III$ ; $\ce{O^2-}$ : $-II$).
\item \textbf{Fluor} : toujours $-I$ (l'élément le plus électronégatif).
\item \textbf{Oxygène} : généralement $-II$. \textbf{Exceptions :} peroxydes ($\ce{O2^2-}$, $\ce{H2O2}$) $\to$ $-I$ ; superoxydes ($\ce{O2-}$) $\to$ $-1/2$ ; fluorures d'oxygène ($\ce{OF2}$) $\to$ $+II$.
\item \textbf{Hydrogène} : généralement $+I$. \textbf{Exception :} hydrures métalliques ($\ce{NaH}$, $\ce{CaH2}$) $\to$ $-I$.
\item \textbf{Halogènes (Cl, Br, I)} : généralement $-I$, sauf combinés à l'oxygène ou au fluor où ils sont positifs.
\item \textbf{Somme des n.o.} : Dans une espèce neutre, la somme vaut $0$ ; dans un ion, elle vaut la charge de l'ion.
\end{enumerate}
\end{propriete}

\begin{methode}[Calculer le n.o. d'un élément dans une espèce]
\begin{enumerate}
\item Identifier les éléments dont le n.o. est \textbf{fixe} par les règles 1 à 6.
\item Poser $x$ = n.o. de l'élément recherché.
\item Écrire l'équation : $\sum$ (n.o. connus) $+$ (coefficient $\times$ $x$) $=$ charge globale.
\item Résoudre pour $x$.
\end{enumerate}
\end{methode}

\begin{exemplebox}[Calculs d'entraînement]
\begin{enumerate}
\item \ce{MnO4-} : O = $-II$. $x + 4(-II) = -1 \Rightarrow x = +VII$.
\item \ce{Cr2O7^2-} : O = $-II$. $2x + 7(-II) = -2 \Rightarrow 2x = 12 \Rightarrow x = +VI$.
\item \ce{H2S2O8} (peroxydisulfate) : H = $+I$, O = $-II$ (sauf liaison O--O). Structure $\ce{HO-O-S(=O)2-O-O-S(=O)2-OH}$. 2 atomes O en liaison O--O ($-I$), 6 atomes O doubles liaisons ($-II$). $2(+I) + 2(-I) + 6(-II) + 2x = 0 \Rightarrow 2 - 2 - 12 + 2x = 0 \Rightarrow x = +VI$.
\item \ce{Fe3O4} (magnétite) : O = $-II$. $3x + 4(-II) = 0 \Rightarrow 3x = 8 \Rightarrow x = +8/3$. C'est une \textbf{valeur moyenne} : $\ce{Fe3O4}$ s'écrit $\ce{FeO.Fe2O3}$ (mélange $\ce{Fe^2+}$ et $\ce{Fe^3+}$).
\end{enumerate}
\end{exemplebox}

\begin{attention}[Pièges fréquents]
\begin{itemize}
\item Ne pas confondre \textbf{nombre d'oxydation} (entier formel) et \textbf{charge réelle} (partielle, non entière dans les liaisons covalentes).
\item Dans $\ce{H2O2}$, l'oxygène est à $-I$ (peroxyde), pas $-II$.
\item Dans $\ce{NH3}$, l'hydrogène est à $+I$ (moins électronégatif que N), donc N = $-III$.
\item Pour les ions polyatomiques, ne pas oublier la charge globale dans la somme !
\end{itemize}
\end{attention}

\courssub{Nombres d'oxydation et réactions redox}

Dans la section précédente, tu as appris à \cle{calculer le nombre d'oxydation} de n'importe quel élément. Tu possèdes désormais un véritable « thermomètre » de l'état d'oxydation. Voyons comment il permet de reconnaître et caractériser une réaction d'oxydoréduction, y compris quand la définition électronique est inapplicable directement.

\courssub{Une idée simple : suivre les variations du nombre d'oxydation}

Reprenons la combustion du magnésium dans le dioxygène :

\[ \ce{2 Mg + O2 -> 2 MgO} \]

Calculons les nombres d'oxydation :
\begin{itemize}
\item Réactifs : \ce{Mg} (corps simple) $\to$ n.o.(Mg) $= 0$ ; \ce{O2} (corps simple) $\to$ n.o.(O) $= 0$.
\item Produit \ce{MgO} : O (règle 4) $\to$ $-II$ ; somme nulle $\to$ n.o.(Mg) $= +II$.
\end{itemize}

Le n.o. du magnésium \textbf{monte} de $0$ à $+II$ (il est oxydé) ; celui de l'oxygène \textbf{descend} de $0$ à $-II$ (il est réduit). C'est une loi générale.

\begin{definition}[Oxydation et réduction : définitions par les nombres d'oxydation]
\begin{itemize}
\item Un élément subit une \cle{oxydation} si son \textbf{nombre d'oxydation augmente} au cours de la réaction.
\item Un élément subit une \cle{réduction} si son \textbf{nombre d'oxydation diminue} au cours de la réaction.
\item Le \cle{réducteur} est l'espèce qui contient l'élément dont le n.o. \textbf{augmente} (c'est lui qui est oxydé).
\item L'\cle{oxydant} est l'espèce qui contient l'élément dont le n.o. \textbf{diminue} (c'est lui qui est réduit).
\end{itemize}
\end{definition}

\begin{aretenir}[Moyen mnémotechnique : l'escalier]
L'\textbf{oxydation}, c'est la \textbf{montée} du n.o. (flèche $\nearrow$) ; la \textbf{réduction}, c'est la \textbf{descente} (flèche $\searrow$).
\textbf{Attention :} Le \textbf{réducteur est oxydé} (son n.o. augmente), l'\textbf{oxydant est réduit} (son n.o. diminue). Le nom de l'espèce désigne ce qu'elle \emph{fait subir à l'autre}, pas ce qu'elle subit elle-même.
\end{aretenir}

\courssub{Le critère de reconnaissance d'une réaction redox}

\begin{propriete}[Critère de reconnaissance universel]
Une réaction chimique est une \cle{réaction d'oxydoréduction} \textbf{si et seulement si} le nombre d'oxydation d'\textbf{au moins un élément} varie entre les réactifs et les produits. Si aucun n.o. ne change, la réaction n'est pas redox.
\end{propriete}

\begin{methode}[Montrer qu'une réaction est (ou n'est pas) redox]
\begin{enumerate}
\item Écrire l'équation-bilan \textbf{équilibrée}.
\item Calculer le \textbf{n.o. de chaque élément} dans chaque réactif et produit.
\item \textbf{Comparer} les valeurs avant/après, élément par élément.
\item Conclure : si au moins un n.o. a changé $\to$ réaction redox (identifier oxydé/réduit, oxydant/réducteur) ; sinon $\to$ pas redox.
\end{enumerate}
\end{methode}

\courssub{Exemples d'application du critère}

\textbf{Exemple 1 : Combustion du magnésium (reprise).}
\[ \ce{2 Mg + O2 -> 2 MgO} \]

\begin{center}
\begin{tabularx}{\linewidth}{|l|X|X|X|}
\hline
\rowcolor{popblueL} \textbf{Élément} & \textbf{n.o. (réactifs)} & \textbf{n.o. (produits)} & \textbf{Conclusion} \\
\hline
Mg & $0$ (corps simple) & $+II$ (dans \ce{MgO}) & $\nearrow$ \textbf{Oxydation} \\
\hline
O & $0$ (dans \ce{O2}) & $-II$ (dans \ce{MgO}) & $\searrow$ \textbf{Réduction} \\
\hline
\end{tabularx}
\end{center}

\ce{Mg} est le \cle{réducteur}, \ce{O2} est l'\cle{oxydant}.

\begin{popfigure}
\begin{center}
\begin{tikzpicture}[scale=1.0, every node/.style={font=\small}]
\node[font=\large] (eq) at (0,0) {$\ce{2 Mg} \;+\; \ce{O2} \;\longrightarrow\; \ce{2 MgO}$};
\node[above=0.15cm of eq, xshift=-2.6cm, popblue] {\footnotesize $0$};
\node[above=0.15cm of eq, xshift=-0.9cm, popgreen] {\footnotesize $0$};
\node[above=0.15cm of eq, xshift=1.7cm, popblue] {\footnotesize $+II$};
\node[above=0.15cm of eq, xshift=2.5cm, popgreen] {\footnotesize $-II$};
\draw[-{Stealth}, very thick, poporange] (-2.9,0.75) .. controls (-2.9,1.6) and (1.4,1.6) .. (1.4,0.75)
  node[midway, above, poporange, font=\small\bfseries] {oxydation : $0 \rightarrow +II$};
\draw[-{Stealth}, very thick, poppurple] (-0.7,-0.35) .. controls (-0.7,-1.3) and (2.6,-1.3) .. (2.6,-0.35)
  node[midway, below, poppurple, font=\small\bfseries] {réduction : $0 \rightarrow -II$};
\end{tikzpicture}
\end{center}
\legende{Variations des n.o. : Mg monte (oxydation), O descend (réduction).}
\end{popfigure}

\textbf{Exemple 2 : Synthèse du chlorure d'hydrogène.}
\[ \ce{H2 + Cl2 -> 2 HCl} \]

\begin{itemize}
\item Réactifs : \ce{H2}, \ce{Cl2} corps simples $\to$ n.o.(H) $= 0$, n.o.(Cl) $= 0$.
\item Produit \ce{HCl} : H $= +I$ (règle 5) $\to$ Cl $= -I$ (somme nulle).
\end{itemize}
H : $0 \to +I$ ($\nearrow$ \textbf{oxydation}) $\to$ \ce{H2} \textbf{réducteur}.
Cl : $0 \to -I$ ($\searrow$ \textbf{réduction}) $\to$ \ce{Cl2} \textbf{oxydant}.

\begin{attention}[Réaction redox sans ions !]
Dans $\ce{H2 + Cl2 -> 2 HCl}$, \textbf{aucun ion n'apparaît} (liaison H--Cl covalente polarisée). Impossible d'écrire des demi-équations avec $\ce{e-}$ libres. Pourtant, les n.o. prouvent sans ambiguïté qu'il s'agit d'une réaction redox. Le \cle{nombre d'oxydation} est l'outil \textbf{universel} : il fonctionne même quand le transfert d'électrons n'est que partiel. Ne cherche pas systématiquement des ions pour conclure !
\end{attention}

\textbf{Contre-exemple : Décomposition thermique du calcaire (CIMENCAM, Figuil).}
\[ \ce{CaCO3 ->[{900\,^\circ C}] CaO + CO2} \]

\begin{center}
\begin{tabularx}{\linewidth}{|l|X|X|X|}
\hline
\rowcolor{popblueL} \textbf{Élément} & \textbf{n.o. dans \ce{CaCO3}} & \textbf{n.o. dans produits} & \textbf{Variation ?} \\
\hline
Ca & $+II$ & $+II$ (\ce{CaO}) & Non \\
\hline
C & $+IV$ & $+IV$ (\ce{CO2}) & Non \\
\hline
O & $-II$ & $-II$ (\ce{CaO}, \ce{CO2}) & Non \\
\hline
\end{tabularx}
\end{center}
(Calcul C dans \ce{CaCO3} : $(+II) + x + 3(-II) = 0 \Rightarrow x = +IV$).
Aucun n.o. ne change : \textbf{pas de réaction redox} (simple décomposition thermique).

\begin{exoresolu}[Combustion du méthane]
\textbf{Énoncé.} Le méthane \ce{CH4} brûle : $\ce{CH4 + 2 O2 -> CO2 + 2 H2O}$. Montrer que c'est une réaction redox. Identifier les acteurs.
\tcblower
\textbf{Solution.}

\textbf{1. n.o. réactifs :}
\ce{CH4} : H $= +I$ $\Rightarrow$ $x + 4(+I) = 0 \Rightarrow$ n.o.(C) $= -IV$.
\ce{O2} : n.o.(O) $= 0$.

\textbf{2. n.o. produits :}
\ce{CO2} : O $= -II$ $\Rightarrow$ $x + 2(-II) = 0 \Rightarrow$ n.o.(C) $= +IV$.
\ce{H2O} : H $= +I$, O $= -II$.

\textbf{3. Comparaison :}
\begin{align*}
\text{Carbone : } & -IV \longrightarrow +IV \quad (\text{augmentation de 8 : \textbf{oxydation}}) \\
\text{Oxygène : } & 0 \longrightarrow -II \quad (\text{diminution de 2 : \textbf{réduction}}) \\
\text{Hydrogène : } & +I \longrightarrow +I \quad (\text{inchangé})
\end{align*}

\textbf{Conclusion.} C'est une \textbf{réaction redox}. \ce{CH4} (C oxydé) est le \cle{réducteur} ; \ce{O2} (O réduit) est l'\cle{oxydant}. Cohérence : gain total n.o. C ($+8$) = perte totale n.o. O ($4 \times 2 = 8$).
\end{exoresolu}

\courssub{Équivalence avec la définition électronique}

\begin{propriete}[Lien n.o. $\leftrightarrow$ électrons]
\begin{itemize}
\item \textbf{Augmenter} son n.o. $\Leftrightarrow$ \textbf{perdre des électrons} (oxydation).
\item \textbf{Diminuer} son n.o. $\Leftrightarrow$ \textbf{gagner des électrons} (réduction).
\end{itemize}
La variation totale du n.o. de l'élément oxydé est égale, en valeur absolue, à la variation totale du n.o. de l'élément réduit : \textbf{conservation des électrons échangés}.
\end{propriete}

Vérification sur Mg : $\ce{Mg -> Mg^2+ + 2 e-}$ (perte 2 e$^-$, n.o. $0 \to +II$, $\Delta = +2$).
$\ce{O2 + 4 e- -> 2 O^2-}$ (gain 2 e$^-$ par atome O, n.o. $0 \to -II$, $\Delta = -2$).

\courssub{Arbre de décision pour le Probatoire}

\begin{popfigure}
\begin{center}
\begin{tikzpicture}[
  node distance=0.55cm and 1.1cm,
  qbox/.style={draw=popblue, thick, rounded corners, fill=popblueL, text width=4.8cm, align=center, font=\small, inner sep=5pt},
  abox/.style={draw=popdark, thick, rounded corners, fill=popcream, text width=4.8cm, align=center, font=\small, inner sep=5pt},
  ok/.style={draw=popgreen, thick, rounded corners, fill=popgreenL, text width=4.4cm, align=center, font=\small\bfseries, inner sep=5pt},
  ko/.style={draw=poporange, thick, rounded corners, fill=poporangeL, text width=4.4cm, align=center, font=\small\bfseries, inner sep=5pt},
  arr/.style={-{Stealth}, thick, popdark}
]
\node[qbox] (start) {Équation-bilan équilibrée};
\node[abox, below=of start] (calc) {Calculer le n.o. de \textbf{chaque élément} (réactifs \& produits)};
\node[qbox, below=of calc] (quest) {Au moins un n.o. a-t-il changé ?};
\node[ok, below left=0.7cm and -0.6cm of quest] (yes) {Réaction \textbf{REDOX} : n.o. $\nearrow$ = oxydation (réducteur) ; n.o. $\searrow$ = réduction (oxydant)};
\node[ko, below right=0.7cm and -0.6cm of quest] (no) {Pas redox (ex. : \ce{CaCO3 -> CaO + CO2})};
\draw[arr] (start) -- (calc);
\draw[arr] (calc) -- (quest);
\draw[arr] (quest) -| node[pos=0.25, above, font=\small\bfseries, popgreen]{OUI} (yes);
\draw[arr] (quest) -| node[pos=0.25, above, font=\small\bfseries, poporange]{NON} (no);
\end{tikzpicture}
\end{center}
\legende{Arbre de décision : la réaction est-elle une oxydoréduction ?}
\end{popfigure}

\courssub{Complément Probatoire : les réactions de dismutation}

\begin{definition}[Dismutation (ou disproportionation)]
Une \cle{dismutation} est une réaction redox où \textbf{un même élément}, présent dans une \textbf{seule espèce} de départ, est \textbf{à la fois oxydé et réduit} : une partie de ses atomes voit son n.o. augmenter, l'autre diminuer.
\end{definition}

Exemple classique : action du dichlore sur l'eau (désinfection eau de boisson) :
\[ \ce{Cl2 + H2O <=> HCl + HOCl} \]
Analyse n.o. du Cl :
\begin{itemize}
\item \ce{Cl2} : n.o.(Cl) $= 0$.
\item \ce{HCl} : n.o.(Cl) $= -I$ ($\searrow$ \textbf{réduction}).
\item \ce{HOCl} (acide hypochloreux) : H $= +I$, O $= -II$ $\Rightarrow$ $+I - II + x = 0 \Rightarrow$ n.o.(Cl) $= +I$ ($\nearrow$ \textbf{oxydation}).
\end{itemize}
Le chlore ($0$) est \textbf{à la fois oxydé ($+I$) et réduit ($-I$)} : c'est une dismutation. L'eau de Javel (marchés de Douala/Yaoundé) repose sur la réaction analogue avec la soude : $\ce{Cl2 + 2 NaOH -> NaCl + NaClO + H2O}$.

\begin{savaistu}[Pourquoi « oxydation » ?]
Le mot vient de l'\textbf{oxygène}. Au \textsc{xviii}\textsuperscript{e} siècle, Lavoisier montra que combustions et rouilles sont combinaisons avec ce gaz, qu'il nomma « oxygène » (« engendreur d'acides »). L'oxygène fut le \textbf{premier oxydant} étudié. On découvrit plus tard que $\ce{Cl2}$, $\ce{MnO4-}$, $\ce{Cr2O7^2-}$... jouent le même rôle sans contenir d'oxygène : le vocabulaire resta, le concept se généralisa (électrons, puis n.o.).
\end{savaistu}

\courssub{Équilibrer une équation redox grâce aux nombres d'oxydation}

La méthode de référence au Probatoire est la \textbf{méthode des demi-équations (ion-électron)}, qui utilise les n.o. pour identifier les couples redox et le nombre d'électrons échangés.

\begin{methode}[Méthode des demi-équations (milieu acide)]
\begin{enumerate}
\item \textbf{Identifier les couples} : calculer les n.o. ; repérer l'élément oxydé (n.o. $\nearrow$) et l'élément réduit (n.o. $\searrow$). Écrire les couples \ce{Ox/Red}.
\item \textbf{Écrire les deux demi-équations} (oxydation et réduction) sous forme \ce{Ox + n e- <=> Red} (ou inverse pour oxydation).
\item \textbf{Équilibrer chaque demi-équation} :
  \begin{itemize}
  \item Atomes autres que O et H.
  \item Atomes O en ajoutant \ce{H2O}.
  \item Atomes H en ajoutant \ce{H+}.
  \item Charges en ajoutant des \textbf{électrons \ce{e-}}.
  \end{itemize}
\item \textbf{Égaliser les électrons} : multiplier les demi-équations pour que le nombre d'\ce{e-} perdus = nombre d'\ce{e-} gagnés.
\item \textbf{Additionner} les deux demi-équations et simplifier (\ce{e-}, \ce{H+}, \ce{H2O} communs).
\item \textbf{Vérifier} : bilan atomes et charges.
\item \textbf{Milieu basique} : ajouter \ce{OH-} des deux côtés pour neutraliser les \ce{H+} restants ($\ce{H+ + OH- -> H2O}$).
\end{enumerate}
\end{methode}

\begin{exoresolu}[Dosage du fer(II) par le permanganate (milieu acide)]
\textbf{Énoncé.} Équilibrer la réaction entre l'ion permanganate \ce{MnO4-} et l'ion fer(II) \ce{Fe^2+} en milieu acide (\ce{H2SO4}).
\tcblower
\textbf{Solution.}

\textbf{1. Couples et n.o. :}
\ce{MnO4-} : Mn $= +VII$ ; \ce{Mn^2+} : Mn $= +II$ ($\searrow$ réduction). Couple \ce{MnO4-}/\ce{Mn^2+}.
\ce{Fe^2+} : Fe $= +II$ ; \ce{Fe^3+} : Fe $= +III$ ($\nearrow$ oxydation). Couple \ce{Fe^3+}/\ce{Fe^2+}.

\textbf{2. Demi-équations :}
\begin{itemize}
\item Réduction : $\ce{MnO4- -> Mn^2+}$
  \begin{itemize}
  \item Mn équilibré.
  \item 4 O $\to$ ajouter $4 \ce{H2O}$ à droite : $\ce{MnO4- -> Mn^2+ + 4 H2O}$.
  \item 8 H à gauche $\to$ ajouter $8 \ce{H+}$ : $\ce{MnO4- + 8 H+ -> Mn^2+ + 4 H2O}$.
  \item Charges : gauche $(-1) + 8 = +7$ ; droite $+2$. Ajouter $5 \ce{e-}$ à gauche : $\ce{MnO4- + 8 H+ + 5 e- -> Mn^2+ + 4 H2O}$.
  \end{itemize}
\item Oxydation : $\ce{Fe^2+ -> Fe^3+}$
  \begin{itemize}
  \item Fe équilibré. Pas de O, pas de H.
  \item Charges : gauche $+2$, droite $+3$. Ajouter $1 \ce{e-}$ à droite : $\ce{Fe^2+ -> Fe^3+ + e-}$.
  \end{itemize}
\end{itemize}

\textbf{3. Égaliser les électrons :} $5 \times$ (Oxydation) : $\ce{5 Fe^2+ -> 5 Fe^3+ + 5 e-}$.

\textbf{4. Additionner :}
\[ \ce{MnO4- + 8 H+ + 5 e- + 5 Fe^2+ -> Mn^2+ + 4 H2O + 5 Fe^3+ + 5 e-} \quad \text{(les 5 électrons s'annulent)} \]

\textbf{5. Équation-bilan finale :}
\[ \ce{MnO4- + 5 Fe^2+ + 8 H+ -> Mn^2+ + 5 Fe^3+ + 4 H2O} \]
\textbf{Vérification :} Atomes : Mn(1), O(4), Fe(5), H(8). Charges : G $(-1) + 5(+2) + 8(+1) = +15$ ; D $(+2) + 5(+3) = +17$ ? \textbf{Erreur !}
Recalcul charges G : $-1 + 10 + 8 = 17$. D : $2 + 15 = 17$. \textbf{Correct.}
\end{exoresolu}

\begin{attention}[Erreurs classiques à l'équilibrage]
\begin{itemize}
\item Oublier d'équilibrer les atomes O et H \textbf{avant} les charges.
\item Mettre les \ce{e-} du mauvais côté (réduction = \ce{e-} à gauche ; oxydation = \ce{e-} à droite).
\item Ne pas multiplier la demi-équation d'oxydation pour égaler les électrons.
\item Oublier de simplifier les espèces communes (\ce{H2O}, \ce{H+}, \ce{e-}) après addition.
\end{itemize}
\end{attention}

\courssub{Applications industrielles : sidérurgie et aluminothermie}

\courssub{Sidérurgie : le haut fourneau (CIMENCAM, Figuil / SONARA contextuel)}

La production de la fonte (fer impur) à partir de minerai de fer (hématite \ce{Fe2O3}, magnétite \ce{Fe3O4}) se fait dans un \textbf{haut fourneau} (hauteur $\sim \SI{30}{m}$). Le \textbf{coke} (C) est le réducteur (et combustible). L'air chaud (\ce{O2}) est soufflé par le bas (tuères).

\begin{center}
\begin{tikzpicture}[scale=0.7, every node/.style={font=\small}]
% Structure fourneau
\draw[thick, popdark] (0,0) -- (0,8) -- (1,9) -- (4,9) -- (5,8) -- (5,0);
\draw[thick, popdark] (1,9) -- (1,10) -- (4,10) -- (4,9); % gueulard
% Zones
\fill[poporangeL] (0.2,0.2) rectangle (4.8,2.5); node[popdark] at (2.5,1.3) {\textbf{Zone de combustion} (\SI{1800}{\celsius})};
\fill[popgoldL] (0.2,2.5) rectangle (4.8,5.5); node[popdark] at (2.5,4) {\textbf{Zone de réduction} (\SI{800}{\celsius}--\SI{1000}{\celsius})};
\fill[popgreenL] (0.2,5.5) rectangle (4.8,8); node[popdark] at (2.5,6.7) {\textbf{Zone de chargement} (minerai + coke + fondant)};
% Flèches air
\draw[-{Stealth}, thick, popblue] (-1,1) -- (0.2,1) node[left] {Air chaud (\ce{O2})};
% Flèches fonte/laitier
\draw[-{Stealth}, thick, poporange] (2.5,0) -- (2.5,-0.5) node[below] {Fonte \ce{Fe} + Laitier};
% Réactions texte
\node[align=left, font=\tiny, popdark] at (6.5, 1.5) {
\textbf{Bas (combustion) :}\\
$\ce{C + O2 -> CO2}$ $\Delta H \ll 0$\\
$\ce{CO2 + C -> 2 CO}$ (endothermique)\\
};
\node[align=left, font=\tiny, popdark] at (6.5, 4) {
\textbf{Milieu (réduction) :}\\
$\ce{Fe2O3 + 3 CO -> 2 Fe + 3 CO2}$\\
$\ce{Fe2O3 + 3 H2 -> 2 Fe + 3 H2O}$\\
$\ce{FeO + CO -> Fe + CO2}$\\
};
\node[align=left, font=\tiny, popdark] at (6.5, 6.5) {
\textbf{Haut :}\\
Chargement : \ce{Fe2O3}, C, \ce{CaCO3}\\
$\ce{CaCO3 -> CaO + CO2}$\\
$\ce{CaO + SiO2 -> CaSiO3}$ (laitier)
};
\end{tikzpicture}
\end{center}

\begin{aretenir}[Réactions clés du haut fourneau]
\begin{itemize}
\item \textbf{Combustion (bas) :} $\ce{C + O2 -> CO2}$ (exothermique, chaleur) ; $\ce{CO2 + C <=> 2 CO}$ (endothermique, monte).
\item \textbf{Réduction (milieu) :} $\ce{Fe2O3 + 3 CO -> 2 Fe + 3 CO2}$ (principal). $\ce{CO}$ est le réducteur gazeux (n.o. C : $+II \to +IV$). $\ce{Fe2O3}$ est l'oxydant (Fe $+III \to 0$).
\item \textbf{Fondant :} Calcaire $\ce{CaCO3} \to \ce{CaO + CO2}$ ; $\ce{CaO + SiO2 -> CaSiO3}$ (laitier, flotte sur la fonte).
\end{itemize}
\end{aretenir}

\courssub{Aluminothermie (soudure rails, chimie du feu)}

Réaction très exothermique entre l'oxyde de fer(III) et l'aluminium en poudre (mélange \og thermite \fg{}). Initiée par un amorce (magnésium + perchlorate ou fil chauffant). Température $\sim \SI{2500}{\celsius}$ $\to$ fer fondu.

\[ \ce{Fe2O3(s) + 2 Al(s) -> Al2O3(s) + 2 Fe(l)} \qquad \Delta H^\circ_r = \SI{-851,5}{kJ.mol^{-1}} \]

Analyse n.o. : Fe $+III \to 0$ (réduction, \ce{Fe2O3} oxydant) ; Al $0 \to +III$ (oxydation, \ce{Al} réducteur).
\textbf{Applications :} Soudure des rails (chemins de fer Camrail), réparation pièces moulées, grenades incendiaires, allumage propergols.

\begin{savaistu}[Pourquoi l'aluminium ?]
L'Al a une affinité pour l'oxygène plus forte que le Fe (diagramme d'Ellingham). La réaction est spontanée et violemment exothermique. L'\ce{Al2O3} (corindon) formé est très stable.
\end{savaistu}

\courssub{Préparation industrielle des acides sulfurique et nitrique}

Deux piliers de l'industrie chimique lourde (engrais, détergents, explosifs, raffinage SONARA...).

\courssub{Acide sulfurique \ce{H2SO4} : Procédé Contact (BASF, 1915)}

\begin{enumerate}
\item \textbf{Production du \ce{SO2} :} Combustion du soufre (ou roasting de pyrite \ce{FeS2}).
\[ \ce{S(s) + O2(g) -> SO2(g)} \quad \text{ou} \quad \ce{4 FeS2 + 11 O2 -> 2 Fe2O3 + 8 SO2} \]
\item \textbf{Oxydation du \ce{SO2} en \ce{SO3} (étape clé, équilibre) :}
\[ \ce{2 SO2(g) + O2(g) <=>[\text{catalyseur } \ce{V2O5}] [\SI{450}{\celsius}] 2 SO3(g)} \quad \Delta H^\circ_r = \SI{-197}{kJ.mol^{-1}} \]
\begin{itemize}
\item Catalyseur : pentoxyde de vanadium \ce{V2O5} sur support silice (remplace platine).
\item Conditions : \SI{450}{\celsius} (compromis cinétique/thermodynamique), pression $\sim \SI{1}{bar}$.
\item Rendement $\sim 98\%$ (excès d'\ce{O2}, retrait \ce{SO3}).
\end{itemize}
\item \textbf{Absorption du \ce{SO3} :} Pas directement dans l'eau (brouillard d'acide). On utilise de l'\textbf{oléum} (\ce{H2S2O7}).
\[ \ce{SO3 + H2SO4 -> H2S2O7} \quad \text{(oléum)} \]
\[ \ce{H2S2O7 + H2O -> 2 H2SO4} \]
\end{enumerate}

\begin{aretenir}[Bilan global Contact]
\[ \ce{S + 3/2 O2 + H2O -> H2SO4} \]
Variations n.o. : S ($0 \to +IV \to +VI$) oxydé ; O ($0 \to -II$) réduit. \ce{O2} oxydant, S/\ce{SO2} réducteur.
\end{aretenir}

\courssub{Acide nitrique \ce{HNO3} : Procédé Ostwald (1902)}

Partant de l'ammoniac \ce{NH3} (synthèse Haber-Bosch : $\ce{N2 + 3 H2 <=> 2 NH3}$).

\begin{enumerate}
\item \textbf{Oxydation catalytique de l'ammoniac (étape clé) :}
\[ \ce{4 NH3(g) + 5 O2(g) ->[\text{Pt/Rh } \SI{900}{\celsius}] 4 NO(g) + 6 H2O(g)} \quad \Delta H^\circ_r = \SI{-905}{kJ.mol^{-1}} \]
\begin{itemize}
\item Catalyseur : gaze \ce{Pt}/\ce{Rh} (90/10) très fin.
\item Très exothermique : auto-entretenu, récupération chaleur (vapeur).
\item Sélectivité vers \ce{NO} (éviter $\ce{N2}$, $\ce{N2O}$).
\end{itemize}
\item \textbf{Oxydation du \ce{NO} en \ce{NO2} (refroidissement) :}
\[ \ce{2 NO(g) + O2(g) -> 2 NO2(g)} \quad \text{(rapide, exothermique, favorisé à basse T)} \]
\item \textbf{Absorption du \ce{NO2} dans l'eau :}
\[ \ce{3 NO2(g) + H2O(l) -> 2 HNO3(aq) + NO(g)} \]
Le \ce{NO} non réagi est recyclé (réoxydé en \ce{NO2} par l'air).
\end{enumerate}

\begin{aretenir}[Bilan global Ostwald]
\[ \ce{NH3 + 2 O2 -> HNO3 + H2O} \]
Variations n.o. : N ($-III \to +II \to +IV \to +V$) oxydé en 3 étapes ; O ($0 \to -II$) réduit. \ce{O2} oxydant, \ce{NH3} réducteur.
\end{aretenir}

\begin{popfigure}
\begin{center}
\begin{tikzpicture}[node distance=1.2cm, >=Stealth, thick,
  box/.style={draw=popblue, fill=popblueL, rounded corners, minimum width=2.5cm, minimum height=1cm, align=center, font=\small},
  arr/.style={->, popdark, thick},
  elabel/.style={above, font=\tiny, align=center}]
% Contact
\node[box] (S) {Soufre \ce{S}};
\node[box, right=of S] (SO2) {\ce{SO2}};
\node[box, right=of SO2] (SO3) {\ce{SO3}};
\node[box, right=of SO3] (Oleum) {Ol\'eum \ce{H2S2O7}};
\node[box, right=of Oleum] (H2SO4) {\ce{H2SO4}};
\draw[arr] (S) -- node[elabel] {\ce{+ O2}} (SO2);
\draw[arr] (SO2) -- node[elabel] {\ce{+ 1/2 O2}\\ $\ce{V2O5}$, $450^\circ\mathrm{C}$} (SO3);
\draw[arr] (SO3) -- node[elabel] {\ce{+ H2SO4}} (Oleum);
\draw[arr] (Oleum) -- node[elabel] {\ce{+ H2O}} (H2SO4);
% Ostwald (below)
\node[box, below=2cm of S] (NH3) {Ammoniac \ce{NH3}};
\node[box, right=of NH3] (NO) {\ce{NO}};
\node[box, right=of NO] (NO2) {\ce{NO2}};
\node[box, right=of NO2] (HNO3) {\ce{HNO3}};
\draw[arr] (NH3) -- node[elabel] {\ce{+ 5/4 O2}\\ $\ce{Pt/Rh}$, $900^\circ\mathrm{C}$} (NO);
\draw[arr] (NO) -- node[elabel] {\ce{+ 1/2 O2}} (NO2);
\draw[arr] (NO2) -- node[elabel] {\ce{+ H2O}\\ recyclage \ce{NO}} (HNO3);
% Labels
\node[popdark, font=\bfseries] at ($(S)!0.5!(SO2)$) {Procédés Contact};
\node[popdark, font=\bfseries] at ($(NH3)!0.5!(NO)$) {Procédés Ostwald};
\end{tikzpicture}
\end{center}
\legende{Synthèse des chaînes industrielles : Acide sulfurique (haut) et Acide nitrique (bas).}
\end{popfigure}

\begin{exercice}[Entraînement Probatoire]
\begin{enumerate}
\item \textbf{Équilibrage par demi-équations (milieu acide) :} $\ce{Cr2O7^2- + Fe^2+ + H+ -> Cr^3+ + Fe^3+ + H2O}$.
\item \textbf{Reconnaissance redox :} Pour chaque réaction, dire si c'est redox (justifier par n.o.) et identifier oxydant/réducteur le cas échéant.
  \begin{enumerate}
  \item $\ce{CuO + H2 -> Cu + H2O}$ (réduction minerai cuivre).
  \item $\ce{2 KClO3 ->[MnO2] 2 KCl + 3 O2}$ (préparation $\ce{O2}$ au labo).
  \item $\ce{AgNO3 + NaCl -> AgCl v + NaNO3}$.
  \end{enumerate}
\item \textbf{Sidérurgie :} Dans le haut fourneau, la réduction de \ce{FeO} par \ce{CO} suit : $\ce{FeO + CO -> Fe + CO2}$. Calculer les n.o. et vérifier le bilan électronique.
\item \textbf{Acide nitrique :} Écrire les 3 équations-bilan du procédé Ostwald. Indiquer pour chaque étape l'élément oxydé, l'élément réduit, l'oxydant et le réducteur.
\end{enumerate}
\end{exercice}

\begin{corrige}[Exercice]
\textbf{1.} Couples : \ce{Cr2O7^2-}/\ce{Cr^3+} (Cr $+VI \to +III$, gain $3\ce{e-}$ par Cr, $6\ce{e-}$ total) et \ce{Fe^3+}/\ce{Fe^2+} (Fe $+II \to +III$, perte $1\ce{e-}$).
Demi-éq. réd. : $\ce{Cr2O7^2- + 14 H+ + 6 e- -> 2 Cr^3+ + 7 H2O}$.
Demi-éq. ox. : $\ce{6 Fe^2+ -> 6 Fe^3+ + 6 e-}$.
Bilan : $\ce{Cr2O7^2- + 6 Fe^2+ + 14 H+ -> 2 Cr^3+ + 6 Fe^3+ + 7 H2O}$.

\textbf{2.}
a) $\ce{CuO + H2 -> Cu + H2O}$. Cu $+II \to 0$ ($\searrow$ réd) ; H $0 \to +I$ ($\nearrow$ ox). \textbf{Redox}. Oxydant : \ce{CuO} ; Réducteur : \ce{H2}.
b) $\ce{2 KClO3 -> 2 KCl + 3 O2}$. Cl $+V \to -I$ ($\searrow$ réd) ; O $-II \to 0$ ($\nearrow$ ox). \textbf{Redox} (dismutation interne du \ce{ClO3-}). Oxydant/Réducteur : \ce{KClO3} (lui-même).
c) $\ce{AgNO3 + NaCl -> AgCl + NaNO3}$. Ag $+I$, N $+V$, O $-II$, Na $+I$, Cl $-I$ : \textbf{tous inchangés}. \textbf{Pas redox} (précipitation).

\textbf{3.} \ce{FeO} : Fe $+II$, O $-II$. \ce{CO} : C $+II$, O $-II$. Produits : Fe $0$, \ce{CO2} : C $+IV$.
Fe $+II \to 0$ ($\searrow$ réd, gain $2\ce{e-}$). C $+II \to +IV$ ($\nearrow$ ox, perte $2\ce{e-}$). Bilan électronique équilibré ($2\ce{e-}$).

\textbf{4.}
(1) $\ce{4 NH3 + 5 O2 -> 4 NO + 6 H2O}$ : N $-III \to +II$ (ox, \ce{NH3} réd) ; O $0 \to -II$ (réd, \ce{O2} ox).
(2) $\ce{2 NO + O2 -> 2 NO2}$ : N $+II \to +IV$ (ox, \ce{NO} réd) ; O $0 \to -II$ (réd, \ce{O2} ox).
(3) $\ce{3 NO2 + H2O -> 2 HNO3 + NO}$ : N $+IV \to +V$ (ox, 2 atomes) et $+IV \to +II$ (réd, 1 atome). \textbf{Dismutation du \ce{NO2}}.
\end{corrige}

\begin{aretenir}[Synthèse de la leçon 10]
\begin{itemize}
\item \textbf{Voie sèche :} Combustions $\ce{Mg}$, $\ce{H2}$, $\ce{C}$ dans $\ce{O2}$/$\ce{Cl2}$ $\to$ transferts d'électrons nets.
\item \textbf{Nombre d'oxydation :} Charge formelle (règles conventionnelles). Outil universel.
\item \textbf{Critère redox :} Variation d'au moins un n.o. $\Leftrightarrow$ réaction redox.
\item \textbf{Analyse :} n.o. $\nearrow$ = oxydation (réducteur) ; n.o. $\searrow$ = réduction (oxydant).
\item \textbf{Équilibrage :} Méthode des demi-équations (identifier couples par n.o., équilibrer atomes/charges/électrons, additionner).
\item \textbf{Dismutation :} Même élément oxydé et réduit ($\ce{Cl2/H2O}$, $\ce{H2O2}$, $\ce{NO2/H2O}$).
\item \textbf{Industrie :}
  \begin{itemize}
  \item Sidérurgie : Haut fourneau, \ce{Fe2O3} réduit par \ce{CO} (venant du coke).
  \item Aluminothermie : \ce{Fe2O3 + 2 Al -> 2 Fe + Al2O3} (soudure rails, $\Delta H \ll 0$).
  \item Acide sulfurique : Contact ($\ce{S -> SO2 -> SO3 -> H2SO4}$, cat. \ce{V2O5}).
  \item Acide nitrique : Ostwald ($\ce{NH3 -> NO -> NO2 -> HNO3}$, cat. \ce{Pt/Rh}).
  \end{itemize}
\end{itemize}
\end{aretenir}

\coursec{Équilibrer une équation redox grâce aux nombres d'oxydation}

Dans la section précédente, tu as appris à \emph{reconnaître} une réaction d'oxydoréduction en suivant les variations des nombres d'oxydation : si au moins un élément voit son n.o. augmenter (oxydation) pendant qu'un autre voit le sien diminuer (réduction), la réaction est bien une réaction redox. Nous allons maintenant franchir une étape supplémentaire : utiliser ces mêmes variations de n.o. pour \emph{équilibrer} les équations-bilan. C'est un outil d'une redoutable efficacité, particulièrement pour les réactions par voie sèche (combustions, métallurgie) où la méthode des demi-équations électroniques, pensée pour les solutions aqueuses, devient maladroite.

\courssub{Le principe : la conservation des électrons échangés}

Imaginons un marché de Mokolo où un commerçant échange des marchandises : ce qu'il cède d'un côté doit exactement correspondre à ce qu'il reçoit de l'autre, sinon l'échange est déséquilibré. En oxydoréduction, la « marchandise » échangée, ce sont les \cle{électrons}. L'élément qui s'oxyde \emph{perd} des électrons (son n.o. augmente), l'élément qui se réduit \emph{gagne} des électrons (son n.o. diminue). Comme les électrons ne peuvent ni apparaître ni disparaître, le nombre d'électrons perdus doit être exactement égal au nombre d'électrons gagnés.

\begin{propriete}[Principe fondamental de l'équilibrage par les n.o.]
Dans toute réaction d'oxydoréduction, la \cle{somme des augmentations} des nombres d'oxydation (tous les atomes oxydés confondus) est égale à la \cle{somme des diminutions} des nombres d'oxydation (tous les atomes réduits confondus) :
\[
\sum \Delta(\text{n.o.})_{\text{augmentations}} \;=\; \sum \left|\Delta(\text{n.o.})\right|_{\text{diminutions}}
\]
Cette égalité traduit simplement la \cle{conservation des électrons} échangés au cours de la réaction.
\end{propriete}

Toute la méthode découle de cette unique phrase : on va choisir les coefficients stœchiométriques de manière que les électrons « partis » d'un côté soient exactement « récupérés » de l'autre.

\courssub{La méthode en quatre étapes}

\begin{methode}[Équilibrer une équation redox par les nombres d'oxydation]
\begin{enumerate}
\item \textbf{Écrire le squelette de l'équation} (formules des réactifs et des produits) et \textbf{calculer les n.o.} de tous les éléments ; repérer ceux dont le n.o. \emph{change}.
\item \textbf{Déterminer la variation de n.o.} pour l'élément oxydé (augmentation) et pour l'élément réduit (diminution), \emph{par atome}, puis \emph{par formule} en multipliant par le nombre d'atomes de l'élément dans la formule.
\item \textbf{Ajuster les coefficients} des espèces redox pour égaliser les variations totales : on utilise le \cle{plus petit multiple commun} (PPMC) des deux variations et on « croise » les coefficients.
\item \textbf{Équilibrer les autres atomes} (ceux qui ne changent pas de n.o. : O, H, spectateurs\ldots), puis \textbf{vérifier} la conservation des masses (et des charges si des ions interviennent).
\end{enumerate}
\end{methode}

\begin{popfigure}
\begin{center}
\begin{tikzpicture}[
  font=\small,
  etape/.style={rectangle, rounded corners=4pt, draw=popblue, fill=popblueL, thick,
                text width=0.82\linewidth, align=left, inner sep=6pt},
  fleche/.style={-{Stealth[length=3mm]}, thick, popdark}
]
\node[etape] (e1) {\textbf{Étape 1 —} Écrire le squelette de l'équation et calculer les n.o. de tous les éléments ; repérer ceux qui varient.};
\node[etape, below=0.55cm of e1] (e2) {\textbf{Étape 2 —} Calculer la variation de n.o. de l'élément oxydé ($\uparrow$) et de l'élément réduit ($\downarrow$), par atome puis par formule.};
\node[etape, below=0.55cm of e2] (e3) {\textbf{Étape 3 —} Égaliser les variations totales grâce au PPMC : croiser les coefficients devant l'oxydant et le réducteur.};
\node[etape, below=0.55cm of e3] (e4) {\textbf{Étape 4 —} Équilibrer les autres atomes (O, H\ldots) puis vérifier la conservation des masses et des charges.};
\draw[fleche] (e1) -- (e2);
\draw[fleche] (e2) -- (e3);
\draw[fleche] (e3) -- (e4);
\end{tikzpicture}
\end{center}
\legende{Organigramme des quatre étapes de la méthode d'équilibrage par les nombres d'oxydation.}
\end{popfigure}

\begin{attention}[La variation se calcule par atome, puis par formule]
Piège classique : quand une formule contient \emph{plusieurs} atomes de l'élément qui change de n.o., il faut multiplier la variation par le nombre d'atomes concernés. Ainsi, dans \ce{O2}, chaque atome d'oxygène passe de 0 à $-$II, soit une variation de 2 \emph{par atome}, donc de $2 \times 2 = 4$ \emph{par molécule} \ce{O2}. De même, dans \ce{Fe2O3}, les deux atomes de fer passent chacun de $+$III à 0 : la diminution totale est $2 \times 3 = 6$ par formule \ce{Fe2O3}. Oublier ce facteur est l'erreur la plus fréquente au Probatoire !
\end{attention}

\courssub{Exemple guidé 1 : l'aluminothermie}

Considérons la réaction utilisée pour souder les rails de chemin de fer (ligne Douala--Ngaoundéré) : l'aluminium réduit l'oxyde de fer(III). Le squelette de l'équation est :
\[
\ce{Al + Fe2O3 -> Al2O3 + Fe}
\]

\textbf{Étape 1 — Calcul des n.o.} L'aluminium est un corps simple : n.o.(Al) $= 0$ ; il devient $+$III dans \ce{Al2O3}. Le fer est $+$III dans \ce{Fe2O3} (car n.o.(O) $= -$II : $2x + 3\times(-2) = 0$ donne $x = +3$) ; il devient 0 dans le corps simple \ce{Fe}. L'oxygène reste à $-$II : il ne participe pas au transfert d'électrons.

\textbf{Étape 2 — Variations.} Al : de 0 à $+$III, augmentation de \textbf{3} par atome. Fe : de $+$III à 0, diminution de \textbf{3} par atome, soit $2 \times 3 = 6$ par formule \ce{Fe2O3}.

\textbf{Étape 3 — Égalisation.} Augmentation totale $=$ diminution totale. Prenons 2 atomes d'Al : l'augmentation vaut $2 \times 3 = 6$, exactement la diminution des 2 atomes de Fe contenus dans une formule \ce{Fe2O3}. On place donc le coefficient \textbf{2} devant Al, et \textbf{1} devant \ce{Fe2O3} (ce qui impose 2 devant Fe).

\textbf{Étape 4 — Autres atomes et vérification.} Il reste à équilibrer l'aluminium et l'oxygène dans les produits : 2 Al donnent 1 \ce{Al2O3}, qui contient 3 O, fournis par 1 \ce{Fe2O3}. L'équation équilibrée s'écrit :
\[
\ce{2Al + Fe2O3 ->[\text{haute température}] Al2O3 + 2Fe}
\]
Vérification : 2 Al, 2 Fe et 3 O de chaque côté. La réaction, amorcée par un ruban de magnésium enflammé, est fortement exothermique : le fer obtenu est \emph{liquide} (température voisine de \SI{2000}{\celsius}), ce qui permet de souder les rails.

\begin{popfigure}
\begin{center}
\begin{tikzpicture}[font=\small]
% Axe vertical
\draw[-{Stealth}, thick, popdark] (0,-3.2) -- (0,3.6) node[above] {variation de n.o.};
\draw[popmuted, dashed] (-0.2,0) -- (7.6,0);
% Barre Al (montante, +3 par atome, x2 = 6)
\fill[poporange!80] (0.6,0) rectangle (1.8,3);
\node[below, align=center] at (1.2,0) {2 Al\\ $0 \rightarrow +\mathrm{III}$};
\node[above, poporange!60!black, font=\bfseries] at (1.2,3) {$+6$};
% Barre Fe (descendante, -3 par atome, x2 = -6)
\fill[popblue!75] (3.2,0) rectangle (4.4,-3);
\node[above, align=center] at (3.8,0) {2 Fe\\ $+\mathrm{III} \rightarrow 0$};
\node[below, popblue!60!black, font=\bfseries] at (3.8,-3) {$-6$};
% Egalité
\node[font=\Large\bfseries, popdark] at (5.6,0) {$\Rightarrow$};
\node[align=center, popgreen!50!black, font=\bfseries] at (6.9,0.4) {$|+6| = |-6|$};
\node[align=center, popmuted] at (6.9,-0.6) {électrons\\ conservés};
\end{tikzpicture}
\end{center}
\legende{Bilan des variations de n.o. dans l'aluminothermie : la hausse totale (2 Al, $+6$) compense exactement la baisse totale (2 Fe, $-6$).}
\end{popfigure}

\courssub{Exemple guidé 2 : l'oxydation catalytique de l'ammoniac}

Cette réaction est la première étape de la préparation industrielle de l'acide nitrique (procédé Ostwald, utilisé par les engrais chimiques). Le squelette est :
\[
\ce{NH3 + O2 -> NO + H2O}
\]

\textbf{Étape 1 — Calcul des n.o.} Dans \ce{NH3} : n.o.(N) $= -$III (car $x + 3\times(+1) = 0$). Dans \ce{NO} : n.o.(N) $= +$II. Dans \ce{O2} : n.o.(O) $= 0$ ; dans \ce{NO} et \ce{H2O} : n.o.(O) $= -$II. L'hydrogène reste à $+$I.

\textbf{Étape 2 — Variations.} N : de $-$III à $+$II, augmentation de \textbf{5} par atome. O : de 0 à $-$II, diminution de 2 par atome, soit \textbf{4} par molécule \ce{O2}.

\textbf{Étape 3 — Égalisation.} PPMC de 5 et 4 : c'est 20. Il faut donc 4 atomes de N (augmentation $4 \times 5 = 20$) et 5 molécules \ce{O2} (diminution $5 \times 4 = 20$) :
\[
\ce{4NH3 + 5O2 -> NO + H2O} \quad \text{(coefficients redox placés)}
\]

\textbf{Étape 4 — Autres atomes et vérification.} 4 N imposent 4 \ce{NO}. L'hydrogène : $4 \times 3 = 12$ H à gauche, donc 6 \ce{H2O}. Comptons l'oxygène à droite : $4 + 6 = 10$ O, exactement ce qu'apportent 5 \ce{O2}. L'équation équilibrée est :
\[
\ce{4NH3 + 5O2 ->[Pt,\ \SI{850}{\celsius}] 4NO + 6H2O}
\]
Vérification complète : 4 N, 12 H et 10 O de chaque côté. La réaction s'effectue sur une toile de \cle{platine} (catalyseur) portée à environ \SI{850}{\celsius} ; elle est exothermique, ce qui entretient la température sans chauffage extérieur une fois amorcée.

\begin{attention}[Un réactif peut jouer un double rôle]
Dans certaines réactions, un même réactif intervient à la fois comme \cle{oxydant} et comme \emph{simple fournisseur d'ions} (rôle acidifiant ou complexant). Seule une \emph{partie} de ce réactif change de n.o. ! Par exemple, dans l'action de l'acide nitrique sur le cuivre, une partie des ions nitrate est réduite en \ce{NO}, mais l'autre partie reste \ce{NO3-} et sert simplement à former le sel \ce{Cu(NO3)2}. Il ne faut donc appliquer l'égalisation des variations \emph{qu'à la fraction qui change de n.o.}, puis compléter les coefficients à l'étape 4.
\end{attention}

\begin{exoresolu}[Action de l'acide nitrique dilué sur le cuivre]
Énoncé : un morceau de cuivre (fils électriques récupérés) est plongé dans de l'acide nitrique dilué. Il se forme du nitrate de cuivre(II), un dégagement de monoxyde d'azote \ce{NO} (gaz incolore qui brunit à l'air) et de l'eau. Équilibrer l'équation-bilan par la méthode des nombres d'oxydation :
\[
\ce{Cu + HNO3 -> Cu(NO3)2 + NO + H2O}
\]
\tcblower
\textbf{Solution détaillée.}

\textbf{Étape 1 — n.o. des éléments.} Cu : 0 (corps simple) $\rightarrow$ $+$II dans \ce{Cu(NO3)2}. N : $+$V dans \ce{HNO3} $\rightarrow$ $+$II dans \ce{NO}. Attention : dans \ce{Cu(NO3)2}, l'azote des nitrates est \emph{resté} à $+$V ; seule la partie de \ce{HNO3} transformée en \ce{NO} est réduite.

\textbf{Étape 2 — Variations.} Cu : augmentation de \textbf{2} par atome. N réduit : diminution de \textbf{3} par atome.

\textbf{Étape 3 — Égalisation.} PPMC de 2 et 3 : c'est 6. Il faut 3 Cu ($3 \times 2 = 6$ électrons cédés) et 2 N réduits ($2 \times 3 = 6$ électrons gagnés), donc 2 \ce{HNO3} \emph{en tant qu'oxydant} :
\[
\ce{3Cu + 2HNO3 -> 3Cu(NO3)2 + 2NO + H2O} \quad \text{(partie redox)}
\]

\textbf{Étape 4 — Compléter.} Les 3 \ce{Cu(NO3)2} contiennent $3 \times 2 = 6$ ions nitrate supplémentaires : il faut donc 6 \ce{HNO3} de plus pour leur rôle \emph{acidifiant}, soit $2 + 6 = 8$ \ce{HNO3} au total. L'hydrogène : 8 H à gauche imposent 4 \ce{H2O}. Vérifions l'oxygène : à gauche $8 \times 3 = 24$ O ; à droite $6 \times 3 + 2 + 4 = 24$ O. L'équation équilibrée est :
\[
\ce{3Cu + 8HNO3 -> 3Cu(NO3)2 + 2NO + 4H2O}
\]
Vérification finale : 3 Cu, 8 H, 8 N et 24 O de chaque côté. La réaction, lente à froid, s'accélère à chaud ; le gaz \ce{NO} s'oxyde au contact de l'air en \ce{NO2} roux, toxique : manipulation sous hotte !
\end{exoresolu}

\courssub{Vérification systématique et lien avec la méthode des demi-équations}

Une équation redox équilibrée doit passer le \emph{crash-test} final : on compte, un par un, les atomes de \textbf{chaque} élément des deux côtés de la flèche, et si des ions interviennent, on vérifie aussi la conservation de la \cle{charge électrique} totale. Ce contrôle prend trente secondes et élimine la quasi-totalité des erreurs.

Tu remarqueras que cette méthode et celle des demi-équations électroniques (vue pour les réactions en solution aqueuse) reposent sur le \emph{même} principe : la conservation des électrons. Elles conduisent donc nécessairement au \emph{même} résultat. La méthode des n.o. est simplement une comptabilité plus directe, qui évite d'écrire les demi-équations.

\begin{center}
\begin{tabularx}{\linewidth}{|l|X|X|}
\hline
\rowcolor{popblueL} \textbf{Critère} & \textbf{Demi-équations} & \textbf{Nombres d'oxydation} \\
\hline
Cadre privilégié & Réactions ioniques en solution & Réactions sans ions (voie sèche, combustions) \\
\hline
Outil & Électrons écrits explicitement & Variations de n.o. (électrons implicites) \\
\hline
Espèces \ce{H2O}, \ce{H+} & À équilibrer soigneusement & Souvent absentes du problème \\
\hline
Résultat & \multicolumn{2}{c|}{Identique : même équation-bilan équilibrée} \\
\hline
\end{tabularx}
\end{center}

\begin{savaistu}[La méthode préférée des métallurgistes]
Pour les réactions sans ions — combustions, grillage des minerais, réduction des oxydes dans les hauts fourneaux — la méthode des nombres d'oxydation est nettement plus rapide que celle des demi-équations : impossible, par exemple, d'écrire une demi-équation électronique « propre » pour la combustion du carbone \ce{C + O2 -> CO2}, car il n'y a ni ions, ni solvant. Les ingénieurs de la sidérurgie (CIMENCAM, fonderies) raisonnent quotidiennement en variations de n.o. pour ajuster les charges de coke et de minerai. C'est aussi cette logique qui sous-tend le procédé Ostwald (acide nitrique) et le procédé de contact (acide sulfurique) que nous étudierons dans la dernière section de cette leçon.
\end{savaistu}

\begin{aretenir}[L'essentiel de la section]
\begin{itemize}
\item Principe : la somme des augmentations de n.o. égale la somme des diminutions de n.o. (conservation des électrons).
\item Méthode en 4 étapes : calculer les n.o. ; déterminer les variations (par atome \emph{puis} par formule) ; égaliser grâce au PPMC en croisant les coefficients ; équilibrer les autres atomes et vérifier.
\item Attention aux formules contenant plusieurs atomes de l'élément (\ce{O2} : variation 4 ; \ce{Fe2O3} : variation 6) et aux réactifs à double rôle (comme \ce{HNO3}).
\item Toujours terminer par le comptage exhaustif des atomes (et des charges).
\item Cette méthode et celle des demi-équations donnent le même résultat ; elle est particulièrement adaptée aux réactions par voie sèche.
\end{itemize}
\end{aretenir}

\begin{exercice}[Pour t'entraîner]
Équilibre par la méthode des nombres d'oxydation les équations suivantes :
\begin{enumerate}
\item $\ce{Fe2O3 + CO -> Fe + CO2}$ (réaction du haut fourneau) ;
\item $\ce{NH3 + CuO -> N2 + Cu + H2O}$ ;
\item $\ce{SO2 + O2 -> SO3}$ (étape clé du procédé de contact).
\end{enumerate}
\end{exercice}

\begin{corrige}[Exercice : équilibrages par les n.o.]
\begin{enumerate}
\item Fe : de $+$III à 0, diminution de 3 par atome, soit 6 par \ce{Fe2O3}. C : de $+$II à $+$IV, augmentation de 2 par atome. PPMC $= 6$ : il faut 3 CO. D'où $\ce{Fe2O3 + 3CO -> 2Fe + 3CO2}$.
\item N : de $-$III à 0, augmentation de 3 par atome, soit 6 pour former 1 \ce{N2}. Cu : de $+$II à 0, diminution de 2. PPMC $= 6$ : il faut 3 CuO. D'où $\ce{2NH3 + 3CuO -> N2 + 3Cu + 3H2O}$.
\item S : de $+$IV à $+$VI, augmentation de 2. O : diminution de 4 par \ce{O2}. PPMC $= 4$ : il faut 2 \ce{SO2} pour 1 \ce{O2}. D'où $\ce{2SO2 + O2 -> 2SO3}$ (catalyseur \ce{V2O5}, environ \SI{450}{\celsius}).
\end{enumerate}
\end{corrige}

\coursec{Applications industrielles : sidérurgie et aluminothermie}

Dans la section précédente, tu as appris à équilibrer une équation d'oxydoréduction grâce aux nombres d'oxydation. À quoi cela sert-il concrètement ? Réponse : à comprendre — et à piloter — certaines des plus grandes industries du monde. Chaque rail de chemin de fer, chaque poutre d'immeuble, chaque carrosserie de voiture provient d'une réaction redox : la \cle{réduction d'un oxyde métallique}. C'est toute la métallurgie. Dans cette section, nous allons voir deux procédés emblématiques : la \cle{sidérurgie} (production du fer au haut fourneau) et l'\cle{aluminothermie} (réduction d'oxydes métalliques par l'aluminium). Tu constateras que les nombres d'oxydation y sont un outil d'analyse irremplaçable.

\courssub{La sidérurgie : du minerai au fer}

\begin{definition}[Sidérurgie]
La \cle{sidérurgie} est l'ensemble des procédés industriels d'élaboration du \textbf{fer} et de ses dérivés (\textbf{fonte}, \textbf{acier}) à partir des \textbf{minerais de fer}, qui sont essentiellement des oxydes : l'hématite \ce{Fe2O3} et la magnétite \ce{Fe3O4}.
\end{definition}

\textbf{L'intuition d'abord.} Le fer n'existe presque jamais à l'état natif dans la nature : il est « prisonnier » de l'oxygène dans les oxydes. Obtenir du fer, c'est donc \emph{arracher l'oxygène} au minerai, autrement dit \textbf{réduire} l'oxyde de fer. Il faut pour cela un réducteur bon marché, disponible en grandes quantités et efficace à haute température : l'industrie a choisi le \textbf{carbone} (sous forme de coke) et surtout le \textbf{monoxyde de carbone} \ce{CO} qu'il permet de générer.

\courssub{Le haut fourneau : principe et équations clés}

Le \cle{haut fourneau} est un réacteur vertical de 25 à \SI{35}{m} de haut, fonctionnant en continu. On introduit par le haut (le \emph{gueulard}) un mélange de trois charges :
\begin{itemize}
\item le \textbf{minerai de fer} (\ce{Fe2O3} ou \ce{Fe3O4}) ;
\item le \textbf{coke} (carbone presque pur, obtenu par distillation de la houille) ;
\item le \textbf{calcaire} \ce{CaCO3}, appelé \emph{fondant}.
\end{itemize}
À la base, on insuffle de l'air préchauffé (vers \SI{1200}{\celsius}) par des ouvertures appelées \emph{tuyères}. Les gaz chauds montent, les charges descendent : c'est un fonctionnement \textbf{à contre-courant}, qui garantit un excellent contact entre gaz réducteurs et minerai.

\begin{popfigure}
\begin{center}
\begin{tikzpicture}[scale=0.95]
% Corps du haut fourneau
\draw[very thick, popdark, fill=popcream] (-1.6,6.8) -- (-1.6,5.6) -- (-2.3,3.4) -- (-1.5,0.4) -- (-1.1,0) -- (1.1,0) -- (1.5,0.4) -- (2.3,3.4) -- (1.6,5.6) -- (1.6,6.8) -- cycle;
% Charges en haut
\draw[->, very thick, popblue] (0,8.6) -- (0,7.0);
\node[popblue, align=center] at (0,9.15) {\textbf{Charges :} minerai \ce{Fe2O3} + coke \ce{C} + calcaire \ce{CaCO3}};
% Zones
\fill[popblueL] (-1.55,4.6) rectangle (1.55,5.5);
\node[popdark] at (0,5.05) {\small Zone de réduction : \ce{Fe2O3 + 3 CO -> 2 Fe + 3 CO2}};
\fill[poporangeL] (-2.05,2.6) rectangle (2.05,4.4);
\node[popdark] at (0,3.85) {\small Zone chaude : \ce{CO2 + C -> 2 CO}};
\node[popdark] at (0,3.1) {\small \ce{CaCO3 -> CaO + CO2} puis laitier};
\fill[poppinkL] (-1.35,0.5) rectangle (1.35,2.4);
\node[popdark] at (0,1.9) {\small Combustion : \ce{C + O2 -> CO2}};
\node[popdark] at (0,1.35) {\small ($T \approx \SI{2000}{\celsius}$)};
% Tuyères
\draw[->, very thick, poporange] (-3.6,1.0) -- (-1.5,1.0);
\node[poporange, left, align=right] at (-3.7,1.0) {\small Air chaud\\ \small (tuyères)};
\draw[->, very thick, poporange] (3.6,1.0) -- (1.5,1.0);
% Sorties basses
\draw[->, very thick, popgreen] (1.2,0.25) -- (3.4,0.25);
\node[popgreen, right, align=left] at (3.5,0.55) {\small \textbf{Laitier} (surnage)};
\draw[->, very thick, poppurple] (0.4,-0.05) -- (0.4,-1.1);
\node[poppurple, right, align=left] at (0.6,-0.85) {\small \textbf{Fonte liquide} (fer + carbone)};
% Gaz de gueulard
\draw[->, thick, popmuted] (1.6,6.4) -- (3.2,6.4);
\node[popmuted, right] at (3.2,6.4) {\small Gaz (\ce{CO}, \ce{CO2}, \ce{N2})};
\end{tikzpicture}
\end{center}
\legende{Schéma simplifié du haut fourneau : les charges descendent, les gaz réducteurs montent ; la fonte et le laitier sont récupérés à la base.}
\end{popfigure}

\textbf{Étape 1 : combustion du coke aux tuyères.} L'air insufflé à la base enflamme le coke ; cette combustion très exothermique fournit la chaleur nécessaire à tout le procédé :
\[ \ce{C + O2 -> CO2} \qquad (\text{fortement exothermique, } T \approx \SI{2000}{\celsius}) \]

\textbf{Étape 2 : formation du réducteur \ce{CO}.} En montant, le dioxyde de carbone rencontre le coke incandescent et est réduit en monoxyde de carbone (réaction de Boudouard) :
\[ \ce{CO2 + C -> 2 CO} \qquad (\text{endothermique, } T > \SI{1000}{\celsius}) \]

\textbf{Étape 3 : réduction du minerai.} Dans la partie supérieure du haut fourneau, le monoxyde de carbone réduit l'oxyde de fer. Pour l'hématite :
\[ \ce{Fe2O3 + 3 CO -> 2 Fe + 3 CO2} \qquad (\text{vers } \SI{900}{\celsius}) \]
Pour la magnétite (oxyde mixte \ce{FeO*Fe2O3}, nombre d'oxydation moyen du fer $= +8/3$) :
\[ \ce{Fe3O4 + 4 CO -> 3 Fe + 4 CO2} \]

\begin{attention}[Qui réduit réellement le minerai ?]
Dans le haut fourneau, le \textbf{réducteur direct} du minerai de fer est le \cle{monoxyde de carbone} \ce{CO}, et non le carbone directement. Le coke joue un double rôle \emph{indirect} : il fournit la chaleur (combustion) et il régénère en permanence le \ce{CO} à partir du \ce{CO2}. Écrire \ce{Fe2O3 + 3 C -> 2 Fe + 3 CO} n'est pas la description fidèle du procédé industriel.
\end{attention}

\begin{exemplebox}[Analyse redox de la réduction de l'hématite]
Vérifions avec les nombres d'oxydation que \ce{Fe2O3 + 3 CO -> 2 Fe + 3 CO2} est bien une réaction d'oxydoréduction.
\begin{itemize}
\item Dans \ce{Fe2O3} : $2 \times n.o.(\text{Fe}) + 3 \times (-2) = 0$, donc $n.o.(\text{Fe}) = +3$.
\item Dans \ce{CO} : $n.o.(\text{C}) + (-2) = 0$, donc $n.o.(\text{C}) = +2$.
\item À droite : Fe corps simple, $n.o.(\text{Fe}) = 0$ ; dans \ce{CO2} : $n.o.(\text{C}) + 2 \times (-2) = 0$, donc $n.o.(\text{C}) = +4$.
\end{itemize}
Le fer passe de $+3$ à $0$ : il \textbf{gagne} 3 électrons par atome, il est \textbf{réduit} (\ce{Fe2O3} est l'oxydant). Le carbone passe de $+2$ à $+4$ : il \textbf{perd} 2 électrons par atome, il est \textbf{oxydé} (\ce{CO} est le réducteur). Le bilan électronique est cohérent : $2 \times 3\ e^-$ gagnés par le fer $= 3 \times 2\ e^-$ cédés par le carbone.
\end{exemplebox}

\courssub{Le rôle du calcaire : l'élimination de la gangue}

Le minerai n'est jamais pur : il contient une \textbf{gangue}, constituée essentiellement de silice \ce{SiO2} et d'alumine, inutilisable et encombrante. C'est là qu'intervient le calcaire. Sous l'effet de la chaleur, il se décompose :
\[ \ce{CaCO3 -> CaO + CO2} \qquad (\text{décomposition thermique, vers } \SI{900}{\celsius}) \]
L'oxyde de calcium \ce{CaO} (chaux vive), qui est un oxyde \emph{basique}, réagit alors avec la silice, oxyde \emph{acide} :
\[ \ce{CaO + SiO2 -> CaSiO3} \]
Le silicate de calcium \ce{CaSiO3} formé constitue le \cle{laitier} : liquide à la température du haut fourneau, \textbf{moins dense que la fonte}, il surnage et peut être soutiré séparément. Le laitier n'est pas un déchet : il est valorisé en cimenterie (fabrication des ciments, comme ceux de CIMENCAM à Douala et Figuil).

\courssub{De la fonte à l'acier}

Le produit brut du haut fourneau est la \cle{fonte} : du fer contenant encore 3 à 5 \% de carbone (ainsi que du silicium, du phosphore, du soufre). Ce carbone en excès la rend dure mais \textbf{cassante}. Pour obtenir l'\cle{acier} (0,02 à 2 \% de carbone), on procède à l'\textbf{affinage} : dans un convertisseur, on insuffle du dioxygène qui oxyde le carbone excédentaire (\ce{C + O2 -> CO2}) et les impuretés. On obtient ainsi un métal à la fois résistant et ductile, matériau de base de la construction moderne.

\courssub{L'aluminothermie : une réduction éclaire}

\textbf{L'intuition.} L'aluminium est un métal très « avide » d'oxygène : son oxyde \ce{Al2O3} est extrêmement stable. Conséquence spectaculaire : l'aluminium en poudre peut \emph{arracher l'oxygène} à de nombreux oxydes métalliques moins stables, en dégageant une chaleur considérable. C'est l'\cle{aluminothermie}.

\begin{experience}[Réaction d'aluminothermie (démonstration)]
\textbf{Dispositif et protocole.} Dans un creuset réfractaire posé au-dessus d'un sable de protection, on mélange intimement de la poudre d'aluminium et de l'oxyde de fer(III) \ce{Fe2O3} (mélange appelé \emph{thermite}). On amorce la réaction avec un ruban de magnésium enflammé ou un mélange amorce (\ce{BaO2} + Al).\par
\textbf{Observations.} Après amorçage, la réaction se propage d'elle-même avec une vive incandescence, des projections et une élévation de température spectaculaire (environ \SI{3000}{\celsius}). À la fin, on recueille une goutte de \textbf{fer liquide} au fond du creuset, recouverte d'une scorie blanche d'alumine.\par
\textbf{Interprétation.} L'aluminium réduit l'oxyde de fer ; la réaction est si exothermique que le fer formé fond (température de fusion du fer : \SI{1538}{\celsius}).
\end{experience}

L'équation-bilan s'écrit :
\[ \ce{2 Al + Fe2O3 -> Al2O3 + 2 Fe} \qquad (\text{très exothermique, amorçage nécessaire}) \]

\begin{exoresolu}[Analyse redox de l'aluminothermie]
\textbf{Énoncé.} À l'aide des nombres d'oxydation, montrer que la réaction \ce{2 Al + Fe2O3 -> Al2O3 + 2 Fe} est une réaction d'oxydoréduction. Identifier l'oxydant, le réducteur, l'espèce oxydée et l'espèce réduite, puis vérifier la conservation du transfert d'électrons.
\tcblower
\textbf{Solution détaillée.}
\begin{itemize}
\item \textbf{Nombres d'oxydation initiaux.} Al est un corps simple : $n.o.(\text{Al}) = 0$. Dans \ce{Fe2O3} : $2 \times n.o.(\text{Fe}) + 3 \times (-2) = 0$ donne $n.o.(\text{Fe}) = +3$.
\item \textbf{Nombres d'oxydation finaux.} Dans \ce{Al2O3} : $2 \times n.o.(\text{Al}) + 3 \times (-2) = 0$ donne $n.o.(\text{Al}) = +3$. Fe est un corps simple : $n.o.(\text{Fe}) = 0$.
\item \textbf{Variations.} L'aluminium passe de $0$ à $+3$ : son nombre d'oxydation \emph{augmente}, il est \textbf{oxydé} ; c'est donc le \textbf{réducteur}. Le fer passe de $+3$ à $0$ : son nombre d'oxydation \emph{diminue}, il est \textbf{réduit} ; \ce{Fe2O3} est donc l'\textbf{oxydant}.
\item \textbf{Bilan électronique.} Chaque atome d'aluminium cède 3 électrons ; il y en a 2, soit $6\ e^-$ cédés. Chaque ion fer(III) gagne 3 électrons ; il y en a 2, soit $6\ e^-$ gagnés. Le transfert est équilibré : l'équation est bien cohérente.
\end{itemize}
\textbf{Conclusion.} \ce{Al} est le réducteur (oxydé en \ce{Al2O3}) et \ce{Fe2O3} est l'oxydant (réduit en \ce{Fe}).
\end{exoresolu}

\begin{popfigure}
\begin{center}
\begin{tikzpicture}[scale=0.95]
% Creuset
\draw[very thick, popdark, fill=popcream] (-1.8,3.2) -- (-1.2,0.6) -- (1.2,0.6) -- (1.8,3.2);
\draw[very thick, popdark] (-1.8,3.2) -- (1.8,3.2);
% Mélange thermite
\fill[poporangeL] (-1.55,2.2) -- (-1.2,0.65) -- (1.2,0.65) -- (1.55,2.2) -- cycle;
\node[popdark] at (0,1.45) {\small Mélange Al + \ce{Fe2O3}};
% Amorce
\draw[very thick, poppink] (0,3.3) -- (0,2.3);
\node[poppink, above] at (0,3.35) {\small Amorce (ruban de Mg enflammé)};
% Étincelles
\foreach \a in {30,70,110,150} {
  \draw[->, popgold, thick] (0,2.4) -- (\a:1.1);
}
% Goutte de fer
\draw[very thick, poppurple, fill=poppurpleL] (0,0.6) -- (0,-0.4);
\draw[poppurple, fill=poppurple] (0,-0.55) circle (0.14);
\node[poppurple, right, align=left] at (0.25,-0.55) {\small Fer en fusion ($\approx \SI{3000}{\celsius}$)};
% Rail
\draw[very thick, popdark, fill=popmuted!30] (-3.2,-1.5) rectangle (3.2,-1.1);
\draw[very thick, popdark] (-2.6,-1.1) -- (-2.6,-0.9) -- (2.6,-0.9) -- (2.6,-1.1);
\draw[poppurple, fill=poppurpleL] (-0.35,-0.9) rectangle (0.35,-1.5);
\node[popdark, below] at (0,-1.55) {\small Joint entre deux rails à souder};
\end{tikzpicture}
\end{center}
\legende{Creuset d'aluminothermie : la réaction amorcée produit du fer liquide qui coule dans le moule entourant le joint des rails.}
\end{popfigure}

\courssub{Applications de l'aluminothermie}

\begin{itemize}
\item \textbf{Soudure des rails de chemin de fer.} Le creuset est placé au-dessus d'un moule entourant le joint entre deux rails ; le fer liquide produit coule dans l'interstice et soude les rails en refroidissant. C'est la \emph{soudure aluminothermique}, encore utilisée sur les voies ferrées du monde entier.
\item \textbf{Obtention de métaux purs.} L'aluminothermie permet de préparer des métaux difficiles à réduire par le carbone, comme le chrome et le manganèse :
\[ \ce{Cr2O3 + 2 Al -> Al2O3 + 2 Cr} \qquad ; \qquad \ce{3 MnO2 + 4 Al -> 2 Al2O3 + 3 Mn} \]
\item \textbf{Autres usages.} Incendiaires et pyrotechnie (forte chaleur dégagée), soudures de fortune en milieu isolé.
\end{itemize}

\begin{savaistu}[Pourquoi le fer obtenu est-il liquide ?]
La réaction \ce{2 Al + Fe2O3 -> Al2O3 + 2 Fe} dégage environ \SI{850}{kJ} par mole de \ce{Fe2O3} : la température atteint localement \SI{3000}{\celsius}, très au-dessus de la température de fusion du fer (\SI{1538}{\celsius}). Le fer est donc obtenu \textbf{liquide} et coule par gravité : c'est exactement ce qu'il faut pour remplir le moule d'une soudure de rail. C'est la grande stabilité de l'alumine \ce{Al2O3} qui rend la réaction si exothermique.
\end{savaistu}

\begin{savaistu}[Le Cameroun, futur pays métallurgique ?]
Le Cameroun possède d'importants gisements de \textbf{minerai de fer à Mbalam} (région de l'Est, estimés à plusieurs centaines de millions de tonnes) et de \textbf{bauxite à Minim-Martap et Ngaoundal} (Adamaoua), matière première de l'aluminium — déjà transformée à Édéa par ALUCAM à partir d'alumine importée. L'exploitation de ces ressources pourrait un jour doter le pays d'une véritable filière sidérurgique et aluminique : les réactions redox que tu étudies ici sont au cœur de ces enjeux industriels nationaux.
\end{savaistu}

\courssub{Généralisation : le principe de la métallurgie}

\begin{aretenir}[Principe général de la métallurgie]
La \cle{métallurgie} consiste à extraire un métal de son minerai en \textbf{réduisant son oxyde} (ou un autre composé) à l'aide d'un \textbf{réducteur adapté} :
\[ \text{oxyde métallique} + \text{réducteur} \ce{->} \text{métal} + \text{oxyde du réducteur} \]
Le choix du réducteur dépend de la stabilité de l'oxyde à détruire et de considérations économiques :
\begin{itemize}
\item \textbf{carbone / monoxyde de carbone} : bon marché, utilisé pour le fer (haut fourneau), le zinc, l'étain ;
\item \textbf{dihydrogène} \ce{H2} : réducteur propre (il donne de l'eau), utilisé pour des métaux de haute pureté (tungstène, molybdène) ;
\item \textbf{aluminium} : réducteur très puissant (aluminothermie), pour Cr, Mn et les soudures ;
\item les métaux dont l'oxyde est \emph{très} stable (Al, Mg, Na\ldots) ne peuvent pas être obtenus par réduction chimique : on recourt à l'\textbf{électrolyse}, une oxydoréduction forcée par le courant électrique.
\end{itemize}
Dans tous les cas, l'analyse par les nombres d'oxydation reste la même : le métal voit son nombre d'oxydation \emph{diminuer} jusqu'à zéro (réduction), tandis que celui du réducteur \emph{augmente} (oxydation).
\end{aretenir}

\begin{exercice}[Pour t'entraîner]
\begin{enumerate}
\item Dans le haut fourneau, la magnétite est réduite selon \ce{Fe3O4 + 4 CO -> 3 Fe + 4 CO2}. Déterminer le nombre d'oxydation du fer dans \ce{Fe3O4}, identifier l'oxydant et le réducteur, et vérifier le bilan du transfert d'électrons.
\item Écrire et équilibrer l'équation de l'aluminothermie de l'oxyde de chrome(III) \ce{Cr2O3}. Indiquer les variations de nombres d'oxydation.
\item Pourquoi ne peut-on pas obtenir l'aluminium par réduction de \ce{Al2O3} au monoxyde de carbone ?
\end{enumerate}
\end{exercice}

\begin{corrige}[Exercice]
\textbf{1.} Dans \ce{Fe3O4} : $3 \times n.o.(\text{Fe}) + 4 \times (-2) = 0$, donc $n.o.(\text{Fe}) = +8/3$ (valeur moyenne : \ce{Fe3O4} est un oxyde mixte contenant du fer(II) et du fer(III)). Le fer passe de $+8/3$ à $0$ : il est réduit, \ce{Fe3O4} est l'oxydant. Le carbone passe de $+2$ (dans \ce{CO}) à $+4$ (dans \ce{CO2}) : il est oxydé, \ce{CO} est le réducteur. Bilan : $3 \times 8/3 = 8\ e^-$ gagnés par le fer, $4 \times 2 = 8\ e^-$ cédés par le carbone : le transfert est équilibré.\par
\textbf{2.} \ce{Cr2O3 + 2 Al -> 2 Cr + Al2O3}. Le chrome passe de $+3$ à $0$ (réduction), l'aluminium de $0$ à $+3$ (oxydation) ; $2 \times 3\ e^-$ échangés de part et d'autre.\par
\textbf{3.} L'oxyde d'aluminium \ce{Al2O3} est \emph{plus} stable que \ce{CO2} et \ce{CO} : le monoxyde de carbone n'est pas un réducteur assez puissant pour arracher l'oxygène de l'alumine aux températures accessibles. C'est précisément pourquoi l'aluminium, lui, réduit les autres oxydes (aluminothermie), et pourquoi on obtient l'aluminium par électrolyse de l'alumine fondue.
\end{corrige}

Tu maîtrises désormais le cœur de la métallurgie : réduire un oxyde, c'est faire « descendre » le nombre d'oxydation d'un métal jusqu'à zéro. Dans la section suivante, nous appliquerons exactement la même logique — mais en sens inverse, avec des \emph{oxydations} ménagées — à la préparation industrielle de deux acides majeurs : l'acide sulfurique et l'acide nitrique.

\coursec{Préparation industrielle des acides sulfurique et nitrique}

Dans la section précédente, nous avons vu comment l'oxydoréduction par voie sèche permet d'extraire les métaux de leurs minerais : le fer dans le haut fourneau, le chrome ou le manganèse par aluminothermie. Mais la réduction des oxydes métalliques n'est pas la seule grande application industrielle des réactions redox. Deux des produits chimiques les plus fabriqués au monde, l'\cle{acide sulfurique} \ce{H2SO4} et l'\cle{acide nitrique} \ce{HNO3}, sont obtenus par des \emph{chaînes d'oxydations successives} où le dioxygène de l'air joue le rôle d'oxydant. Le nombre d'oxydation va nous servir de fil conducteur : nous allons suivre, étape par étape, l'ascension du n.o. du soufre (de $0$ à $+\mathrm{VI}$) et celui de l'azote (de $-\mathrm{III}$ à $+\mathrm{V}$).

\courssub{Pourquoi ces deux acides sont-ils si importants ?}

\begin{savaistu}[Un baromètre de l'industrie]
La production mondiale d'acide sulfurique dépasse \textbf{200 millions de tonnes par an}. On dit souvent que la consommation d'acide sulfurique d'un pays est un \emph{indicateur de son développement industriel}, tant cet acide est utilisé partout : fabrication des engrais phosphatés (superphosphates), raffinage du pétrole (comme à la SONARA à Limbé), batteries au plomb des automobiles, traitement des minerais, industrie chimique. L'acide nitrique, lui, sert surtout à fabriquer les \cle{engrais azotés} (nitrate d'ammonium \ce{NH4NO3}) et les explosifs. Deux acides, deux piliers de l'agriculture et de l'industrie modernes.
\end{savaistu}

L'idée économique centrale est simple : le dioxygène de l'air est un \textbf{oxydant gratuit et abondant}. Toute la stratégie industrielle consiste à faire passer, par étapes contrôlées, un élément bon marché (le soufre $S$, l'ammoniac \ce{NH3}) d'un état d'oxydation bas à un état d'oxydation élevé, celui de l'acide désiré.

\begin{methode}[Analyser une étape industrielle en trois gestes]
Pour chaque étape d'un procédé industriel, applique toujours la même démarche :
\begin{enumerate}
\item \textbf{Écrire} l'équation-bilan équilibrée de l'étape, avec ses conditions (catalyseur, température).
\item \textbf{Calculer} les nombres d'oxydation de l'élément clé (S ou N) avant et après, ainsi que celui de l'oxygène.
\item \textbf{Conclure} : l'espèce dont le n.o. \emph{augmente} est le réducteur (elle est oxydée) ; celle dont le n.o. \emph{diminue} est l'oxydant (elle est réduite). En général, c'est \ce{O2} qui passe de $0$ à $-\mathrm{II}$.
\end{enumerate}
\end{methode}

\courssub{Préparation de l'acide sulfurique : le procédé de contact}

Le procédé de contact part du \cle{soufre} (ou de la \cle{pyrite} \ce{FeS2}, un minerai très répandu dont le grillage donne \ce{SO2}) et comporte \textbf{trois étapes}.

\textbf{Étape 1 : combustion du soufre.} On brûle le soufre fondu dans un courant d'air sec :
\[ \ce{S + O2 -> SO2} \]
Analyse redox : le soufre passe de n.o. $0$ à $+\mathrm{IV}$ (dans \ce{SO2} : $x + 2\times(-2)=0 \Rightarrow x=+4$). Il est \emph{oxydé} ; le dioxygène passe de $0$ à $-\mathrm{II}$, il est \emph{réduit}. C'est une combustion classique, totale et exothermique.

\textbf{Étape 2 : oxydation catalytique de \ce{SO2}.} C'est l'étape délicate. La réaction
\[ \ce{2 SO2 + O2 <=>[V2O5,\ \SI{450}{\celsius}] 2 SO3} \]
est \emph{lente} et \emph{réversible}. Pour la rendre rapide et rentable, on utilise un catalyseur, le \cle{pentoxyde de divanadium} \ce{V2O5}, vers \SI{450}{\celsius} (c'est le « contact » du gaz avec le catalyseur solide qui a donné son nom au procédé). Le soufre passe de $+\mathrm{IV}$ à $+\mathrm{VI}$ : nouvelle oxydation par \ce{O2}.

\textbf{Étape 3 : absorption de \ce{SO3}.} Le trioxyde de soufre réagit avec l'eau selon :
\[ \ce{SO3 + H2O -> H2SO4} \]
Ici, aucun n.o. ne change (S reste à $+\mathrm{VI}$, O à $-\mathrm{II}$, H à $+\mathrm{I}$) : ce n'est \emph{pas} une réaction redox, mais une simple hydratation, extrêmement exothermique.

\begin{attention}[Pourquoi n'absorbe-t-on pas \ce{SO3} directement dans l'eau pure ?]
La réaction $\ce{SO3 + H2O -> H2SO4}$ est si violente que, dans l'eau pure, elle produit une \textbf{brume d'acide} (fines gouttelettes de \ce{H2SO4} en suspension) qui ne se condense pas et s'échappe de l'installation : perte de produit et pollution. En pratique, on absorbe donc \ce{SO3} dans de l'\textbf{acide sulfurique concentré} (on obtient de l'\emph{oléum}, \ce{H2S2O7}), puis on \emph{dilue} prudemment ensuite. Retiens : absorption dans l'acide, dilution après.
\end{attention}

\begin{popfigure}
\begin{center}
\begin{tikzpicture}[font=\small, >=Stealth,
  etape/.style={draw=popblue, fill=popblueL, rounded corners=2pt, minimum height=8mm, inner sep=4pt},
  conv/.style={draw=poporange, fill=poporangeL, rounded corners=2pt, minimum height=8mm, inner sep=4pt}]
\node[etape] (S) {\ce{S}};
\node[etape, right=14mm of S] (SO2) {\ce{SO2}};
\node[conv, right=16mm of SO2] (C) {Convertisseur};
\node[etape, right=16mm of C] (SO3) {\ce{SO3}};
\node[conv, right=16mm of SO3] (T) {Tour d'absorption};
\node[etape, right=16mm of T] (H) {\ce{H2SO4}};
\draw[->, thick, popdark] (S) -- node[above, font=\scriptsize]{\ce{O2}, combustion} (SO2);
\draw[->, thick, popdark] (SO2) -- node[above, font=\scriptsize]{\ce{O2}} (C);
\draw[->, thick, popdark] (C) -- node[above, font=\scriptsize]{\ce{V2O5}, \SI{450}{\celsius}} (SO3);
\draw[->, thick, popdark] (SO3) -- (T);
\draw[->, thick, popdark] (T) -- node[above, font=\scriptsize]{dilution} (H);
\node[below=3mm of S, font=\scriptsize\itshape, popmuted] {n.o. $= 0$};
\node[below=3mm of SO2, font=\scriptsize\itshape, popmuted] {n.o. $= +\mathrm{IV}$};
\node[below=3mm of SO3, font=\scriptsize\itshape, popmuted] {n.o. $= +\mathrm{VI}$};
\node[below=3mm of H, font=\scriptsize\itshape, popmuted] {n.o. $= +\mathrm{VI}$};
\end{tikzpicture}
\end{center}
\legende{Chaîne de production de l'acide sulfurique par le procédé de contact : le soufre est oxydé en deux temps ($0 \to +\mathrm{IV} \to +\mathrm{VI}$), puis \ce{SO3} est absorbé dans l'acide concentré.}
\end{popfigure}

\begin{aretenir}[Procédé de contact]
Trois étapes : combustion $\ce{S -> SO2}$, oxydation catalytique $\ce{2SO2 + O2 <=> 2SO3}$ sur \ce{V2O5} à \SI{450}{\celsius}, absorption $\ce{SO3 + H2O -> H2SO4}$ (via l'acide concentré). Le soufre est le réducteur, le dioxygène de l'air est l'oxydant.
\end{aretenir}

\courssub{Préparation de l'acide nitrique : le procédé Ostwald}

Le procédé \cle{Ostwald} part de l'\cle{ammoniac} \ce{NH3} (lui-même fabriqué par synthèse Haber-Bosch à partir de \ce{N2} et \ce{H2}) et comporte également \textbf{trois étapes}.

\textbf{Étape 1 : oxydation catalytique de l'ammoniac.} On fait passer un mélange d'ammoniac et d'air sur une \emph{gaze de platine-rhodium} chauffée (environ \SI{850}{\celsius}) :
\[ \ce{4 NH3 + 5 O2 ->[Pt/Rh] 4 NO + 6 H2O} \]
Analyse redox : dans \ce{NH3}, l'azote a pour n.o. $-\mathrm{III}$ ($x + 3\times(+1) = 0 \Rightarrow x = -3$) ; dans \ce{NO}, il vaut $+\mathrm{II}$. L'azote est donc \emph{oxydé} (perte de 5 électrons par atome) ; \ce{O2} est réduit de $0$ à $-\mathrm{II}$. Vérifie l'équilibre électronique : $4$ atomes N perdent $4\times 5 = 20$ électrons, et $5$ molécules \ce{O2} gagnent $5\times 4 = 20$ électrons.

\textbf{Étape 2 : oxydation du monoxyde d'azote.} À l'air, le gaz incolore \ce{NO} s'oxyde spontanément en dioxyde d'azote \ce{NO2}, gaz roux :
\[ \ce{2 NO + O2 -> 2 NO2} \]
L'azote passe de $+\mathrm{II}$ à $+\mathrm{IV}$.

\textbf{Étape 3 : absorption dans l'eau.} Le dioxyde d'azote barbote dans l'eau :
\[ \ce{3 NO2 + H2O -> 2 HNO3 + NO} \]
Cette fois, il n'y a plus de dioxygène : c'est \ce{NO2} qui joue \emph{à la fois} le rôle d'oxydant et de réducteur. C'est une \cle{dismutation} (ou disproportionation).

\begin{exemplebox}[Montrons que l'étape 3 est une dismutation]
Dans \ce{NO2} : $x + 2\times(-2) = 0 \Rightarrow x = +\mathrm{IV}$. Dans \ce{HNO3} : $1 + x + 3\times(-2) = 0 \Rightarrow x = +\mathrm{V}$. Dans \ce{NO} : $x + (-2) = 0 \Rightarrow x = +\mathrm{II}$. Sur 3 atomes d'azote à $+\mathrm{IV}$, \textbf{deux sont oxydés} à $+\mathrm{V}$ (dans \ce{HNO3}) et \textbf{un est réduit} à $+\mathrm{II}$ (dans \ce{NO}). Un même élément, dans une même espèce, subit à la fois oxydation et réduction : c'est la définition d'une dismutation.
\end{exemplebox}

\begin{attention}[Le recyclage du NO : économie et environnement]
Le \ce{NO} formé à l'étape 3 n'est pas perdu : il est \textbf{recyclé} vers l'étape 2, où il est réoxydé en \ce{NO2}. Ce recyclage a un double intérêt : il améliore le rendement global en acide nitrique, et il évite de rejeter les oxydes d'azote \ce{NO} et \ce{NO2} dans l'atmosphère. Ces gaz, notés \ce{NO_x}, sont des polluants majeurs : ils provoquent des irritations respiratoires et contribuent aux \textbf{pluies acides} (ils forment \ce{HNO3} dans les nuages), qui acidifient les sols et les lacs et dégradent les bâtiments.
\end{attention}

\begin{popfigure}
\begin{center}
\begin{tikzpicture}[font=\small, >=Stealth,
  etape/.style={draw=poppurple, fill=poppurpleL, rounded corners=2pt, minimum height=8mm, inner sep=4pt},
  conv/.style={draw=poporange, fill=poporangeL, rounded corners=2pt, minimum height=8mm, inner sep=4pt}]
\node[etape] (NH3) {\ce{NH3}};
\node[conv, right=14mm of NH3] (G) {Gaze Pt/Rh};
\node[etape, right=14mm of G] (NO) {\ce{NO}};
\node[etape, right=14mm of NO] (NO2) {\ce{NO2}};
\node[conv, right=14mm of NO2] (A) {Absorption \ce{H2O}};
\node[etape, right=14mm of A] (H) {\ce{HNO3}};
\draw[->, thick, popdark] (NH3) -- node[above, font=\scriptsize]{\ce{O2}} (G);
\draw[->, thick, popdark] (G) -- (NO);
\draw[->, thick, popdark] (NO) -- node[above, font=\scriptsize]{\ce{O2}} (NO2);
\draw[->, thick, popdark] (NO2) -- (A);
\draw[->, thick, popdark] (A) -- (H);
\draw[->, thick, popgreen, dashed] (A.south) to[out=-90,in=-90] node[below, font=\scriptsize]{recyclage du \ce{NO}} (NO.south);
\node[above=3mm of NH3, font=\scriptsize\itshape, popmuted] {$-\mathrm{III}$};
\node[above=3mm of NO, font=\scriptsize\itshape, popmuted] {$+\mathrm{II}$};
\node[above=3mm of NO2, font=\scriptsize\itshape, popmuted] {$+\mathrm{IV}$};
\node[above=3mm of H, font=\scriptsize\itshape, popmuted] {$+\mathrm{V}$};
\end{tikzpicture}
\end{center}
\legende{Chaîne de production de l'acide nitrique par le procédé Ostwald. Le monoxyde d'azote formé lors de l'absorption est recyclé.}
\end{popfigure}

\begin{popfigure}
\begin{center}
\begin{tikzpicture}
\begin{axis}[
  width=0.85\linewidth, height=5.2cm,
  xlabel={Étape du procédé Ostwald}, ylabel={n.o. de l'azote},
  symbolic x coords={NH3, NO, NO2, HNO3},
  xticklabels={\ce{NH3}, \ce{NO}, \ce{NO2}, \ce{HNO3}},
  xtick=data, ymin=-4, ymax=6,
  grid=both, grid style={popline!40},
  ytick={-3,-2,...,5},
]
\addplot[popcurve, mark=*, mark size=2.5pt] coordinates {(NH3,-3) (NO,2) (NO2,4) (HNO3,5)};
\node[font=\scriptsize, popmuted, anchor=south] at (axis cs:NH3,-3) {$-\mathrm{III}$};
\node[font=\scriptsize, popmuted, anchor=south] at (axis cs:NO,2) {$+\mathrm{II}$};
\node[font=\scriptsize, popmuted, anchor=south] at (axis cs:NO2,4) {$+\mathrm{IV}$};
\node[font=\scriptsize, popmuted, anchor=south] at (axis cs:HNO3,5) {$+\mathrm{V}$};
\end{axis}
\end{tikzpicture}
\end{center}
\legende{Ascension du nombre d'oxydation de l'azote au cours du procédé Ostwald : de $-\mathrm{III}$ dans l'ammoniac à $+\mathrm{V}$ dans l'acide nitrique.}
\end{popfigure}

\courssub{Analyse redox globale et exercice de synthèse}

\begin{aretenir}[L'idée directrice des deux procédés]
Dans les deux procédés, un élément bon marché (S ou N) est \textbf{oxydé par étapes successives} jusqu'à l'état d'oxydation maximal ($+\mathrm{VI}$ pour S, $+\mathrm{V}$ pour N), l'oxydant étant le \textbf{dioxygène de l'air}, gratuit. Les catalyseurs (\ce{V2O5}, Pt/Rh) accélèrent les étapes lentes ; les absorptions finales transforment les oxydes en acides.
\end{aretenir}

\begin{exoresolu}[Bilan électronique de l'étape clé du procédé Ostwald]
Énoncé. On considère l'étape $\ce{4 NH3 + 5 O2 -> 4 NO + 6 H2O}$. (1) Déterminer le n.o. de l'azote dans \ce{NH3} et dans \ce{NO}. (2) Identifier l'oxydant et le réducteur. (3) Vérifier la conservation du nombre d'électrons échangés.
\tcblower
Solution. (1) Dans \ce{NH3} : $x + 3\times(+1) = 0 \Rightarrow x = -\mathrm{III}$. Dans \ce{NO} : $x + (-2) = 0 \Rightarrow x = +\mathrm{II}$. (2) Le n.o. de N augmente ($-\mathrm{III} \to +\mathrm{II}$) : \ce{NH3} est le \textbf{réducteur}, il est oxydé. Le n.o. de O diminue ($0 \to -\mathrm{II}$) : \ce{O2} est l'\textbf{oxydant}, il est réduit. (3) Chaque atome N perd $5$ électrons ; $4$ atomes N perdent $20$ électrons. Chaque molécule \ce{O2} gagne $4$ électrons ($2\times 2$) ; $5$ molécules gagnent $20$ électrons. Les électrons perdus par le réducteur sont exactement ceux gagnés par l'oxydant : l'équation est cohérente.
\end{exoresolu}

\begin{exercice}[Exercice 1]
(1) Calculer le n.o. du soufre dans \ce{SO2}, \ce{SO3} et \ce{H2SO4}. (2) Montrer que l'étape $\ce{SO3 + H2O -> H2SO4}$ n'est pas une réaction d'oxydoréduction. (3) Dans le procédé Ostwald, écrire l'équation-bilan globale de la transformation de \ce{NH3} en \ce{HNO3} en combinant les trois étapes (le \ce{NO} recyclé doit disparaître).
\end{exercice}

\begin{corrige}[Exercice 1]
(1) \ce{SO2} : $x - 4 = 0 \Rightarrow x = +\mathrm{IV}$ ; \ce{SO3} : $x - 6 = 0 \Rightarrow x = +\mathrm{VI}$ ; \ce{H2SO4} : $2 + x - 8 = 0 \Rightarrow x = +\mathrm{VI}$. (2) S reste à $+\mathrm{VI}$, O à $-\mathrm{II}$, H à $+\mathrm{I}$ : aucun n.o. ne varie, donc ce n'est pas une réaction redox. (3) On combine les étapes en tenant compte du recyclage du \ce{NO} :
\begin{itemize}
\item Étape 1 : $\ce{4 NH3 + 5 O2 -> 4 NO + 6 H2O}$
\item Étape 2 (multipliée par 3 pour fournir 6 \ce{NO2}) : $\ce{6 NO + 3 O2 -> 6 NO2}$
\item Étape 3 (multipliée par 2 pour consommer 6 \ce{NO2}) : $\ce{6 NO2 + 2 H2O -> 4 HNO3 + 2 NO}$
\end{itemize}
En sommant : $\ce{4 NH3 + 5 O2 + 6 NO + 3 O2 + 6 NO2 + 2 H2O -> 4 NO + 6 H2O + 6 NO2 + 4 HNO3 + 2 NO}$.
On simplifie les espèces intermédiaires (\ce{6 NO2} s'annulent ; \ce{6 NO} à gauche et \ce{4 NO + 2 NO} à droite s'annulent ; \ce{2 H2O} à gauche et \ce{6 H2O} à droite donnent \ce{4 H2O} à droite).
Bilan global : $\ce{4 NH3 + 8 O2 -> 4 HNO3 + 4 H2O}$, soit en divisant par 4 :
\[ \ce{NH3 + 2 O2 -> HNO3 + H2O} \]
Vérification des n.o. : N passe de $-\mathrm{III}$ à $+\mathrm{V}$ (perte de $8\,e^-$ par atome N), et $2\,\ce{O2}$ gagnent $8\,e^-$ ($4 \times 2$) : le bilan électronique est équilibré.
\end{corrige}

\begin{attention}[Erreurs fréquentes]
\begin{itemize}
\item Ne confonds pas \emph{oxydation} et \emph{acidification} : l'étape d'absorption ($\ce{SO3 + H2O}$) n'est pas redox, car aucun n.o. ne change.
\item Dans $\ce{3 NO2 + H2O -> 2 HNO3 + NO}$, n'oublie pas que seuls \textbf{deux} atomes d'azote sur trois finissent dans l'acide ; le troisième est réduit en \ce{NO}, d'où l'importance du recyclage pour le rendement.
\item Le catalyseur (\ce{V2O5} ou Pt/Rh) ne figure \emph{pas} dans l'équation-bilan : il accélère la réaction sans être consommé.
\item Pour le bilan global Ostwald, n'écris pas une équation d'étape 3 modifiée arbitrairement (ex. $\ce{4 NO2 + 2 H2O -> 4 HNO3 + 2 NO}$) qui ne serait pas équilibrée. Respecte la stœchiométrie $\ce{3 NO2 -> 2 HNO3 + NO}$ et gère le recyclage par combinaison algébrique.
\end{itemize}
\end{attention}

Ainsi, grâce aux nombres d'oxydation, tu peux désormais lire une filière industrielle entière comme une histoire d'électrons : le soufre et l'azote « montent » patiemment l'échelle des états d'oxydation, poussés par le dioxygène de l'air, jusqu'à devenir les deux acides qui nourrissent l'agriculture et font tourner l'industrie mondiale.

\newpage

\coursec{Travaux pratiques}

\coursec{Leçon 10 : Oxydoréduction par voie sèche et nombres d'oxydation}

\courssub{Objectifs du TP}

À l'issue de cette séance de travaux pratiques, tu dois être capable de :

\begin{itemize}
\item réaliser en toute sécurité la combustion du magnésium dans l'air et celle du carbone dans le dioxygène ;
\item identifier expérimentalement les produits formés : l'oxyde de magnésium \ce{MgO} (caractère basique) et le dioxyde de carbone \ce{CO2} (test à l'eau de chaux) ;
\item écrire les équations-bilan de ces combustions par voie sèche ;
\item calculer le nombre d'oxydation de chaque élément avant et après la réaction ;
\item montrer, à l'aide des nombres d'oxydation, que ces combustions sont des réactions d'\cle{oxydoréduction} et identifier l'oxydant, le réducteur, l'oxydation et la réduction.
\end{itemize}

\courssub{Matériel et produits}

\begin{center}
\rowcolors{2}{popblueL}{white}
\begin{tabularx}{\linewidth}{|l|X|}
\hline
\rowcolor{popblue!40} \textbf{Matériel} & \textbf{Produits} \\
\hline
1 bécher de \SI{250}{mL} & Ruban de magnésium (environ \SI{10}{cm}) \\
\hline
1 pince en fer (ou pince à creuset) & Morceau de charbon de bois (environ \SI{2}{g}) \\
\hline
1 flacon à gaz avec couvercle en verre & Eau oxygénée à 10 volumes (environ \SI{20}{mL}) \\
\hline
1 cuillère à combustion (ou fil de fer recourbé) & Dioxyde de manganèse \ce{MnO2} en poudre (une pointe de spatule) \\
\hline
1 pipette Pasteur, 1 agitateur en verre & Eau de chaux (solution saturée de \ce{Ca(OH)2}) \\
\hline
Papier pH (ou quelques gouttes de BBT) & Eau distillée \\
\hline
Bec Bunsen (ou lampe à alcool) & Lunettes de protection, gants, blouse \\
\hline
\end{tabularx}
\end{center}

\begin{savaistu}[Et si un produit manque ?]
Dans beaucoup de lycées camerounais, on peut s'organiser avec les ressources locales : le \cle{charbon de bois} de cuisine convient parfaitement comme source de carbone ; le ruban de magnésium peut être remplacé par de la laine de fer fine (la combustion est moins éclatante mais observable) ; si l'eau oxygénée manque, le dioxygène peut être récupéré par déplacement d'eau à partir de la décomposition de \ce{KMnO4} chauffé ; à défaut de papier pH, le jus de \cle{bissap} (infusion de fleurs d'hibiscus) est un excellent indicateur coloré naturel : il vire au vert-jaune en milieu basique.
\end{savaistu}

\courssub{Sécurité}

\begin{attention}[Consignes de sécurité impératives]
\begin{itemize}
\item \textbf{Pictogramme « comburant »} (flamme sur un cercle) : le dioxygène ne brûle pas mais il \emph{active violemment} les combustions. Tenir les flacons de \ce{O2} loin de toute flamme et de toute graisse.
\item \textbf{Pictogramme « corrosif / irritant »} : l'eau oxygénée et l'eau de chaux sont irritantes pour la peau et les yeux. Porter \textbf{blouse, gants et lunettes} pendant toute la séance.
\item \textbf{Danger pour les yeux} : la combustion du magnésium émet une lumière blanche très intense, riche en ultraviolets. \textbf{Ne jamais regarder directement la flamme} : observer de biais ou à travers une vitre.
\item Tenir le ruban de magnésium avec une \textbf{pince}, jamais à la main ; effectuer la combustion \textbf{au-dessus d'un bécher} pour recueillir la poudre et éviter tout contact avec la paillasse.
\item Le flacon contenant le charbon incandescent devient chaud : le manipuler avec des gants ou un chiffon sec.
\item En cas de projection dans les yeux : rincer abondamment à l'eau pendant plusieurs minutes et prévenir immédiatement l'enseignant.
\end{itemize}
\end{attention}

\courssub{Schéma du dispositif}

\begin{popfigure}
\begin{center}
\begin{tikzpicture}[scale=0.95, every node/.style={font=\small}]
% --- Dispositif 1 : combustion du magnésium ---
\begin{scope}
% bécher
\draw[thick, popdark] (-1.2,0) -- (-1.2,1.8) (-1.2,0) -- (1.2,0) (1.2,0) -- (1.2,1.8);
\draw[thick, popdark] (-1.35,1.8) -- (-1.2,1.8) (1.2,1.8) -- (1.35,1.8);
% poudre blanche
\fill[popgold!60] (-0.9,0.02) .. controls (-0.3,0.35) and (0.3,0.35) .. (0.9,0.02) -- cycle;
\node[popdark] at (0,-0.45) {poudre blanche de \ce{MgO}};
% pince + ruban
\draw[thick, popmuted] (0.15,3.6) -- (0.05,2.6);
\draw[thick, popmuted] (0.35,3.6) -- (0.25,2.6);
\draw[very thick, popblue] (0.15,2.6) -- (0.15,2.0);
\node[popblue, right] at (0.4,2.3) {ruban de Mg};
% flamme
\fill[poporange] (0.15,1.95) .. controls (-0.15,1.6) and (0.0,1.35) .. (0.15,1.15) .. controls (0.35,1.4) and (0.45,1.65) .. (0.15,1.95);
\fill[popgold] (0.15,1.7) .. controls (0.02,1.5) and (0.08,1.35) .. (0.15,1.25) .. controls (0.25,1.4) and (0.28,1.55) .. (0.15,1.7);
% étoiles de lumière
\foreach \a in {30,90,150,210,270,330}{
  \draw[popgold, thick] ({0.15+0.55*cos(\a)},{1.6+0.45*sin(\a)}) -- ({0.15+0.8*cos(\a)},{1.6+0.6*sin(\a)});}
\node[popdark, align=center] at (0,4.1) {\textbf{Expérience 1 : combustion du magnésium}};
\end{scope}
% --- Dispositif 2 : combustion du carbone ---
\begin{scope}[xshift=8.2cm]
% flacon à gaz
\draw[thick, popdark] (-1.1,0) -- (-1.1,2.6) (-1.1,0) -- (1.1,0) (1.1,0) -- (1.1,2.6);
\draw[thick, popdark] (-1.1,2.6) -- (-0.45,2.6) (0.45,2.6) -- (1.1,2.6);
\draw[thick, popdark] (-0.45,2.6) -- (-0.45,3.0) (0.45,2.6) -- (0.45,3.0);
\draw[thick, popdark] (-0.45,3.0) -- (0.45,3.0);
% gaz O2
\node[poppurple] at (0,2.1) {\ce{O2}};
% cuillère à combustion + charbon incandescent
\draw[thick, popmuted] (0.0,4.2) -- (0.0,1.35);
\draw[thick, popmuted] (-0.35,1.35) arc (180:360:0.35);
\fill[poporange!80] (-0.18,1.32) circle (0.16);
\fill[popgold] (-0.18,1.32) circle (0.07);
\node[poporange, right] at (0.15,1.0) {charbon incandescent};
% eau de chaux
\draw[popblue!50, fill=popblue!15] (-1.05,0.02) -- (-1.05,0.55) -- (1.05,0.55) -- (1.05,0.02) -- cycle;
\node[popblue, below] at (0,0.02) {eau de chaux};
\node[popdark, align=center] at (0,4.6) {\textbf{Expérience 2 : combustion du carbone}};
\end{scope}
\end{tikzpicture}
\end{center}
\legende{Dispositifs expérimentaux : à gauche, combustion du ruban de magnésium au-dessus d'un bécher ; à droite, introduction du charbon incandescent dans un flacon de dioxygène, puis ajout d'eau de chaux.}
\end{popfigure}

\courssub{Protocole}

\textbf{Préparation du dioxygène (à faire en premier).}

\begin{enumerate}
\item Verser environ \SI{20}{mL} d'eau oxygénée à 10 volumes dans un erlenmeyer propre.
\item Ajouter une pointe de spatule de dioxyde de manganèse \ce{MnO2} : un dégagement gazeux immédiat se produit selon $\ce{2H2O2 ->[MnO2] 2H2O + O2}$.
\item Recueillir le gaz dans le flacon à gaz par déplacement d'air (le dioxygène est plus dense que l'air) pendant environ 2 minutes, puis fermer avec le couvercle en verre. Vérifier la présence de \ce{O2} avec une bûchette en point incandescent : elle se \emph{ravive}.
\end{enumerate}

\textbf{Séance 1 : combustion du magnésium.}

\begin{enumerate}
\setcounter{enumi}{3}
\item Découper un ruban de magnésium d'environ \SI{10}{cm} et le nettoyer rapidement avec du papier abrasif pour éliminer la couche d'oxyde superficielle.
\item Saisir le ruban avec la pince, placer le bécher en dessous, puis enflammer l'extrémité libre à la flamme du bec Bunsen. \textbf{Observer de biais}, sans regarder directement la lumière.
\item Laisser la poudre blanche formée tomber dans le bécher. Ajouter environ \SI{50}{mL} d'eau distillée, agiter pour mettre la poudre en suspension.
\item Tremper une baguette de verre dans la suspension et déposer une goutte sur du papier pH (ou ajouter 2 gouttes de BBT). Noter la couleur et la valeur du pH.
\end{enumerate}

\textbf{Séance 2 : combustion du carbone.}

\begin{enumerate}
\setcounter{enumi}{7}
\item Placer un morceau de charbon de bois (environ \SI{2}{g}) dans la cuillère à combustion et le chauffer fortement à la flamme jusqu'à ce qu'il soit \emph{incandescent} (rougeoyant).
\item Introduire rapidement la cuillère dans le flacon de dioxygène et refermer le couvercle. Observer la combustion.
\item Une fois la combustion terminée et le flacon tiédi, verser environ \SI{20}{mL} d'eau de chaux, refermer et \textbf{agiter vigoureusement}. Observer.
\end{enumerate}

\courssub{Observations et résultats attendus}

\begin{center}
\rowcolors{2}{popgreenL}{white}
\begin{tabularx}{\linewidth}{|X|X|X|}
\hline
\rowcolor{popgreen!45} \textbf{Étape} & \textbf{Observation} & \textbf{Interprétation} \\
\hline
Combustion du Mg & Lumière blanche éblouissante, dégagement de chaleur ; dépôt d'une poudre blanche & Formation d'oxyde de magnésium \ce{MgO} ; réaction très exothermique \\
\hline
Suspension de \ce{MgO} + papier pH & Papier pH : bleu-vert, pH $\approx 9$ à $10$ (BBT : bleu) & \ce{MgO} est un oxyde basique : $\ce{MgO + H2O -> Mg(OH)2}$ \\
\hline
Charbon dans \ce{O2} & Le charbon brûle vivement avec des étincelles, bien plus que dans l'air & Le dioxygène pur active la combustion \\
\hline
Ajout d'eau de chaux & L'eau de chaux se \textbf{trouble} (précipité blanc laiteux) & Présence de \ce{CO2} : $\ce{CO2 + Ca(OH)2 -> CaCO3 + H2O}$ \\
\hline
\end{tabularx}
\end{center}

\courssub{Exploitation}

\begin{exoresolu}[Analyse des deux combustions par les nombres d'oxydation]
\textbf{Questions.}
\begin{enumerate}
\item Écrire les équations-bilan des combustions du magnésium et du carbone dans le dioxygène.
\item Attribuer le nombre d'oxydation (n.o.) de chaque élément avant et après chaque réaction.
\item Montrer que ces réactions sont des réactions d'oxydoréduction ; identifier l'oxydant, le réducteur, l'oxydation et la réduction.
\end{enumerate}
\tcblower
\textbf{1. Équations-bilan.}
\[ \ce{2Mg + O2 -> 2MgO} \qquad\qquad \ce{C + O2 -> CO2} \]

\textbf{2. Nombres d'oxydation.} Rappel des règles : n.o. d'un corps simple $= 0$ ; n.o. de l'oxygène dans un oxyde $= -II$.

\begin{center}
\begin{tabularx}{\linewidth}{|l|X|X|}
\hline
\rowcolor{popblue!40} \textbf{Élément} & \textbf{Avant la réaction} & \textbf{Après la réaction} \\
\hline
Magnésium & \ce{Mg} : n.o. $= 0$ & \ce{Mg} dans \ce{MgO} : n.o. $= +II$ \\
\hline
Oxygène & \ce{O2} : n.o. $= 0$ & \ce{O} dans \ce{MgO} : n.o. $= -II$ \\
\hline
Carbone & \ce{C} : n.o. $= 0$ & \ce{C} dans \ce{CO2} : n.o. $= +IV$ \\
\hline
Oxygène & \ce{O2} : n.o. $= 0$ & \ce{O} dans \ce{CO2} : n.o. $= -II$ \\
\hline
\end{tabularx}
\end{center}

\textbf{3. Caractère redox.} Dans les deux réactions, les nombres d'oxydation de certains éléments \textbf{varient} : ce sont bien des réactions d'\cle{oxydoréduction} (par voie sèche, car elles ont lieu sans solution aqueuse).

\begin{itemize}
\item Le n.o. du magnésium passe de $0$ à $+II$ (celui du carbone de $0$ à $+IV$) : il \emph{augmente}, donc \ce{Mg} (resp. \ce{C}) \textbf{s'oxyde} : c'est le \textbf{réducteur}. Il subit une \textbf{oxydation} (perte d'électrons : $\ce{Mg -> Mg^2+ + 2e-}$).
\item Le n.o. de l'oxygène passe de $0$ à $-II$ : il \emph{diminue}, donc \ce{O2} \textbf{se réduit} : c'est l'\textbf{oxydant}. Il subit une \textbf{réduction} (gain d'électrons : $\ce{O2 + 4e- -> 2O^2-}$).
\end{itemize}

Vérification par les demi-équations électroniques pour le magnésium : $\ce{2Mg -> 2Mg^2+ + 4e-}$ et $\ce{O2 + 4e- -> 2O^2-}$ ; la somme redonne bien $\ce{2Mg + O2 -> 2MgO}$.
\end{exoresolu}

\begin{exercice}[Pour aller plus loin]
Le fer brûle également dans le dioxygène en donnant l'oxyde de fer(II,III) (magnétite) \ce{Fe3O4}. Équilibrer l'équation-bilan, déterminer le n.o. du fer dans \ce{Fe3O4} (on montrera qu'il est fractionnaire, égal à $+8/3$, car \ce{Fe3O4} contient deux ions \ce{Fe^3+} et un ion \ce{Fe^2+}), puis identifier l'oxydant et le réducteur.
\end{exercice}

\begin{corrige}[Exercice « Pour aller plus loin »]
Équation équilibrée : $\ce{3Fe + 2O2 -> Fe3O4}$. Dans \ce{Fe3O4}, la somme des n.o. est nulle : $3x + 4\times(-II) = 0$ d'où $x = +8/3 \approx +2,67$. Ce n.o. moyen fractionnaire s'explique : \ce{Fe3O4} est un oxyde mixte \ce{FeO.Fe2O3}, contenant un fer au n.o. $+II$ et deux fers au n.o. $+III$. Le n.o. du fer augmente (de $0$ à $+8/3$) : le fer est le \textbf{réducteur}, il s'oxyde ; le n.o. de l'oxygène diminue (de $0$ à $-II$) : \ce{O2} est l'\textbf{oxydant}, il se réduit.
\end{corrige}

\courssub{Conclusion}

Ce TP a permis d'observer deux combustions par \cle{voie sèche} : celle du magnésium, éblouissante, qui fournit un oxyde basique \ce{MgO} (pH $\approx 10$ en suspension), et celle du carbone, qui produit \ce{CO2}, mis en évidence par le trouble de l'eau de chaux. Dans les deux cas, le calcul des nombres d'oxydation montre une \textbf{augmentation} du n.o. du réducteur (\ce{Mg} : $0 \to +II$ ; \ce{C} : $0 \to +IV$) et une \textbf{diminution} de celui de l'oxydant (\ce{O} : $0 \to -II$). Une combustion dans le dioxygène est donc toujours une réaction d'\cle{oxydoréduction}, même en l'absence de toute solution : le transfert d'électrons s'effectue directement entre les réactifs, à l'état solide ou gazeux. Ce principe est à la base de nombreux procédés industriels, comme la réduction des minerais de fer en sidérurgie ou l'aluminothermie, étudiés dans la suite de la leçon.

\newpage

\coursec{Méthodes et exercices résolus}

\coursec{Leçon 10 : Oxydoréduction par voie sèche et nombres d'oxydation}

\courssub{1. Les réactions d'oxydoréduction par voie sèche}

\begin{definition}[Réaction d'oxydoréduction par voie sèche]
Une \cle{réaction d'oxydoréduction par voie sèche} est une réaction chimique mettant en jeu un transfert d'électrons entre réactifs, qui se déroule \emph{sans solvant aqueux}, généralement à haute température. Il s'agit le plus souvent de \cle{combustions} d'éléments (métaux ou non-métaux) dans le dioxygène $\ce{O2}$ ou le dichlore $\ce{Cl2}$.
\end{definition}

\begin{exemplebox}[Principales combustions par voie sèche au programme]
\begin{itemize}
\item \textbf{Magnésium dans le dioxygène} : $\ce{2Mg + O2 -> 2MgO}$ (flamme blanche éblouissante, poudre blanche).
\item \textbf{Dihydrogène dans le dioxygène} : $\ce{2H2 + O2 -> 2H2O}$ (flamme pâle, formation d'eau).
\item \textbf{Dihydrogène dans le dichlore} : $\ce{H2 + Cl2 -> 2HCl}$ (réaction explosive à la lumière, formation de chlorure d'hydrogène gazeux).
\item \textbf{Carbone (graphite/coke) dans le dioxygène} : $\ce{C + O2 -> CO2}$ (combustion complète) ; $\ce{2C + O2 -> 2CO}$ (combustion incomplète, fourneau).
\end{itemize}
Dans tous ces cas, un corps simple (n.o. $=0$) se transforme en composé (n.o. $\neq 0$) : c'est la signature d'une réaction redox.
\end{exemplebox}

\begin{methode}[Exploiter quantitativement une réaction de combustion par voie sèche]
\begin{enumerate}
\item \textbf{Écrire et équilibrer} l'équation-bilan de la combustion (vérifier la conservation des atomes et des charges).
\item \textbf{Calculer les quantités de matière initiales} :
      $n = \dfrac{m}{M}$ pour un solide (carbone, magnésium...),
      $n = \dfrac{V}{V_m}$ pour un gaz (dioxygène, dichlore...) aux conditions données.
\item \textbf{Construire un tableau d'avancement} et déterminer l'avancement maximal $x_{\max}$ (réactif limitant = celui qui donne le plus petit $x_{\max}$).
\item \textbf{Établir le bilan final} : quantités de matière et masses des produits formés, réactif en excès restant.
\end{enumerate}
\textbf{Réflexe} : le produit d'une combustion par voie sèche est un solide (oxyde, chlorure) ou un gaz (dioxyde de carbone, chlorure d'hydrogène) ; la stœchiométrie 1/1 ou 2/1 est fréquente, vérifie-la avant tout calcul.
\end{methode}

\begin{exoresolu}[Combustion du carbone dans le dioxygène]
On brûle complètement $m = \SI{6,0}{g}$ de carbone (graphite) dans un flacon contenant $V = \SI{4,0}{L}$ de dioxygène, mesuré dans les conditions où $V_m = \SI{24,0}{L.mol^{-1}}$. Donnée : $M(\ce{C}) = \SI{12,0}{g.mol^{-1}}$.
\begin{enumerate}
\item Écrire l'équation-bilan de la combustion complète du carbone.
\item Déterminer le réactif limitant.
\item Calculer la masse de dioxyde de carbone formé et la masse de carbone non réagi.
\end{enumerate}
\tcblower
\textbf{1. Équation-bilan}
\[ \ce{C + O2 -> CO2} \]
Vérification redox : le carbone passe de $0$ (corps simple) à $+IV$ (dans $\ce{CO2}$) : \cle{oxydation} ; l'oxygène passe de $0$ à $-II$ : \cle{réduction}.

\textbf{2. Réactif limitant}
Quantités initiales :
\[ n_0(\ce{C}) = \frac{6,0}{12,0} = \SI{0,50}{mol} \qquad ; \qquad n_0(\ce{O2}) = \frac{4,0}{24,0} = \SI{0,167}{mol} \]

\begin{center}
\begin{tabularx}{\linewidth}{|l|X|X|X|X|}
\hline
\rowcolor{popblueL} État & Avancement & \ce{C} & \ce{O2} & \ce{CO2} \\
\hline
Initial & $0$ & $0,50$ & $0,167$ & $0$ \\
\hline
En cours & $x$ & $0,50 - x$ & $0,167 - x$ & $x$ \\
\hline
Final & $x_{\max}$ & $0,50 - x_{\max}$ & $0,167 - x_{\max}$ & $x_{\max}$ \\
\hline
\end{tabularx}
\end{center}

Si $\ce{C}$ limitant : $x_{\max} = 0,50\ \text{mol}$.
Si $\ce{O2}$ limitant : $x_{\max} = 0,167\ \text{mol}$.
La plus petite valeur l'emporte : $x_{\max} = \SI{0,167}{mol}$. Le \cle{dioxygène est le réactif limitant}.

\textbf{3. Bilan final}
$n(\ce{CO2}) = x_{\max} = \SI{0,167}{mol}$.
$M(\ce{CO2}) = 12,0 + 2\times 16,0 = \SI{44,0}{g.mol^{-1}}$.
\[ m(\ce{CO2}) = 0,167 \times 44,0 \approx \SI{7,3}{g} \]
Carbone restant : $n(\ce{C}) = 0,50 - 0,167 = \SI{0,333}{mol}$ soit $m = 0,333 \times 12,0 \approx \SI{4,0}{g}$.

\textbf{Conclusion} : le dioxygène est limitant ; il se forme $\SI{7,3}{g}$ de $\ce{CO2}$ et il reste $\SI{4,0}{g}$ de carbone non brûlé.
\end{exoresolu}

\begin{popfigure}
\begin{tikzpicture}[scale=0.9, every node/.style={font=\small}]
% Blast furnace schematic
\draw[popdark, thick] (0,0) -- (0,6) -- (1,6.5) -- (3,6.5) -- (4,6) -- (4,0) -- cycle;
\draw[popdark, thick] (1,6.5) -- (1,7.5) -- (3,7.5) -- (3,6.5);
% Zones
\fill[poporangeL] (0,0) rectangle (4,1.5) node[pos=0.5, align=center, text=popdark] {Foyer\\Combustion coke\\$\ce{C + O2 -> CO2}$\\(T $\approx \SI{2000}{\celsius}$)};
\fill[poppinkL] (0,1.5) rectangle (4,3.5) node[pos=0.5, align=center, text=popdark] {Bosse\\Réduction directe\\$\ce{Fe2O3 + 3C -> 2Fe + 3CO}$};
\fill[popgoldL] (0,3.5) rectangle (4,5.5) node[pos=0.5, align=center, text=popdark] {Ventre\\Réduction indirecte\\$\ce{Fe2O3 + 3CO -> 2Fe + 3CO2}$\\(T $\approx \SI{800}{\celsius}$)};
\fill[popgreenL] (0,5.5) rectangle (4,6) node[pos=0.5, align=center, text=popdark] {Chargement\\Minerai + Coke + Fondant};
% Arrows
\draw[-Stealth, popblue, thick] (4.2,0.7) -- (5,0.7) node[right] {Fonte liquide};
\draw[-Stealth, popblue, thick] (4.2,1.2) -- (5,1.2) node[right] {Laitier};
\draw[-Stealth, popblue, thick] (2,7.5) -- (2,8) node[above] {Gaz $\ce{CO2, N2}$};
\draw[-Stealth, poporange, thick] (-0.5,3) -- (0,3) node[left] {Air chaud};
\draw[-Stealth, popdark, thick] (-0.5,5.8) -- (0,5.8) node[left] {Charge};
\end{tikzpicture}
\legende{Schéma simplifié d'un haut fourneau (sidérurgie) : zones de température et réactions principales.}
\end{popfigure}

\courssub{2. Le nombre d'oxydation : définition et règles de calcul}

\begin{definition}[Nombre d'oxydation]
Le \cle{nombre d'oxydation} (n.o.) d'un atome dans une espèce chimique est la charge électrique que cet atome \emph{porterait} si l'on considérait toutes les liaisons comme \emph{ioniques} (les électrons de liaison attribués à l'atome le plus électronégatif). Il s'exprime en \cle{chiffres romains} précédés du signe ($+I$, $-II$, $+VI$, $0$, etc.).
\end{definition}

\begin{aretenir}[Règles conventionnelles de calcul du nombre d'oxydation (ordre de priorité)]
\begin{enumerate}
\item \textbf{Corps simple} : n.o. $= 0$ (ex. $\ce{Fe}$, $\ce{O2}$, $\ce{Cl2}$, $\ce{S8}$).
\item \textbf{Ion monoatomique} : n.o. $=$ charge de l'ion (ex. $\ce{Fe^3+} \to +III$, $\ce{Cl-} \to -I$).
\item \textbf{Valeurs de référence (quasi systématiques)} :
      \begin{itemize}
      \item Hydrogène : $+I$ (sauf dans les \emph{hydrures} ioniques type $\ce{NaH}$ où il est $-I$).
      \item Oxygène : $-II$ (sauf dans les \emph{peroxoïdes} $\ce{H2O2}$, $\ce{Na2O2}$ où il est $-I$, et dans $\ce{OF2}$ où il est $+II$).
      \item Fluor : toujours $-I$ (plus électronégatif).
      \item Métaux alcalins (Groupe 1) : $+I$ ; Alcalino-terreux (Groupe 2) : $+II$ ; Aluminium : $+III$.
      \end{itemize}
\item \textbf{Règle de la somme} : La somme algébrique des n.o. de \emph{tous} les atomes de l'espèce est égale à sa \cle{charge globale} ($0$ pour une molécule neutre, $q$ pour un ion de charge $q$).
\end{enumerate}
\end{aretenir}

\begin{methode}[Déterminer le nombre d'oxydation d'un élément dans une espèce]
\begin{enumerate}
\item Identifier la nature de l'espèce (corps simple, ion, composé).
\item Attribuer les n.o. connus (règles 1 à 3 ci-dessus).
\item Poser $x$ = n.o. de l'élément cherché.
\item Appliquer la règle de la somme : $\sum (\text{n.o.} \times \text{nombre d'atomes}) = \text{charge de l'espèce}$.
\item Résoudre l'équation et exprimer le résultat en \textbf{chiffres romains} avec signe ($+VI$, $-II$, $0$...).
\item \textbf{Vérifier} : refaire la somme pour retomber sur la charge de l'espèce.
\end{enumerate}
\end{methode}

\begin{exoresolu}[Détermination de nombres d'oxydation]
Déterminer le nombre d'oxydation : a) du soufre dans l'acide sulfurique $\ce{H2SO4}$ ; b) de l'azote dans l'ion nitrate $\ce{NO3-}$ ; c) du manganèse dans l'ion permanganate $\ce{MnO4-}$ ; d) du carbone dans le dioxyde de carbone $\ce{CO2}$ ; e) du chrome dans l'ion dichromate $\ce{Cr2O7^2-}$.
\tcblower
On applique la règle de la somme avec $\mathrm{n.o.}(\ce{O}) = -II$ et $\mathrm{n.o.}(\ce{H}) = +I$.

\textbf{a) Soufre dans $\ce{H2SO4}$} (molécule neutre, somme $= 0$) :
\[ 2\times(+I) + x + 4\times(-II) = 0 \;\Rightarrow\; 2 + x - 8 = 0 \;\Rightarrow\; x = +VI \]
$\mathrm{n.o.}(\ce{S}) = +VI$.

\textbf{b) Azote dans $\ce{NO3-}$} (ion de charge $-1$, somme $= -1$) :
\[ x + 3\times(-II) = -1 \;\Rightarrow\; x - 6 = -1 \;\Rightarrow\; x = +V \]
$\mathrm{n.o.}(\ce{N}) = +V$.

\textbf{c) Manganèse dans $\ce{MnO4-}$} (somme $= -1$) :
\[ x + 4\times(-II) = -1 \;\Rightarrow\; x = +VII \]
$\mathrm{n.o.}(\ce{Mn}) = +VII$.

\textbf{d) Carbone dans $\ce{CO2}$} (somme $= 0$) :
\[ x + 2\times(-II) = 0 \;\Rightarrow\; x = +IV \]
$\mathrm{n.o.}(\ce{C}) = +IV$.

\textbf{e) Chrome dans $\ce{Cr2O7^2-}$} (somme $= -2$) :
\[ 2x + 7\times(-II) = -2 \;\Rightarrow\; 2x - 14 = -2 \;\Rightarrow\; 2x = 12 \;\Rightarrow\; x = +VI \]
$\mathrm{n.o.}(\ce{Cr}) = +VI$ (chaque atome de Cr).

\textbf{Conclusion :} les degrés d'oxydation maximaux pour \ce{S}, \ce{N}, \ce{Mn}, \ce{Cr} sont atteints dans ces oxydes/oxoanions.
\end{exoresolu}

\courssub{3. Nombres d'oxydation et réactions redox}

\begin{propriete}[Critère des nombres d'oxydation]
Une réaction chimique est une \cle{réaction d'oxydoréduction} si et seulement si le nombre d'oxydation d'au moins un élément \textbf{change} au cours de la réaction.
\begin{itemize}
\item Si le n.o. \textbf{augmente} : l'élément est \cle{oxydé} (perte d'électrons) $\rightarrow$ c'est le \cle{réducteur}.
\item Si le n.o. \textbf{diminue} : l'élément est \cle{réduit} (gain d'électrons) $\rightarrow$ c'est l'\cle{oxydant}.
\end{itemize}
La présence d'un \textbf{corps simple} parmi les réactifs ou les produits est un indice fort (n.o. passe de $0$ à $\neq 0$).
\end{propriete}

\begin{methode}[Montrer qu'une réaction est une réaction d'oxydoréduction]
\begin{enumerate}
\item \textbf{Attribuer} le n.o. de \emph{chaque} élément dans \emph{chaque} espèce (réactifs et produits) grâce aux règles de la leçon 2.
\item \textbf{Comparer} les n.o. d'un même élément avant (réactifs) et après (produits).
\item \textbf{Conclure} :
      \begin{itemize}
      \item Variation $\uparrow$ (ex $0 \to +II$) $\rightarrow$ \textbf{Oxydation} (réducteur).
      \item Variation $\downarrow$ (ex $0 \to -II$) $\rightarrow$ \textbf{Réduction} (oxydant).
      \item Au moins une variation $\rightarrow$ \textbf{Réaction redox}.
      \item Aucune variation $\rightarrow$ \textbf{Pas de redox} (réaction acido-basique, précipitation, complexe...).
      \end{itemize}
\end{enumerate}
\end{methode}

\begin{exoresolu}[Combustion du magnésium dans le dioxygène]
Équation : $\ce{2Mg + O2 -> 2MgO}$. Montrer que c'est une réaction redox. Identifier oxydant et réducteur.
\tcblower
\textbf{Attribution des n.o.} :
\begin{itemize}
\item Réactifs : $\ce{Mg}$ et $\ce{O2}$ corps simples $\Rightarrow$ $\mathrm{n.o.}(\ce{Mg})=0$, $\mathrm{n.o.}(\ce{O})=0$.
\item Produit $\ce{MgO}$ : $\mathrm{n.o.}(\ce{O})=-II$ (règle O), molécule neutre $\Rightarrow$ $\mathrm{n.o.}(\ce{Mg})=+II$.
\end{itemize}

\textbf{Comparaison} :
\begin{center}
\begin{tabularx}{\linewidth}{|l|X|X|X|}
\hline
\rowcolor{popblueL} Élément & n.o. avant & n.o. après & Variation \\
\hline
Magnésium & $0$ (dans $\ce{Mg}$) & $+II$ (dans $\ce{MgO}$) & \textbf{Augmente} : Oxydation \\
\hline
Oxygène & $0$ (dans $\ce{O2}$) & $-II$ (dans $\ce{MgO}$) & \textbf{Diminue} : Réduction \\
\hline
\end{tabularx}
\end{center}

\textbf{Conclusion} : Les n.o. varient $\rightarrow$ \cle{réaction d'oxydoréduction}.
$\ce{Mg}$ est oxydé ($0 \to +II$, perte de $2\ \ce{e-}$) $\rightarrow$ c'est le \cle{réducteur}.
$\ce{O2}$ est réduit ($0 \to -II$, gain de $4\ \ce{e-}$ par molécule) $\rightarrow$ c'est l'\cle{oxydant}.
Demi-équations : $\ce{Mg -> Mg^2+ + 2e-}$ ; $\ce{O2 + 4e- -> 2O^2-}$.
\end{exoresolu}

\courssub{4. Équilibrer une équation redox grâce aux nombres d'oxydation}

\begin{methode}[Méthode des nombres d'oxydation pour équilibrer une équation redox]
\begin{enumerate}
\item \textbf{Calculer} le n.o. des éléments dont le degré d'oxydation varie (à gauche et à droite).
\item \textbf{Déterminer} la variation unitaire : $\Delta_{\text{ox}}$ (électrons perdus par atome oxydé) et $\Delta_{\text{red}}$ (électrons gagnés par atome réduit).
\item \textbf{Égaliser les échanges d'électrons} : multiplier les espèces oxydée et réduite par des coefficients tels que le total d'électrons cédés = total d'électrons captés (PPCM des variations).
\item \textbf{Équilibrer les autres éléments} (conservation de la matière). En milieu aqueux : ajouter $\ce{H2O}$ pour l'O, $\ce{H+}$ (ou $\ce{H3O+}$) pour l'H et la charge.
\item \textbf{Vérifier} : conservation de \emph{chaque} élément et conservation de la \emph{charge électrique}.
\end{enumerate}
\textbf{Réflexe} : La somme des électrons perdus (oxydation) doit \emph{toujours} égaler la somme des électrons gagnés (réduction). C'est la condition \emph{sine qua non} de l'équilibrage redox.
\end{methode}

\begin{exoresolu}[Équilibrage voie sèche : Fer + Dichlore]
Le fer brûle dans le dichlore donnant du chlorure de fer(III). Équilibrer : $\ce{Fe + Cl2 -> FeCl3}$.
\tcblower
\textbf{1. N.o. des éléments variables}
Réactifs (corps simples) : $\mathrm{n.o.}(\ce{Fe})=0$, $\mathrm{n.o.}(\ce{Cl})=0$.
Produit $\ce{FeCl3}$ : $\mathrm{n.o.}(\ce{Cl})=-I \Rightarrow \mathrm{n.o.}(\ce{Fe})=+III$.

\textbf{2. Variations}
$\ce{Fe}$ : $0 \to +III$ $\Rightarrow$ Oxydation, perte de \textbf{3 e-} par atome Fe.
$\ce{Cl}$ : $0 \to -I$ $\Rightarrow$ Réduction, gain de \textbf{1 e-} par atome Cl, soit \textbf{2 e-} par molécule $\ce{Cl2}$.

\textbf{3. Égalisation (PPCM de 3 et 2 = 6 e-)}
$2\ \ce{Fe}$ perdent $2 \times 3 = 6\ \ce{e-}$.
$3\ \ce{Cl2}$ gagnent $3 \times 2 = 6\ \ce{e-}$.

\textbf{4. Équilibrage atomes}
$2\ \ce{Fe}$ donnent $2\ \ce{FeCl3}$ (6 atomes Cl), cohérent avec $3\ \ce{Cl2}$.

\textbf{5. Équation finale}
\[ \ce{2Fe + 3Cl2 -> 2FeCl3} \]
Vérification : Fe : 2/2 ; Cl : 6/6. Charges : 0/0. \checkmark
\end{exoresolu}

\begin{exoresolu}[Équilibrage en milieu aqueux acide : Permanganate + Fer(II)]
En solution acide, l'ion permanganate $\ce{MnO4-}$ oxyde l'ion ferreux $\ce{Fe^2+}$ en ion ferrique $\ce{Fe^3+}$, se réduisant lui-même en $\ce{Mn^2+}$. Équilibrer l'équation.
\tcblower
\textbf{1. N.o. variables}
$\ce{Mn}$ dans $\ce{MnO4-}$ : $x - 8 = -1 \Rightarrow x = +VII$.
$\ce{Mn}$ dans $\ce{Mn^2+}$ : $+II$. Variation : $+VII \to +II$ = gain de \textbf{5 e-} (Réduction).
$\ce{Fe}$ : $+II \to +III$ = perte de \textbf{1 e-} (Oxydation).

\textbf{2. Égalisation électrons (PPCM 5)}
$1\ \ce{MnO4-}$ gagne $5\ \ce{e-}$.
$5\ \ce{Fe^2+}$ perdent $5 \times 1 = 5\ \ce{e-}$.
Coefficients provisoires : $1\ \ce{MnO4-} + 5\ \ce{Fe^2+} \to 1\ \ce{Mn^2+} + 5\ \ce{Fe^3+}$.

\textbf{3. Équilibrage O et H (milieu acide)}
Oxygène : $4$ à gauche $\Rightarrow$ ajouter $4\ \ce{H2O}$ à droite.
Hydrogène : $8$ H à droite $\Rightarrow$ ajouter $8\ \ce{H+}$ (ou $\ce{H3O+}$) à gauche.

\textbf{4. Vérification charges}
Gauche : $-1 + 5\times(+2) + 8\times(+1) = -1 + 10 + 8 = +17$.
Droite : $+2 + 5\times(+3) = +17$. \checkmark

\textbf{Équation finale} :
\[ \ce{MnO4- + 5Fe^2+ + 8H+ -> Mn^2+ + 5Fe^3+ + 4H2O} \]
\end{exoresolu}

\courssub{5. Applications industrielles : sidérurgie et aluminothermie}

\begin{methode}[Analyser un procédé métallurgique d'oxydoréduction]
\begin{enumerate}
\item \textbf{Identifier} le minerai (oxyde métallique = oxydant) et le réducteur (charbon/CO ou Al).
\item \textbf{Écrire} l'équation-bilan équilibrée de la réaction principale de réduction.
\item \textbf{Justifier} par les n.o. : le métal voit son n.o. \emph{diminuer} jusqu'à $0$ (réduction), le réducteur voit son n.o. \emph{augmenter} (oxydation).
\item \textbf{Préciser} les conditions (température, catalyseur) et l'intérêt (fer liquide, soudage \emph{in situ}...).
\end{enumerate}
\end{methode}

\begin{exoresolu}[Sidérurgie et aluminothermie]
\textbf{1. Haut fourneau (Sidérurgie).} Réduction de l'hématite $\ce{Fe2O3}$ par le monoxyde de carbone $\ce{CO}$ :
$\ce{Fe2O3 + 3CO -> 2Fe + 3CO2}$.
Montrer que le fer est réduit et identifier le réducteur.

\textbf{2. Aluminothermie.} Réaction : $\ce{Fe2O3 + 2Al -> 2Fe + Al2O3}$ (très exothermique). Indiquer l'élément oxydé et justifier le nom.
\tcblower
\textbf{1. Sidérurgie}
N.o. :
$\ce{Fe2O3}$ : $\mathrm{n.o.}(\ce{Fe}) = +III$ ($2x - 6 = 0$).
$\ce{Fe}$ (métal) : $\mathrm{n.o.} = 0$.
$\ce{CO}$ : $\mathrm{n.o.}(\ce{C}) = +II$. $\ce{CO2}$ : $\mathrm{n.o.}(\ce{C}) = +IV$.

Fer : $+III \to 0$ \textbf{(diminution)} $\rightarrow$ \cle{Réduction} (gain $3\ \ce{e-}$/atome). C'est l'\textbf{oxydant}.
Carbone : $+II \to +IV$ \textbf{(augmentation)} $\rightarrow$ \cle{Oxydation} (perte $2\ \ce{e-}$/atome). $\ce{CO}$ est le \cle{réducteur}.
Principe : Le coke brûle en bas ($\ce{C + O2 -> CO2}$), puis $\ce{CO2 + C -> 2CO}$ (réaction de Boudouard). Le $\ce{CO}$ monte et réduit le minerai par paliers ($\ce{Fe2O3 -> Fe3O4 -> FeO -> Fe}$).

\textbf{2. Aluminothermie}
$\ce{Al}$ : n.o. $= 0$. $\ce{Al2O3}$ : $2x - 6 = 0 \Rightarrow x = +III$.
Aluminium : $0 \to +III$ \textbf{(augmentation)} $\rightarrow$ \cle{Oxydation}. L'aluminium est le \cle{réducteur}.
Fer : $+III \to 0$ \textbf{(réduction)}.
Le nom \emph{aluminothermie} vient de \textbf{Aluminium} (réducteur) + \textbf{Thermie} (chaleur). $\Delta H \ll 0$, T $> \SI{1500}{\celsius}$ $\rightarrow$ fer \emph{liquide}$\rightarrow$ soudage rails \emph{in situ}.
\end{exoresolu}

\begin{savaistu}[Le fer au Cameroun]
Le Cameroun possède d'importants gisements de fer (Mbalam, Ngaoundal). Bien que l'extraction ne soit pas encore industrielle à grande échelle, la maîtrise de la réduction du fer par le carbone (sidérurgie) ou l'aluminium (aluminothermie) est au cœur des projets de développement de la sidérurgie nationale (usine de Kribi, partenariats avec Sundance Resources, Sinosteel). L'aluminothermie est aussi utilisée localement pour la maintenance des voies ferrées (Camrail).
\end{savaistu}

\courssub{6. Préparation industrielle des acides sulfurique et nitrique}

\begin{methode}[Rédiger les équations-bilan d'une filière d'oxydation progressive]
\begin{enumerate}
\item \textbf{Lister} les étapes successives : oxydation de l'élément clé (S ou N) par paliers jusqu'à son degré maximal ($+VI$ pour S, $+V$ pour N), puis hydratation de l'anhydride.
\item \textbf{Équilibrer} chaque étape (méthode des n.o. pour les étapes redox).
\item \textbf{Préciser} le n.o. de l'élément clé à chaque étape, le catalyseur, la température.
\item \textbf{Rappeler} que l'hydratation finale ($\ce{SO3 + H2O}$, $\ce{NO2 + H2O + O2}$) n'est \emph{pas} une réaction redox (n.o. constant).
\end{enumerate}
\end{methode}

\begin{exoresolu}[Filières industrielles : Acide sulfurique (Contact) et Acide nitrique (Ostwald)]
\textbf{1. Acide sulfurique (Procédé Contact).} Écrire les 3 étapes à partir du soufre, avec n.o.(S).
\textbf{2. Acide nitrique (Procédé Ostwald).} Écrire les étapes à partir de l'ammoniac $\ce{NH3}$, avec n.o.(N).
\tcblower
\textbf{1. Acide sulfurique ($\ce{H2SO4}$) — Chaîne : $\ce{S} \to \ce{SO2} \to \ce{SO3} \to \ce{H2SO4}$}

\textbf{Étape 1 : Combustion du soufre} (four, $\approx \SI{1000}{\celsius}$)
\[ \ce{S + O2 -> SO2} \qquad \mathrm{n.o.}(\ce{S}) : 0 \to +IV \quad \text{(Oxydation)} \]

\textbf{Étape 2 : Oxydation catalytique du $\ce{SO2}$} (convertisseur, $\ce{V2O5}$, $\SI{450}{\celsius}$)
\[ \ce{2SO2 + O2 <=> 2SO3} \qquad \mathrm{n.o.}(\ce{S}) : +IV \to +VI \quad \text{(Oxydation)} \]
\textit{Équilibrage n.o.} : S perd $2\ \ce{e-}$/atome ($2 \times 2 = 4\ \ce{e-}$ pour $2\ \ce{SO2}$). $\ce{O2}$ gagne $4\ \ce{e-}$/molécule. Ratio $2:1$. Équilibre déplacé vers l'avant par excès $\ce{O2}$ (air) et T modérée.

\textbf{Étape 3 : Hydratation / Absorption} (tour à garnissage, acide $\SI{98}{\%}$ puis dilution)
\[ \ce{SO3 + H2O -> H2SO4} \qquad \mathrm{n.o.}(\ce{S}) = +VI \text{ (inchangé)} \]
\textit{Note industrielle} : $\ce{SO3}$ absorbé dans $\ce{H2SO4}$ concentré $\to$ oléum $\ce{H2S2O7}$, puis dilution contrôlée (exothermique, évite le brouillard d'acide).

\textbf{2. Acide nitrique ($\ce{HNO3}$) — Procédé Ostwald. Chaîne : $\ce{NH3} \to \ce{NO} \to \ce{NO2} \to \ce{HNO3}$}

\textbf{Étape 1 : Oxydation catalytique de l'ammoniac} (grilles Pt-Rh, $\SI{850}{\celsius}$)
\[ \ce{4NH3 + 5O2 -> 4NO + 6H2O} \qquad \mathrm{n.o.}(\ce{N}) : -III \to +II \quad \text{(Oxydation)} \]
$\ce{NH3}$ : $x + 3(+I) = 0 \Rightarrow x = -III$. $\ce{NO}$ : $x - 2 = 0 \Rightarrow x = +II$.
Perte $5\ \ce{e-}$/N. $4\ \ce{N}$ perdent $20\ \ce{e-}$. $5\ \ce{O2}$ gagnent $20\ \ce{e-}$ ($4\ \ce{e-}$/$\ce{O2}$).

\textbf{Étape 2 : Oxydation du monoxyde d'azote} (refroidissement $\to$ oxydation air)
\[ \ce{2NO + O2 -> 2NO2} \qquad \mathrm{n.o.}(\ce{N}) : +II \to +IV \quad \text{(Oxydation)} \]

\textbf{Étape 3 : Absorption du dioxyde d'azote} (tour d'absorption, eau + air)
\[ \ce{4NO2 + O2 + 2H2O -> 4HNO3} \qquad \mathrm{n.o.}(\ce{N}) : +IV \to +V \quad \text{(Oxydation)} \]
L'oxygène de l'air assure l'oxydation finale $+IV \to +V$.

\textbf{Conclusion} : Dans les deux filières, l'élément clé (S ou N) est \cle{oxydé par paliers successifs} jusqu'à son degré d'oxydation maximal, puis l'anhydride acide final est hydraté. Ces acides sont la base de l'industrie des engrais (superphosphates, nitrates d'ammonium), vitale pour l'agriculture camerounaise (cacao, coton, banane).
\end{exoresolu}

\begin{popfigure}
\begin{tikzpicture}[scale=0.85, every node/.style={font=\small, align=center}]
% Contact Process Flowchart
\node[draw, popblue, thick, rounded corners, fill=popblueL] (S) at (0,0) {Soufre $\ce{S}$\\n.o. 0};
\node[draw, poporange, thick, rounded corners, fill=poporangeL] (air1) at (3,1.5) {Air ($\ce{O2}$)};
\node[draw, poppurple, thick, rounded corners, fill=poppurpleL] (four) at (3,0) {Four à soufre\\$\approx \SI{1000}{\celsius}$};
\node[draw, popgreen, thick, rounded corners, fill=popgreenL] (SO2) at (6,0) {$\ce{SO2}$\\n.o. +IV};
\node[draw, poppurple, thick, rounded corners, fill=poppurpleL] (conv) at (9,0) {Convertisseur\\$\ce{V2O5}$, $\SI{450}{\celsius}$\\$\ce{2SO2 + O2 <=> 2SO3}$};
\node[draw, poporange, thick, rounded corners, fill=poporangeL] (air2) at (9,1.5) {Air ($\ce{O2}$)};
\node[draw, popgreen, thick, rounded corners, fill=popgreenL] (SO3) at (12,0) {$\ce{SO3}$\\n.o. +VI};
\node[draw, poppink, thick, rounded corners, fill=poppinkL] (abs) at (15,0) {Tour d'absorption\\$\ce{H2SO4}$ $\SI{98}{\%}$\\$\ce{SO3 + H2O -> H2SO4}$};
\node[draw, popblue, thick, rounded corners, fill=popblueL] (H2SO4) at (18,0) {Acide sulfurique\\$\ce{H2SO4}$};

\draw[-Stealth, thick, popdark] (S) -- (four);
\draw[-Stealth, thick, popdark] (air1) -- (four);
\draw[-Stealth, thick, popdark] (four) -- (SO2);
\draw[-Stealth, thick, popdark] (SO2) -- (conv);
\draw[-Stealth, thick, popdark] (air2) -- (conv);
\draw[-Stealth, thick, popdark] (conv) -- (SO3);
\draw[-Stealth, thick, popdark] (SO3) -- (abs);
\draw[-Stealth, thick, popdark] (abs) -- (H2SO4);
\end{tikzpicture}
\legende{Filière de l'acide sulfurique (Procédé Contact) : oxydation progressive du soufre $0 \to +IV \to +VI$.}
\end{popfigure}

\begin{popfigure}
\begin{tikzpicture}[scale=0.85, every node/.style={font=\small, align=center}]
% Ostwald Process Flowchart
\node[draw, popblue, thick, rounded corners, fill=popblueL] (NH3) at (0,0) {Ammoniac $\ce{NH3}$\\n.o. -III};
\node[draw, poporange, thick, rounded corners, fill=poporangeL] (air1) at (3,1.5) {Air ($\ce{O2}$)};
\node[draw, poppurple, thick, rounded corners, fill=poppurpleL] (ox1) at (3,0) {Oxydation catalytique\\Pt-Rh, $\SI{850}{\celsius}$\\$\ce{4NH3 + 5O2 -> 4NO + 6H2O}$};
\node[draw, popgreen, thick, rounded corners, fill=popgreenL] (NO) at (6,0) {$\ce{NO}$\\n.o. +II};
\node[draw, poporange, thick, rounded corners, fill=poporangeL] (air2) at (9,1.5) {Air ($\ce{O2}$)};
\node[draw, poppurple, thick, rounded corners, fill=poppurpleL] (ox2) at (9,0) {Oxydation\\$\ce{2NO + O2 -> 2NO2}$};
\node[draw, popgreen, thick, rounded corners, fill=popgreenL] (NO2) at (12,0) {$\ce{NO2}$\\n.o. +IV};
\node[draw, poppink, thick, rounded corners, fill=poppinkL] (abs) at (15,0) {Tour d'absorption\\Eau + Air ($\ce{O2}$)\\$\ce{4NO2 + O2 + 2H2O -> 4HNO3}$};
\node[draw, popblue, thick, rounded corners, fill=popblueL] (HNO3) at (18,0) {Acide nitrique\\$\ce{HNO3}$};

\draw[-Stealth, thick, popdark] (NH3) -- (ox1);
\draw[-Stealth, thick, popdark] (air1) -- (ox1);
\draw[-Stealth, thick, popdark] (ox1) -- (NO);
\draw[-Stealth, thick, popdark] (NO) -- (ox2);
\draw[-Stealth, thick, popdark] (air2) -- (ox2);
\draw[-Stealth, thick, popdark] (ox2) -- (NO2);
\draw[-Stealth, thick, popdark] (NO2) -- (abs);
\draw[-Stealth, thick, popdark] (abs) -- (HNO3);
\end{tikzpicture}
\legende{Filière de l'acide nitrique (Procédé Ostwald) : oxydation progressive de l'azote $-III \to +II \to +IV \to +V$.}
\end{popfigure}

\begin{attention}[Les 5 erreurs qui coûtent des points au Probatoire]
\begin{enumerate}
\item \textbf{Confondre n.o. et charge de l'ion} : Le n.o. s'écrit en \emph{chiffres romains} avec le signe ($+VI$, $-II$), la charge en chiffres arabes avec le signe après ($2-$, $3+$). Écrire « n.o. du soufre $= +6$ » est toléré, mais « le soufre porte $6+$ » ou « charge $+VI$ » est \textbf{faux}.
\item \textbf{Oublier les exceptions majeures} : Dans $\ce{H2O2}$ (peroxoïde), $\mathrm{n.o.}(\ce{O}) = -I$ (pas $-II$) ; dans $\ce{NaH}$ (hydrure), $\mathrm{n.o.}(\ce{H}) = -I$ (pas $+I$). Appliquer aveuglément « O $= -II$ » ou « H $= +I$ » fausse tout le calcul.
\item \textbf{Mal poser la règle de la somme pour un ion} : La somme des n.o. = \textbf{charge de l'ion}, pas zéro. Ex. $\ce{MnO4-}$ : $x + 4(-II) = -1 \Rightarrow x = +VII$. Écrire $x - 8 = 0$ donne $+VIII$ (aberration).
\item \textbf{Inverser oxydation et réduction} : \textbf{Augmentation} du n.o. = \textbf{Oxydation} (perte $\ce{e-}$) = \textbf{Réducteur}. \textbf{Diminution} du n.o. = \textbf{Réduction} (gain $\ce{e-}$) = \textbf{Oxydant}. Mnémotechnique : \emph{« Augmente = Oxydation = Réducteur »} (AOR).
\item \textbf{Équilibrer une équation redox sans vérifier l'égalité des électrons} : Les coefficients stœchiométriques des couples redox doivent \textbf{d'abord} assurer : (coeff oxydé $\times$ e- perdus) = (coeff réduit $\times$ e- gagnés). Terminer \emph{toujours} par la double vérification : atomes \textbf{et} charges.
\end{enumerate}
\end{attention}

\newpage

\coursec{Exercices}

\courssub{Je vérifie mes connaissances}

\begin{exercice}[QCM : réactions d'oxydoréduction]
Pour chaque affirmation, choisir la (ou les) bonne(s) réponse(s).
\begin{enumerate}
\item Une réaction d'oxydoréduction est une réaction au cours de laquelle :
\begin{enumerate}
\item il y a toujours du dioxygène ;
\item le nombre d'oxydation d'au moins deux éléments varie ;
\item il y a transfert d'électrons ;
\item il y a transfert de protons \ce{H+}.
\end{enumerate}
\item Le nombre d'oxydation de l'oxygène vaut $-2$ :
\begin{enumerate}
\item dans tous les composés ;
\item sauf dans les peroxydes et dans \ce{OF2} ;
\item sauf dans \ce{O2} où il vaut $0$ ;
\item jamais.
\end{enumerate}
\item Lors de la combustion du magnésium dans le dioxygène $\ce{2 Mg + O2 -> 2 MgO}$ :
\begin{enumerate}
\item le magnésium est réduit ;
\item le magnésium est l'oxydant ;
\item le dioxygène est l'oxydant ;
\item le n.o. du magnésium passe de $0$ à $+2$.
\end{enumerate}
\item Une dismutation est une réaction dans laquelle :
\begin{enumerate}
\item deux éléments s'oxydent simultanément ;
\item un même élément est à la fois oxydé et réduit ;
\item l'oxydant et le réducteur sont la même espèce ;
\item le nombre d'oxydation ne varie pas.
\end{enumerate}
\end{enumerate}
\end{exercice}

\begin{exercice}[Vrai ou faux, à justifier]
Répondre par \emph{vrai} ou \emph{faux}, en justifiant chaque réponse à l'aide des nombres d'oxydation.
\begin{enumerate}
\item Dans l'ion sulfate \ce{SO4^2-}, le nombre d'oxydation du soufre est $+4$.
\item La réaction $\ce{CaCO3 -> CaO + CO2}$ est une réaction d'oxydoréduction.
\item Dans \ce{H2}, le nombre d'oxydation de l'hydrogène est $+1$.
\item Lors de la combustion du carbone $\ce{C + O2 -> CO2}$, le carbone est oxydé.
\item Le nombre d'oxydation d'un ion monoatomique est égal à la charge de cet ion.
\end{enumerate}
\end{exercice}

\begin{exercice}[Texte à trous]
Compléter le texte suivant avec les mots ou expressions de la liste : \emph{oxydation ; réduction ; oxydant ; réducteur ; augmente ; diminue ; voie sèche ; électrons}.

Une réaction d'oxydoréduction par \ldots\ldots\ldots\ldots se déroule en l'absence d'eau, souvent à haute température. Au cours d'une \ldots\ldots\ldots\ldots, le nombre d'oxydation d'un élément \ldots\ldots\ldots\ldots ; au cours d'une \ldots\ldots\ldots\ldots, il \ldots\ldots\ldots\ldots. L'espèce qui capte des \ldots\ldots\ldots\ldots est l'\ldots\ldots\ldots\ldots ; celle qui les cède est le \ldots\ldots\ldots\ldots. Dans la combustion du dihydrogène dans le dichlore $\ce{H2 + Cl2 -> 2 HCl}$, le dihydrogène joue le rôle de \ldots\ldots\ldots\ldots.
\end{exercice}

\begin{exercice}[Définitions et règles]
\begin{enumerate}
\item Définir le \cle{nombre d'oxydation} d'un élément dans un composé.
\item Énoncer la règle concernant la somme des nombres d'oxydation des éléments dans : (a) un composé neutre ; (b) un ion polyatomique.
\item Donner le nombre d'oxydation : d'un élément à l'état de corps simple ; de l'hydrogène dans la plupart des composés ; de l'oxygène dans la plupart des composés.
\item Citer deux applications industrielles de l'oxydoréduction par voie sèche étudiées dans cette leçon.
\end{enumerate}
\end{exercice}

\courssub{Je m'entraîne}

\begin{methode}[Équilibrer une équation redox par la méthode des nombres d'oxydation]
\begin{enumerate}
\item Attribuer le nombre d'oxydation (n.o.) à chaque élément dans les réactifs et les produits.
\item Identifier l'élément qui s'oxyde (n.o. augmente) et celui qui se réduit (n.o. diminue).
\item Calculer la variation de n.o. par atome pour l'oxydation ($\Delta n.o._{ox} > 0$) et pour la réduction ($\Delta n.o._{red} < 0$).
\item Trouver les coefficients stœchiométriques des espèces concernées pour que le total des électrons perdus soit égal au total des électrons gagnés (produit en croix des variations absolues).
\item Équilibrer les autres éléments (souvent O par \ce{H2O}, H par \ce{H+} en milieu acide) et vérifier la balance des charges.
\end{enumerate}
\end{methode}

\begin{exercice}[Calcul de nombres d'oxydation]
Déterminer le nombre d'oxydation de l'élément souligné dans chacune des espèces suivantes (l'élément à considérer est indiqué entre parenthèses) :
\begin{enumerate}
\item \ce{H2SO4} (soufre) ;
\item \ce{HNO3} (azote) ;
\item \ce{KMnO4} (manganèse) ;
\item \ce{Cr2O7^2-} (chrome) ;
\item \ce{NH4+} (azote) ;
\item \ce{SO3^2-} (soufre).
\end{enumerate}
\end{exercice}

\begin{exercice}[Combustion du fer dans le dioxygène]
On considère la combustion du fer dans le dioxygène, qui donne l'oxyde magnétique (magnétite) :
\[ \ce{3 Fe + 2 O2 -> Fe3O4} \]
\begin{enumerate}
\item Déterminer le nombre d'oxydation du fer dans \ce{Fe} et dans \ce{Fe3O4} (on pourra écrire la magnétite comme un oxyde mixte \ce{FeO.Fe2O3}).
\item Déterminer le nombre d'oxydation de l'oxygène avant et après la réaction.
\item Montrer qu'il s'agit d'une réaction d'oxydoréduction.
\item Identifier l'oxydant, le réducteur, l'élément oxydé et l'élément réduit.
\end{enumerate}
\end{exercice}

\begin{exercice}[Équilibrer par les nombres d'oxydation]
Équilibrer les équations suivantes par la méthode des nombres d'oxydation, en précisant pour chacune les variations de n.o. :
\begin{enumerate}
\item $\ce{Fe2O3 + CO -> Fe + CO2}$ (réduction en haut fourneau) ;
\item $\ce{Al + Cr2O3 -> Al2O3 + Cr}$ ;
\item $\ce{Cu + O2 -> CuO}$ ;
\item $\ce{SO2 + O2 -> SO3}$.
\end{enumerate}
\end{exercice}

\begin{exercice}[Le haut fourneau et la magnétite]
Un haut fourneau traite de la magnétite \ce{Fe3O4}, réduite par le monoxyde de carbone \ce{CO} en fer, avec dégagement de dioxyde de carbone.
\begin{enumerate}
\item Écrire et équilibrer l'équation-bilan de la réaction.
\item Préciser les variations des nombres d'oxydation du fer et du carbone ; identifier l'oxydant et le réducteur.
\item Calculer la masse de fer obtenue à partir d'une tonne ($m = \SI{1000}{kg}$) de magnétite pure.
\end{enumerate}
\emph{Données :} $M(\ce{Fe}) = \SI{55,8}{g.mol^{-1}}$ ; $M(\ce{O}) = \SI{16,0}{g.mol^{-1}}$.
\end{exercice}

\courssub{Je me prépare au Probatoire}

\begin{savaistu}[L'aluminothermie : une réaction spectaculaire]
La réaction d'aluminothermie ($\ce{2 Al + Fe2O3 -> Al2O3 + 2 Fe}$) libère une enthalpie considérable ($\Delta H^\circ \approx \SI{-852}{kJ.mol^{-1}}$ pour 2 mol Al), portant la température au-delà de \SI{2500}{\celsius}. Le fer fondu (p.f. \SI{1538}{\celsius}) et l'alumine fondue se séparent par différence de densité. Ce procédé, breveté par Goldschmidt (1895), sert encore aujourd'hui pour le soudage \emph{in situ} des rails de chemin de fer (CAMRAIL au Cameroun) et la réparation de pièces moulées massives. Il illustre la puissance des réactions redox par voie sèche où l'aluminium, fort réducteur, « vole » l'oxygène à l'oxyde de fer.
\end{savaistu}

\begin{exercice}[L'aluminothermie : soudure des rails]
Pour souder les rails de chemin de fer, on utilise la réaction d'aluminothermie : un mélange d'aluminium en poudre et d'oxyde de fer(III) \ce{Fe2O3}, une fois amorcé, réagit vivement en dégageant une chaleur considérable ; le fer obtenu, liquide, coule dans le moule et assure la soudure.
\begin{enumerate}
\item Écrire et équilibrer l'équation-bilan de la réaction entre l'aluminium et l'oxyde de fer(III).
\item Déterminer les nombres d'oxydation de l'aluminium et du fer avant et après la réaction. En déduire le nombre d'électrons échangés par atome d'aluminium et par atome de fer.
\item Identifier l'oxydant, le réducteur, l'espèce oxydée et l'espèce réduite.
\item On veut obtenir une masse $m_{\ce{Fe}} = \SI{1,00}{kg}$ de fer pour réaliser une soudure.
\begin{enumerate}
\item Calculer la quantité de matière de fer correspondante.
\item En déduire la masse d'aluminium nécessaire, en supposant la réaction totale.
\item Calculer la masse d'oxyde de fer(III) qu'il faut engager.
\end{enumerate}
\item En réalité, le rendement de l'opération est de $90\,\%$. Quelle masse d'aluminium faut-il prévoir en pratique ?
\end{enumerate}
\emph{Données :} $M(\ce{Al}) = \SI{27,0}{g.mol^{-1}}$ ; $M(\ce{Fe}) = \SI{55,8}{g.mol^{-1}}$ ; $M(\ce{O}) = \SI{16,0}{g.mol^{-1}}$.
\end{exercice}

\begin{savaistu}[Le procédé Ostwald et l'acide nitrique]
L'acide nitrique est à la base de la chimie des nitrates (engrais \ce{NH4NO3}, explosifs). Le procédé Ostwald (1902) utilise un catalyseur platine-rhodium à \SI{900}{\celsius} pour l'étape 1. L'étape 3 (absorption) se fait dans des tours à garnissage où l'eau ruisselle à contre-courant des gaz. Le \ce{NO} formé est recyclé à l'étape 2, ce qui améliore le rendement global. Au Cameroun, l'acide nitrique est importé pour la fabrication d'engrais complexes (NPK) par la société SPC (Société des Produits Chimiques) à Douala.
\end{savaistu}

\begin{exercice}[Le procédé Ostwald : préparation de l'acide nitrique]
L'acide nitrique \ce{HNO3} est préparé industriellement à partir de l'ammoniac \ce{NH3} selon le procédé Ostwald, qui comporte trois étapes :
\begin{itemize}
\item étape 1 : oxydation catalytique de l'ammoniac par le dioxygène en monoxyde d'azote \ce{NO} et eau ;
\item étape 2 : oxydation du monoxyde d'azote en dioxyde d'azote \ce{NO2} ;
\item étape 3 : action de l'eau sur le dioxyde d'azote, donnant l'acide nitrique et du monoxyde d'azote (recyclé).
\end{itemize}
\begin{enumerate}
\item Écrire et équilibrer les équations-bilan des trois étapes.
\item Déterminer le nombre d'oxydation de l'azote dans \ce{NH3}, \ce{NO}, \ce{NO2} et \ce{HNO3}.
\item Montrer que les étapes 1 et 2 sont des oxydations de l'azote.
\item Montrer que l'étape 3 est une \cle{dismutation} : préciser les proportions d'azote oxydé et d'azote réduit.
\item Une usine traite un volume $V = \SI{1000}{m^3}$ d'ammoniac (mesuré dans les conditions où $V_m = \SI{24,0}{L.mol^{-1}}$). En supposant le rendement global de $100\,\%$ et le recyclage total du \ce{NO} de l'étape 3, montrer que l'équation globale peut s'écrire $\ce{NH3 + 2 O2 -> HNO3 + H2O}$, puis calculer la masse d'acide nitrique obtenue.
\end{enumerate}
\emph{Données :} $M(\ce{H}) = \SI{1,0}{g.mol^{-1}}$ ; $M(\ce{N}) = \SI{14,0}{g.mol^{-1}}$ ; $M(\ce{O}) = \SI{16,0}{g.mol^{-1}}$.
\end{exercice}

\begin{savaistu}[Du minerai à l'acide : l'intégration industrielle]
Le couplage « haut fourneau / acide sulfurique » est classique dans les bassins sidérurgiques. La pyrite (\ce{FeS2}), souvent gangue des minerais de fer, est valorisée : son grillage fournit le \ce{SO2} pour l'acide sulfurique (procédé contact) et le résidu \ce{Fe2O3} alimente le haut fourneau. L'acide sulfurique produit sert au décapage des tôles (sidérurgie) et à la fabrication d'engrais (superphosphates). Au Cameroun, la SONARA (raffinage) et les cimenteries (CIMENCAM) sont de gros consommateurs d'acide sulfurique pour le traitement des eaux et la catalyse.
\end{savaistu}

\begin{exercice}[Du minerai de fer à l'acide sulfurique : le complexe industriel]
Une région minière dispose de pyrite de fer \ce{FeS2} et d'hématite \ce{Fe2O3}. On envisage : (A) la production de fonte en haut fourneau ; (B) la fabrication d'acide sulfurique par grillage de la pyrite.

\textbf{Partie A — Réduction de l'hématite.}
\begin{enumerate}
\item Écrire et équilibrer l'équation de la réduction de \ce{Fe2O3} par le monoxyde de carbone.
\item À l'aide des nombres d'oxydation, montrer qu'il s'agit d'une réaction d'oxydoréduction et identifier l'oxydant et le réducteur.
\item Calculer la masse de fer produite par la réduction complète de $m = \SI{5,00}{t}$ d'hématite pure.
\end{enumerate}

\textbf{Partie B — Préparation de l'acide sulfurique.}
Le grillage de la pyrite dans l'air donne l'oxyde de fer(III) et le dioxyde de soufre ; celui-ci est ensuite oxydé en trioxyde de soufre \ce{SO3}, qui, par absorption, conduit à l'acide sulfurique.
\begin{enumerate}
\item Écrire et équilibrer l'équation du grillage : $\ce{FeS2 + O2 -> Fe2O3 + SO2}$.
\item Déterminer les nombres d'oxydation du soufre dans \ce{FeS2}, \ce{SO2}, \ce{SO3} et \ce{H2SO4} (dans \ce{FeS2}, le fer est au n.o. $+2$).
\item Écrire et équilibrer l'équation d'oxydation de \ce{SO2} en \ce{SO3}, puis celle de la formation de l'acide sulfurique à partir de \ce{SO3} et de l'eau.
\item Quelle masse d'acide sulfurique peut-on théoriquement obtenir à partir de $m' = \SI{1,20}{t}$ de pyrite pure, le rendement global étant supposé de $100\,\%$ ?
\end{enumerate}
\emph{Données :} $M(\ce{Fe}) = \SI{55,8}{g.mol^{-1}}$ ; $M(\ce{S}) = \SI{32,1}{g.mol^{-1}}$ ; $M(\ce{O}) = \SI{16,0}{g.mol^{-1}}$ ; $M(\ce{H}) = \SI{1,0}{g.mol^{-1}}$.
\end{exercice}

\newpage

\coursec{Corrigés des exercices}

\begin{corrige}[Exercice 1 — QCM : réactions d'oxydoréduction]
\begin{enumerate}
\item \textbf{Bonnes réponses : b et c.} Une réaction d'oxydoréduction se caractérise par un \cle{transfert d'électrons} (c), qui se traduit par la variation du nombre d'oxydation d'au moins deux éléments (b). Le dioxygène n'est pas indispensable (contre-exemple : $\ce{H2 + Cl2 -> 2 HCl}$), donc (a) est faux. Le transfert de protons \ce{H+} caractérise les réactions acido-basiques, donc (d) est faux.
\item \textbf{Bonnes réponses : b et c.} Le n.o. de l'oxygène vaut $-2$ dans la plupart des composés, mais $-1$ dans les peroxydes (ex. \ce{H2O2}), $+2$ dans \ce{OF2} (le fluor est plus électronégatif), et $0$ dans le corps simple \ce{O2}.
\item \textbf{Bonnes réponses : b et c.} Dans $\ce{2 Mg + O2 -> 2 MgO}$, le n.o. du magnésium passe de $0$ à $+2$ : il est \emph{oxydé}, c'est donc le \emph{réducteur} ((a) et (b) faux). Le n.o. de l'oxygène passe de $0$ à $-2$ : \ce{O2} est réduit, c'est l'\emph{oxydant} (c).
\item \textbf{Bonnes réponses : c et d.} Dans une \cle{dismutation}, un même élément (d'une même espèce) est à la fois oxydé et réduit : l'espèce joue le double rôle d'oxydant et de réducteur.
\end{enumerate}
\textbf{Point de vigilance :} ne pas confondre « oxydation » (perte d'électrons, n.o. qui augmente) avec « réaction avec le dioxygène » : le mot « oxydation » a un sens plus large que « réaction avec l'oxygène ».
\end{corrige}

\begin{corrige}[Exercice 2 — Vrai ou faux, à justifier]
\begin{enumerate}
\item \textbf{Faux.} Dans \ce{SO4^2-} : $n.o.(\ce{S}) + 4 \times (-2) = -2$, d'où $n.o.(\ce{S}) = -2 + 8 = +6$ (et non $+4$).
\item \textbf{Faux.} Dans $\ce{CaCO3 -> CaO + CO2}$ : Ca garde le n.o. $+2$, C garde le n.o. $+4$, O garde le n.o. $-2$. Aucun n.o. ne varie : ce n'est \emph{pas} une réaction d'oxydoréduction (c'est une décomposition thermique).
\item \textbf{Faux.} \ce{H2} est un corps simple : le n.o. de l'hydrogène y est $0$. Il vaut $+1$ dans les \emph{composés} (sauf hydrures métalliques comme \ce{NaH}, où il vaut $-1$).
\item \textbf{Vrai.} Dans $\ce{C + O2 -> CO2}$ : le n.o. du carbone passe de $0$ à $+4$ ; il augmente, donc le carbone est \emph{oxydé} (il cède 4 électrons par atome).
\item \textbf{Vrai.} Par convention, le nombre d'oxydation d'un ion monoatomique est égal à sa charge : $n.o. = +2$ pour \ce{Ca^2+}, $-1$ pour \ce{Cl-}, etc.
\end{enumerate}
\end{corrige}

\begin{corrige}[Exercice 3 — Texte à trous]
Une réaction d'oxydoréduction par \textbf{voie sèche} se déroule en l'absence d'eau, souvent à haute température. Au cours d'une \textbf{oxydation}, le nombre d'oxydation d'un élément \textbf{augmente} ; au cours d'une \textbf{réduction}, il \textbf{diminue}. L'espèce qui capte des \textbf{électrons} est l'\textbf{oxydant} ; celle qui les cède est le \textbf{réducteur}. Dans la combustion du dihydrogène dans le dichlore $\ce{H2 + Cl2 -> 2 HCl}$, le dihydrogène joue le rôle de \textbf{réducteur} (son n.o. passe de $0$ à $+1$ : il est oxydé).
\end{corrige}

\begin{corrige}[Exercice 4 — Définitions et règles]
\begin{enumerate}
\item Le \cle{nombre d'oxydation} (n.o.) d'un élément dans un composé est la \emph{charge électrique fictive} que porterait l'atome de cet élément si toutes les liaisons de la molécule (ou de l'ion) étaient ioniques, les électrons de chaque liaison étant attribués à l'atome le plus électronégatif. On le note en chiffres romains ou en nombre algébrique ($+2$, $-1$\ldots).
\item (a) Dans un \emph{composé neutre}, la somme des n.o. de tous les éléments est \textbf{nulle}. (b) Dans un \emph{ion polyatomique}, la somme des n.o. est égale à la \textbf{charge de l'ion}.
\item Corps simple (ex. \ce{Fe}, \ce{O2}, \ce{H2}) : n.o. $= 0$. Hydrogène dans la plupart des composés : n.o. $= +1$ (sauf hydrures : $-1$). Oxygène dans la plupart des composés : n.o. $= -2$ (sauf peroxydes : $-1$, \ce{OF2} : $+2$, \ce{O2} : $0$).
\item Deux applications industrielles : la \textbf{sidérurgie} (réduction des oxydes de fer en haut fourneau par \ce{CO}) et l'\textbf{aluminothermie} (soudure des rails). On peut aussi citer la préparation des acides sulfurique et nitrique.
\end{enumerate}
\end{corrige}

\begin{corrige}[Exercice 5 — Calcul de nombres d'oxydation]
On applique : somme des n.o. $=$ charge de l'espèce, avec $n.o.(\ce{H}) = +1$, $n.o.(\ce{O}) = -2$, $n.o.(\ce{K}) = +1$.
\begin{enumerate}
\item \ce{H2SO4} : $2(+1) + x + 4(-2) = 0 \Rightarrow x = +6$. $n.o.(\ce{S}) = +6$.
\item \ce{HNO3} : $(+1) + x + 3(-2) = 0 \Rightarrow x = +5$. $n.o.(\ce{N}) = +5$.
\item \ce{KMnO4} : $(+1) + x + 4(-2) = 0 \Rightarrow x = +7$. $n.o.(\ce{Mn}) = +7$.
\item \ce{Cr2O7^2-} : $2x + 7(-2) = -2 \Rightarrow 2x = +12 \Rightarrow x = +6$. $n.o.(\ce{Cr}) = +6$.
\item \ce{NH4+} : $x + 4(+1) = +1 \Rightarrow x = -3$. $n.o.(\ce{N}) = -3$.
\item \ce{SO3^2-} : $x + 3(-2) = -2 \Rightarrow x = +4$. $n.o.(\ce{S}) = +4$.
\end{enumerate}
\textbf{Point de vigilance :} pour un ion, la somme des n.o. est égale à la \emph{charge de l'ion}, pas à zéro — c'est l'erreur la plus fréquente.
\end{corrige}

\begin{corrige}[Exercice 6 — Combustion du fer dans le dioxygène]
\begin{enumerate}
\item Dans \ce{Fe} (corps simple) : $n.o.(\ce{Fe}) = 0$. La magnétite \ce{Fe3O4} s'écrit \ce{FeO.Fe2O3} : elle contient un fer au n.o. $+2$ (dans \ce{FeO}) et deux fers au n.o. $+3$ (dans \ce{Fe2O3}). Le n.o. \emph{moyen} du fer est $\dfrac{(+2) + 2(+3)}{3} = +\dfrac{8}{3}$.
\item Oxygène : dans \ce{O2}, $n.o. = 0$ ; dans \ce{Fe3O4}, $n.o. = -2$.
\item Le n.o. du fer augmente (de $0$ à $+8/3$ en moyenne) et celui de l'oxygène diminue (de $0$ à $-2$) : deux éléments voient leur n.o. varier, c'est donc bien une \textbf{réaction d'oxydoréduction}.
\item Le fer est \emph{oxydé} : c'est le \textbf{réducteur}. Le dioxygène est \emph{réduit} : c'est l'\textbf{oxydant}.
\end{enumerate}
\textbf{Vérification de la conservation des électrons :} 3 atomes Fe perdent au total $8\ \text{e}^-$ (passage de $0$ à $+8/3$ en moyenne, soit $3 \times 8/3 = 8$) ; 2 molécules \ce{O2}, soit 4 atomes O, gagnent $4 \times 2 = 8\ \text{e}^-$. Le bilan est équilibré.
\end{corrige}

\begin{corrige}[Exercice 7 — Équilibrer par les nombres d'oxydation]
\begin{enumerate}
\item $\ce{Fe2O3 + CO -> Fe + CO2}$. Fer : $+3 \to 0$, variation $-3$ par atome ($\times 2$ atomes $= -6$). Carbone : $+2 \to +4$, variation $+2$. Produit en croix : il faut 3 CO pour 1 \ce{Fe2O3} ($3 \times 2 = 6$ électrons échangés) :
\[ \ce{Fe2O3 + 3 CO -> 2 Fe + 3 CO2} \]
\item $\ce{Al + Cr2O3 -> Al2O3 + Cr}$. Aluminium : $0 \to +3$ ($+3$ par atome). Chrome : $+3 \to 0$ ($-3$ par atome, $\times 2 = -6$). Il faut 2 Al ($2 \times 3 = 6$ électrons) :
\[ \ce{2 Al + Cr2O3 -> Al2O3 + 2 Cr} \]
\item $\ce{Cu + O2 -> CuO}$. Cuivre : $0 \to +2$ ($+2$). Oxygène : $0 \to -2$ ($-2$ par atome, $\times 2 = -4$ par \ce{O2}). Il faut 2 Cu par \ce{O2} :
\[ \ce{2 Cu + O2 -> 2 CuO} \]
\item $\ce{SO2 + O2 -> SO3}$. Soufre : $+4 \to +6$ ($+2$). Oxygène : $0 \to -2$ ($-4$ par \ce{O2}). Il faut 2 \ce{SO2} par \ce{O2} :
\[ \ce{2 SO2 + O2 -> 2 SO3} \]
\end{enumerate}
\textbf{Point de vigilance :} la variation de n.o. doit être multipliée par le \emph{nombre d'atomes} de l'élément dans l'espèce (ex. 2 Cr dans \ce{Cr2O3}) avant de faire le produit en croix.
\end{corrige}

\begin{corrige}[Exercice 8 — Le haut fourneau et la magnétite]
\begin{enumerate}
\item Équation-bilan : le fer passe de $+8/3$ (moyenne) à $0$ et le carbone de $+2$ (dans \ce{CO}) à $+4$ (dans \ce{CO2}) :
\[ \ce{Fe3O4 + 4 CO -> 3 Fe + 4 CO2} \]
Vérification : 4 CO apportent 4 C et 4 O ; les 4 O de \ce{Fe3O4} donnent 4 \ce{CO2}. Bilan : Fe : 3/3 ; C : 4/4 ; O : 8/8.
\item Fer : n.o. moyen $+8/3 \to 0$ ; au total, 3 Fe gagnent $8\ \text{e}^-$ : le fer est \emph{réduit}, \ce{Fe3O4} est l'\textbf{oxydant}. Carbone : $+2 \to +4$, chaque C perd $2\ \text{e}^-$ ; 4 C perdent $8\ \text{e}^-$ : le carbone est \emph{oxydé}, \ce{CO} est le \textbf{réducteur}.
\item Masse molaire de la magnétite :
\[ M(\ce{Fe3O4}) = 3 \times 55,8 + 4 \times 16,0 = 167,4 + 64,0 = \SI{231,4}{g.mol^{-1}}. \]
Quantité de magnétite : $n(\ce{Fe3O4}) = \dfrac{\num{1,0e6}}{\num{231,4}} = \SI{4321,5}{mol}$ (avec $m = \SI{1000}{kg} = \num{1,0e6}{g}$).
D'après l'équation, $n(\ce{Fe}) = 3\,n(\ce{Fe3O4}) = \SI{12964,6}{mol}$.
\[ m(\ce{Fe}) = \num{12964,6} \times 55,8 = \SI{723425}{g} \approx \SI{723}{kg}. \]
\textbf{Conclusion :} une tonne de magnétite pure fournit environ $\SI{0,72}{t}$ de fer (titre massique en fer de la magnétite : $\frac{167,4}{231,4} = 72,3\,\%$).
\end{enumerate}
\end{corrige}

\begin{corrige}[Exercice 9 — L'aluminothermie : soudure des rails]
\begin{enumerate}
\item Équation-bilan équilibrée :
\[ \ce{2 Al + Fe2O3 -> Al2O3 + 2 Fe} \]
\item Nombres d'oxydation : Al : $0 \to +3$ (perte de $3\ \text{e}^-$ par atome) ; Fe : $+3 \to 0$ (gain de $3\ \text{e}^-$ par atome). Bilan : 2 Al cèdent $6\ \text{e}^-$, 2 Fe en captent $6$.
\item L'aluminium est \emph{oxydé} : c'est le \textbf{réducteur}. \ce{Fe2O3} (le fer III) est \emph{réduit} : c'est l'\textbf{oxydant}. Espèce oxydée : \ce{Al} ; espèce réduite : \ce{Fe2O3}.
\item 
\begin{enumerate}
\item $n(\ce{Fe}) = \dfrac{m_{\ce{Fe}}}{M(\ce{Fe})} = \dfrac{\num{1000}}{55,8} = \SI{17,92}{mol}$.
\item D'après l'équation, $n(\ce{Al}) = n(\ce{Fe}) = \SI{17,92}{mol}$, d'où :
\[ m(\ce{Al}) = 17,92 \times 27,0 = \SI{483,9}{g} \approx \SI{484}{g}. \]
\item $n(\ce{Fe2O3}) = \dfrac{n(\ce{Fe})}{2} = \SI{8,96}{mol}$ ; $M(\ce{Fe2O3}) = 2 \times 55,8 + 3 \times 16,0 = \SI{159,6}{g.mol^{-1}}$ :
\[ m(\ce{Fe2O3}) = 8,96 \times 159,6 = \SI{1430}{g} \approx \SI{1,43}{kg}. \]
\end{enumerate}
\item Avec un rendement de $90\,\%$, il faut engager davantage d'aluminium :
\[ m_{\text{pratique}}(\ce{Al}) = \dfrac{483,9}{0,90} = \SI{537,6}{g} \approx \SI{538}{g}. \]
\end{enumerate}
\textbf{Barème indicatif (Probatoire) :} équation (1 pt) ; n.o. et électrons échangés (1 pt) ; identification oxydant/réducteur (1 pt) ; $n(\ce{Fe})$ (0,5 pt) ; $m(\ce{Al})$ (0,5 pt) ; $m(\ce{Fe2O3})$ (0,5 pt) ; rendement (0,5 pt).
\textbf{Points de vigilance :} attention au sens du rendement : pour un rendement inférieur à $100\,\%$, on \emph{divise} la masse théorique de réactif par le rendement (on en met plus), on ne multiplie pas. Bien convertir $\SI{1,00}{kg} = \SI{1000}{g}$.
\end{corrige}

\begin{corrige}[Exercice 10 — Le procédé Ostwald : préparation de l'acide nitrique]
\begin{enumerate}
\item Équations des trois étapes :
\begin{align*}
\text{Étape 1 : } & \ce{4 NH3 + 5 O2 -> 4 NO + 6 H2O} \quad (\text{catalyseur Pt-Rh, } \SI{900}{\celsius})\\
\text{Étape 2 : } & \ce{2 NO + O2 -> 2 NO2}\\
\text{Étape 3 : } & \ce{3 NO2 + H2O -> 2 HNO3 + NO}
\end{align*}
\item Nombres d'oxydation de l'azote : dans \ce{NH3} : $x + 3 = 0 \Rightarrow x = -3$ ; dans \ce{NO} : $x - 2 = 0 \Rightarrow x = +2$ ; dans \ce{NO2} : $x - 4 = 0 \Rightarrow x = +4$ ; dans \ce{HNO3} : $1 + x - 6 = 0 \Rightarrow x = +5$.
\item Étape 1 : n.o.(N) passe de $-3$ à $+2$ : il \emph{augmente}, l'azote est oxydé. Étape 2 : n.o.(N) passe de $+2$ à $+4$ : il augmente encore, l'azote est de nouveau oxydé. Les étapes 1 et 2 sont bien des \textbf{oxydations} de l'azote.
\item Étape 3 : dans \ce{NO2}, $n.o.(\ce{N}) = +4$. Après réaction : $+5$ dans \ce{HNO3} (oxydation) et $+2$ dans \ce{NO} (réduction). Le \emph{même élément} azote, au même degré $+4$, est à la fois oxydé et réduit : c'est une \cle{dismutation}. Sur 3 atomes d'azote engagés, \textbf{2 sont oxydés} (en \ce{HNO3}) et \textbf{1 est réduit} (en \ce{NO}) : proportions $2/3$ oxydé, $1/3$ réduit.
\item Avec recyclage total du \ce{NO} de l'étape 3 vers l'étape 2, tout l'azote de \ce{NH3} finit en \ce{HNO3}. En combinant les trois étapes (étape 1 $\times 3$, étape 2 $\times 6$ avec recyclage, étape 3 $\times 4$), les intermédiaires \ce{NO} et \ce{NO2} s'éliminent et l'on obtient :
\[ \ce{NH3 + 2 O2 -> HNO3 + H2O} \]
Vérification : N : 1/1 ; H : 3/3 ; O : 4/4.
Quantité d'ammoniac : $V = \SI{1000}{m^3} = \SI{1,0e6}{L}$ :
\[ n(\ce{NH3}) = \dfrac{V}{V_m} = \dfrac{\num{1,0e6}}{24,0} = \SI{41667}{mol}. \]
D'après l'équation globale, $n(\ce{HNO3}) = n(\ce{NH3}) = \SI{41667}{mol}$.
$M(\ce{HNO3}) = 1,0 + 14,0 + 3 \times 16,0 = \SI{63,0}{g.mol^{-1}}$.
\[ m(\ce{HNO3}) = 41667 \times 63,0 = \SI{2625000}{g} = \SI{2625}{kg} \approx \SI{2,63}{t}. \]
\end{enumerate}
\textbf{Barème indicatif :} 3 équations (1,5 pt) ; n.o. (1 pt) ; oxydations (0,5 pt) ; dismutation justifiée (1 pt) ; équation globale (0,5 pt) ; calcul de $n$ et $m$ (1,5 pt).
\textbf{Point de vigilance :} $\SI{1000}{m^3} = \num{1,0e6}{L}$ (et non $\num{1,0e3}{L}$) : la conversion de volume est le piège classique de cet exercice.
\end{corrige}

\begin{corrige}[Exercice 11 — Du minerai de fer à l'acide sulfurique]
\textbf{Partie A — Réduction de l'hématite.}
\begin{enumerate}
\item Équation équilibrée (voir exercice 7) :
\[ \ce{Fe2O3 + 3 CO -> 2 Fe + 3 CO2} \]
\item Fer : $+3 \to 0$ (diminution : réduction) ; carbone : $+2 \to +4$ (augmentation : oxydation). Deux n.o. varient : c'est une réaction d'\textbf{oxydoréduction}. \ce{Fe2O3} est l'\textbf{oxydant} (il est réduit), \ce{CO} est le \textbf{réducteur} (il est oxydé).
\item $M(\ce{Fe2O3}) = 2 \times 55,8 + 3 \times 16,0 = \SI{159,6}{g.mol^{-1}}$.
\[ n(\ce{Fe2O3}) = \dfrac{\num{5,00e6}}{159,6} = \SI{31328}{mol} \quad (m = \SI{5,00}{t} = \num{5,00e6}{g}). \]
$n(\ce{Fe}) = 2\,n(\ce{Fe2O3}) = \SI{62657}{mol}$, d'où :
\[ m(\ce{Fe}) = 62657 \times 55,8 = \SI{3496258}{g} \approx \SI{3,50}{t}. \]
(Contrôle rapide : titre en fer de l'hématite $= \frac{111,6}{159,6} = 69,9\,\%$ ; $0,699 \times 5,00 = \SI{3,50}{t}$.)
\end{enumerate}
\textbf{Partie B — Préparation de l'acide sulfurique.}
\begin{enumerate}
\item Grillage de la pyrite :
\[ \ce{4 FeS2 + 11 O2 -> 2 Fe2O3 + 8 SO2} \]
Vérification : Fe : 4/4 ; S : 8/8 ; O : 22 = 6 + 16.
\item Nombres d'oxydation du soufre : dans \ce{FeS2} avec Fe au n.o. $+2$ : $+2 + 2x = 0 \Rightarrow x = -1$ (ion disulfure \ce{S2^2-}) ; dans \ce{SO2} : $x = +4$ ; dans \ce{SO3} : $x = +6$ ; dans \ce{H2SO4} : $x = +6$.
\item Oxydation du dioxyde de soufre (procédé contact, catalyseur \ce{V2O5}) :
\[ \ce{2 SO2 + O2 -> 2 SO3} \]
Absorption (simplifiée) :
\[ \ce{SO3 + H2O -> H2SO4} \]
\item Masse molaire de la pyrite : $M(\ce{FeS2}) = 55,8 + 2 \times 32,1 = \SI{120,0}{g.mol^{-1}}$.
\[ n(\ce{FeS2}) = \dfrac{\num{1,20e6}}{120,0} = \SI{10000}{mol} \quad (m' = \SI{1,20}{t} = \num{1,20e6}{g}). \]
Bilan soufre : 1 mol de \ce{FeS2} contient 2 mol de S, donc conduit à 2 mol de \ce{H2SO4} :
\[ n(\ce{H2SO4}) = 2 \times 10000 = \SI{20000}{mol}. \]
$M(\ce{H2SO4}) = 2 \times 1,0 + 32,1 + 4 \times 16,0 = \SI{98,1}{g.mol^{-1}}$.
\[ m(\ce{H2SO4}) = 20000 \times 98,1 = \SI{1962000}{g} \approx \SI{1,96}{t}. \]
\end{enumerate}
\textbf{Barème indicatif (Probatoire) :} Partie A : équation (1 pt), n.o. et identification (1 pt), masse de fer (1 pt). Partie B : grillage (1 pt), n.o. du soufre (1 pt), deux équations (1 pt), masse d'acide (1,5 pt).
\textbf{Points de vigilance :} (1) dans \ce{FeS2}, le soufre est au n.o. $-1$ et non $-2$ : il faut utiliser l'indication « Fe au n.o. $+2$ » ; (2) ne pas oublier le facteur 2 entre \ce{FeS2} et \ce{H2SO4} (deux atomes de soufre par formule) ; (3) convertir les tonnes en grammes avant tout calcul de quantité de matière.
\end{corrige}

\newpage

\coursec{Fiche bilan}

\begin{popfigure}
\begin{tikzpicture}[scale=0.92, every node/.style={transform shape}]
\node[draw=popdark, fill=popblue, text=white, rounded corners=6pt, font=\bfseries\small, align=center, minimum width=3.4cm, minimum height=1.2cm] (C) at (0,0) {Oxydoréduction\\par voie sèche\\et nombres d'oxydation};
\node[draw=popblue, fill=popblueL, rounded corners=4pt, font=\footnotesize, align=center, text width=3.1cm] (N1) at (-5.2,2.6) {\textbf{Réactions par voie sèche}\\ \ce{2Mg + O2 -> 2MgO}\\ \ce{H2 + Cl2 -> 2HCl}\\ \ce{C + O2 -> CO2}};
\node[draw=poporange, fill=poporangeL, rounded corners=4pt, font=\footnotesize, align=center, text width=3.1cm] (N2) at (0,3.4) {\textbf{Nombre d'oxydation}\\ Charge fictive de l'atome\\ si toutes les liaisons\\ étaient ioniques};
\node[draw=popgreen, fill=popgreenL, rounded corners=4pt, font=\footnotesize, align=center, text width=3.1cm] (N3) at (5.2,2.6) {\textbf{Règles de calcul}\\ corps simple : 0\\ H : $+1$ \quad O : $-2$\\ somme $=$ charge};
\node[draw=poppink, fill=poppinkL, rounded corners=4pt, font=\footnotesize, align=center, text width=3.1cm] (N4) at (-5.2,-2.6) {\textbf{Reconnaître une redox}\\ n.o. qui varie\\ oxydation : n.o. $\nearrow$\\ réduction : n.o. $\searrow$};
\node[draw=poppurple, fill=poppurpleL, rounded corners=4pt, font=\footnotesize, align=center, text width=3.1cm] (N5) at (0,-3.4) {\textbf{Équilibrer}\\ égaliser les variations\\ de n.o. (électrons échangés)\\ puis H et O};
\node[draw=popgold, fill=popgoldL, rounded corners=4pt, font=\footnotesize, align=center, text width=3.3cm] (N6) at (5.2,-2.6) {\textbf{Applications}\\ haut fourneau : \ce{Fe2O3 + 3CO -> 2Fe + 3CO2}\\ aluminothermie : \ce{Fe2O3 + 2Al -> Al2O3 + 2Fe}};
\draw[-{Stealth}, thick, popblue] (C) -- (N1);
\draw[-{Stealth}, thick, poporange] (C) -- (N2);
\draw[-{Stealth}, thick, popgreen] (C) -- (N3);
\draw[-{Stealth}, thick, poppink] (C) -- (N4);
\draw[-{Stealth}, thick, poppurple] (C) -- (N5);
\draw[-{Stealth}, thick, popgold] (C) -- (N6);
\end{tikzpicture}
\legende{Carte mentale de la leçon : les six idées-forces à retenir.}
\end{popfigure}

\begin{bilan}[L'essentiel de la leçon]
\textbf{Définitions clés}
\begin{itemize}
\item Une \cle{réaction d'oxydoréduction par voie sèche} est une réaction redox qui se déroule \textbf{sans solvant} (en phase gazeuse ou entre solides), souvent lors de combustions.
\item Le \cle{nombre d'oxydation} (n.o.) d'un élément dans une espèce chimique est la \textbf{charge électrique fictive} que porterait l'atome si tous les électrons des liaisons étaient attribués à l'atome le plus électronégatif.
\item \cle{Oxydation} : augmentation du n.o. (perte d'électrons). \cle{Réduction} : diminution du n.o. (gain d'électrons). Une réaction est redox \textbf{si et seulement si} au moins un n.o. varie.
\end{itemize}

\textbf{Règles de détermination du nombre d'oxydation}
\begin{center}
\begin{tabularx}{\linewidth}{|l|X|}
\hline
\rowcolor{popblueL} \textbf{Cas} & \textbf{Règle} \\
\hline
Corps simple (\ce{Fe}, \ce{O2}, \ce{S8}) & n.o. $= 0$ \\
\hline
Ion monoatomique (\ce{Fe^3+}, \ce{Cl-}) & n.o. $=$ charge de l'ion \\
\hline
Hydrogène combiné & $+1$ (sauf hydrures : $-1$) \\
\hline
Oxygène combiné & $-2$ (sauf peroxydes : $-1$) \\
\hline
Molécule neutre & somme des n.o. $= 0$ \\
\hline
Ion polyatomique & somme des n.o. $=$ charge de l'ion \\
\hline
\end{tabularx}
\end{center}

\textbf{Équations-types à connaître par c\oe ur}
\begin{center}
\begin{tabularx}{\linewidth}{|X|X|}
\hline
\rowcolor{popblueL} \textbf{Situation} & \textbf{Équation-bilan} \\
\hline
Magnésium dans \ce{O2} & \ce{2Mg + O2 -> 2MgO} \\
\hline
Dihydrogène dans \ce{Cl2} & \ce{H2 + Cl2 -> 2HCl} \\
\hline
Carbone dans \ce{O2} & \ce{C + O2 -> CO2} \\
\hline
Haut fourneau (sidérurgie) & \ce{Fe2O3 + 3CO -> 2Fe + 3CO2} \\
\hline
Aluminothermie & \ce{Fe2O3 + 2Al -> Al2O3 + 2Fe} \\
\hline
Grillage de la pyrite & \ce{4FeS2 + 11O2 -> 2Fe2O3 + 8SO2} \\
\hline
Oxydation de \ce{SO2} & \ce{2SO2 + O2 <=> 2SO3} \\
\hline
Hydratation (acide sulfurique) & \ce{SO3 + H2O -> H2SO4} \\
\hline
Oxydation de l'ammoniac & \ce{4NH3 + 5O2 -> 4NO + 6H2O} \\
\hline
Formation de \ce{NO2} & \ce{2NO + O2 -> 2NO2} \\
\hline
Absorption (acide nitrique) & \ce{3NO2 + H2O -> 2HNO3 + NO} \\
\hline
\end{tabularx}
\end{center}

\textbf{Méthodes express}
\begin{itemize}
\item \textbf{Prouver qu'une réaction est redox} : calculer le n.o. de chaque élément avant et après ; si au moins un n.o. varie, c'est une réaction d'oxydoréduction. L'élément dont le n.o. \emph{augmente} est oxydé (c'est le réducteur) ; celui dont le n.o. \emph{diminue} est réduit (c'est l'oxydant).
\item \textbf{Équilibrer par les n.o.} : (1) repérer les éléments dont le n.o. varie ; (2) calculer les variations $\Delta(\text{n.o.})$ ; (3) ajuster les coefficients pour que gain d'électrons $=$ perte d'électrons ; (4) équilibrer les autres atomes (O puis H en dernier).
\item \textbf{Astuce de contrôle} : la somme des variations positives doit être égale, en valeur absolue, à la somme des variations négatives.
\end{itemize}
\end{bilan}

\begin{aretenir}[Checklist avant le Probatoire]
\begin{itemize}
\item[$\square$] Je sais définir une réaction d'oxydoréduction par voie sèche et citer trois exemples de combustions.
\item[$\square$] Je sais écrire les équations de combustion de \ce{Mg} dans \ce{O2}, de \ce{H2} dans \ce{Cl2} et de \ce{C} dans \ce{O2}.
\item[$\square$] Je connais la définition du nombre d'oxydation.
\item[$\square$] Je sais que le n.o. d'un corps simple est nul et que celui d'un ion monoatomique égale sa charge.
\item[$\square$] Je sais que H vaut $+1$ et O vaut $-2$ dans la plupart des composés (et je connais les exceptions).
\item[$\square$] Je sais calculer le n.o. d'un élément dans une molécule ou un ion polyatomique (ex. S dans \ce{SO4^2-} : $+6$).
\item[$\square$] Je sais montrer qu'une réaction est redox en comparant les n.o. avant et après.
\item[$\square$] Je sais identifier l'espèce oxydée, l'espèce réduite, l'oxydant et le réducteur.
\item[$\square$] Je sais équilibrer une équation redox par la méthode des variations de n.o.
\item[$\square$] Je sais décrire le principe du haut fourneau et de l'aluminothermie avec leurs équations.
\item[$\square$] Je sais écrire les trois étapes de la préparation industrielle de l'acide sulfurique.
\item[$\square$] Je sais écrire les trois étapes de la préparation industrielle de l'acide nitrique à partir de l'ammoniac.
\end{itemize}
\end{aretenir}

\findecours

\end{document}
